\documentclass[11pt]{article}

    \usepackage[breakable]{tcolorbox}
    \usepackage{parskip} % Stop auto-indenting (to mimic markdown behaviour)
    
    \usepackage{iftex}
    \ifPDFTeX
    	\usepackage[T1]{fontenc}
    	\usepackage{mathpazo}
    \else
    	\usepackage{fontspec}
    \fi

    % Basic figure setup, for now with no caption control since it's done
    % automatically by Pandoc (which extracts ![](path) syntax from Markdown).
    \usepackage{graphicx}
    % Maintain compatibility with old templates. Remove in nbconvert 6.0
    \let\Oldincludegraphics\includegraphics
    % Ensure that by default, figures have no caption (until we provide a
    % proper Figure object with a Caption API and a way to capture that
    % in the conversion process - todo).
    \usepackage{caption}
    \DeclareCaptionFormat{nocaption}{}
    \captionsetup{format=nocaption,aboveskip=0pt,belowskip=0pt}

    \usepackage[Export]{adjustbox} % Used to constrain images to a maximum size
    \adjustboxset{max size={0.9\linewidth}{0.9\paperheight}}
    \usepackage{float}
    \floatplacement{figure}{H} % forces figures to be placed at the correct location
    \usepackage{xcolor} % Allow colors to be defined
    \usepackage{enumerate} % Needed for markdown enumerations to work
    \usepackage{geometry} % Used to adjust the document margins
    \usepackage{amsmath} % Equations
    \usepackage{amssymb} % Equations
    \usepackage{textcomp} % defines textquotesingle
    % Hack from http://tex.stackexchange.com/a/47451/13684:
    \AtBeginDocument{%
        \def\PYZsq{\textquotesingle}% Upright quotes in Pygmentized code
    }
    \usepackage{upquote} % Upright quotes for verbatim code
    \usepackage{eurosym} % defines \euro
    \usepackage[mathletters]{ucs} % Extended unicode (utf-8) support
    \usepackage{fancyvrb} % verbatim replacement that allows latex
    \usepackage{grffile} % extends the file name processing of package graphics 
                         % to support a larger range
    \makeatletter % fix for grffile with XeLaTeX
    \def\Gread@@xetex#1{%
      \IfFileExists{"\Gin@base".bb}%
      {\Gread@eps{\Gin@base.bb}}%
      {\Gread@@xetex@aux#1}%
    }
    \makeatother

    % The hyperref package gives us a pdf with properly built
    % internal navigation ('pdf bookmarks' for the table of contents,
    % internal cross-reference links, web links for URLs, etc.)
    \usepackage{hyperref}
    % The default LaTeX title has an obnoxious amount of whitespace. By default,
    % titling removes some of it. It also provides customization options.
    \usepackage{titling}
    \usepackage{longtable} % longtable support required by pandoc >1.10
    \usepackage{booktabs}  % table support for pandoc > 1.12.2
    \usepackage[inline]{enumitem} % IRkernel/repr support (it uses the enumerate* environment)
    \usepackage[normalem]{ulem} % ulem is needed to support strikethroughs (\sout)
                                % normalem makes italics be italics, not underlines
    \usepackage{mathrsfs}
    

    
    % Colors for the hyperref package
    \definecolor{urlcolor}{rgb}{0,.145,.698}
    \definecolor{linkcolor}{rgb}{.71,0.21,0.01}
    \definecolor{citecolor}{rgb}{.12,.54,.11}

    % ANSI colors
    \definecolor{ansi-black}{HTML}{3E424D}
    \definecolor{ansi-black-intense}{HTML}{282C36}
    \definecolor{ansi-red}{HTML}{E75C58}
    \definecolor{ansi-red-intense}{HTML}{B22B31}
    \definecolor{ansi-green}{HTML}{00A250}
    \definecolor{ansi-green-intense}{HTML}{007427}
    \definecolor{ansi-yellow}{HTML}{DDB62B}
    \definecolor{ansi-yellow-intense}{HTML}{B27D12}
    \definecolor{ansi-blue}{HTML}{208FFB}
    \definecolor{ansi-blue-intense}{HTML}{0065CA}
    \definecolor{ansi-magenta}{HTML}{D160C4}
    \definecolor{ansi-magenta-intense}{HTML}{A03196}
    \definecolor{ansi-cyan}{HTML}{60C6C8}
    \definecolor{ansi-cyan-intense}{HTML}{258F8F}
    \definecolor{ansi-white}{HTML}{C5C1B4}
    \definecolor{ansi-white-intense}{HTML}{A1A6B2}
    \definecolor{ansi-default-inverse-fg}{HTML}{FFFFFF}
    \definecolor{ansi-default-inverse-bg}{HTML}{000000}

    % commands and environments needed by pandoc snippets
    % extracted from the output of `pandoc -s`
    \providecommand{\tightlist}{%
      \setlength{\itemsep}{0pt}\setlength{\parskip}{0pt}}
    \DefineVerbatimEnvironment{Highlighting}{Verbatim}{commandchars=\\\{\}}
    % Add ',fontsize=\small' for more characters per line
    \newenvironment{Shaded}{}{}
    \newcommand{\KeywordTok}[1]{\textcolor[rgb]{0.00,0.44,0.13}{\textbf{{#1}}}}
    \newcommand{\DataTypeTok}[1]{\textcolor[rgb]{0.56,0.13,0.00}{{#1}}}
    \newcommand{\DecValTok}[1]{\textcolor[rgb]{0.25,0.63,0.44}{{#1}}}
    \newcommand{\BaseNTok}[1]{\textcolor[rgb]{0.25,0.63,0.44}{{#1}}}
    \newcommand{\FloatTok}[1]{\textcolor[rgb]{0.25,0.63,0.44}{{#1}}}
    \newcommand{\CharTok}[1]{\textcolor[rgb]{0.25,0.44,0.63}{{#1}}}
    \newcommand{\StringTok}[1]{\textcolor[rgb]{0.25,0.44,0.63}{{#1}}}
    \newcommand{\CommentTok}[1]{\textcolor[rgb]{0.38,0.63,0.69}{\textit{{#1}}}}
    \newcommand{\OtherTok}[1]{\textcolor[rgb]{0.00,0.44,0.13}{{#1}}}
    \newcommand{\AlertTok}[1]{\textcolor[rgb]{1.00,0.00,0.00}{\textbf{{#1}}}}
    \newcommand{\FunctionTok}[1]{\textcolor[rgb]{0.02,0.16,0.49}{{#1}}}
    \newcommand{\RegionMarkerTok}[1]{{#1}}
    \newcommand{\ErrorTok}[1]{\textcolor[rgb]{1.00,0.00,0.00}{\textbf{{#1}}}}
    \newcommand{\NormalTok}[1]{{#1}}
    
    % Additional commands for more recent versions of Pandoc
    \newcommand{\ConstantTok}[1]{\textcolor[rgb]{0.53,0.00,0.00}{{#1}}}
    \newcommand{\SpecialCharTok}[1]{\textcolor[rgb]{0.25,0.44,0.63}{{#1}}}
    \newcommand{\VerbatimStringTok}[1]{\textcolor[rgb]{0.25,0.44,0.63}{{#1}}}
    \newcommand{\SpecialStringTok}[1]{\textcolor[rgb]{0.73,0.40,0.53}{{#1}}}
    \newcommand{\ImportTok}[1]{{#1}}
    \newcommand{\DocumentationTok}[1]{\textcolor[rgb]{0.73,0.13,0.13}{\textit{{#1}}}}
    \newcommand{\AnnotationTok}[1]{\textcolor[rgb]{0.38,0.63,0.69}{\textbf{\textit{{#1}}}}}
    \newcommand{\CommentVarTok}[1]{\textcolor[rgb]{0.38,0.63,0.69}{\textbf{\textit{{#1}}}}}
    \newcommand{\VariableTok}[1]{\textcolor[rgb]{0.10,0.09,0.49}{{#1}}}
    \newcommand{\ControlFlowTok}[1]{\textcolor[rgb]{0.00,0.44,0.13}{\textbf{{#1}}}}
    \newcommand{\OperatorTok}[1]{\textcolor[rgb]{0.40,0.40,0.40}{{#1}}}
    \newcommand{\BuiltInTok}[1]{{#1}}
    \newcommand{\ExtensionTok}[1]{{#1}}
    \newcommand{\PreprocessorTok}[1]{\textcolor[rgb]{0.74,0.48,0.00}{{#1}}}
    \newcommand{\AttributeTok}[1]{\textcolor[rgb]{0.49,0.56,0.16}{{#1}}}
    \newcommand{\InformationTok}[1]{\textcolor[rgb]{0.38,0.63,0.69}{\textbf{\textit{{#1}}}}}
    \newcommand{\WarningTok}[1]{\textcolor[rgb]{0.38,0.63,0.69}{\textbf{\textit{{#1}}}}}
    
    
    % Define a nice break command that doesn't care if a line doesn't already
    % exist.
    \def\br{\hspace*{\fill} \\* }
    % Math Jax compatibility definitions
    \def\gt{>}
    \def\lt{<}
    \let\Oldtex\TeX
    \let\Oldlatex\LaTeX
    \renewcommand{\TeX}{\textrm{\Oldtex}}
    \renewcommand{\LaTeX}{\textrm{\Oldlatex}}
    % Document parameters
    % Document title
    \title{Building\_a\_Recurrent\_Neural\_Network\_Step\_by\_Step}
    
    
    
    
    
% Pygments definitions
\makeatletter
\def\PY@reset{\let\PY@it=\relax \let\PY@bf=\relax%
    \let\PY@ul=\relax \let\PY@tc=\relax%
    \let\PY@bc=\relax \let\PY@ff=\relax}
\def\PY@tok#1{\csname PY@tok@#1\endcsname}
\def\PY@toks#1+{\ifx\relax#1\empty\else%
    \PY@tok{#1}\expandafter\PY@toks\fi}
\def\PY@do#1{\PY@bc{\PY@tc{\PY@ul{%
    \PY@it{\PY@bf{\PY@ff{#1}}}}}}}
\def\PY#1#2{\PY@reset\PY@toks#1+\relax+\PY@do{#2}}

\expandafter\def\csname PY@tok@w\endcsname{\def\PY@tc##1{\textcolor[rgb]{0.73,0.73,0.73}{##1}}}
\expandafter\def\csname PY@tok@c\endcsname{\let\PY@it=\textit\def\PY@tc##1{\textcolor[rgb]{0.25,0.50,0.50}{##1}}}
\expandafter\def\csname PY@tok@cp\endcsname{\def\PY@tc##1{\textcolor[rgb]{0.74,0.48,0.00}{##1}}}
\expandafter\def\csname PY@tok@k\endcsname{\let\PY@bf=\textbf\def\PY@tc##1{\textcolor[rgb]{0.00,0.50,0.00}{##1}}}
\expandafter\def\csname PY@tok@kp\endcsname{\def\PY@tc##1{\textcolor[rgb]{0.00,0.50,0.00}{##1}}}
\expandafter\def\csname PY@tok@kt\endcsname{\def\PY@tc##1{\textcolor[rgb]{0.69,0.00,0.25}{##1}}}
\expandafter\def\csname PY@tok@o\endcsname{\def\PY@tc##1{\textcolor[rgb]{0.40,0.40,0.40}{##1}}}
\expandafter\def\csname PY@tok@ow\endcsname{\let\PY@bf=\textbf\def\PY@tc##1{\textcolor[rgb]{0.67,0.13,1.00}{##1}}}
\expandafter\def\csname PY@tok@nb\endcsname{\def\PY@tc##1{\textcolor[rgb]{0.00,0.50,0.00}{##1}}}
\expandafter\def\csname PY@tok@nf\endcsname{\def\PY@tc##1{\textcolor[rgb]{0.00,0.00,1.00}{##1}}}
\expandafter\def\csname PY@tok@nc\endcsname{\let\PY@bf=\textbf\def\PY@tc##1{\textcolor[rgb]{0.00,0.00,1.00}{##1}}}
\expandafter\def\csname PY@tok@nn\endcsname{\let\PY@bf=\textbf\def\PY@tc##1{\textcolor[rgb]{0.00,0.00,1.00}{##1}}}
\expandafter\def\csname PY@tok@ne\endcsname{\let\PY@bf=\textbf\def\PY@tc##1{\textcolor[rgb]{0.82,0.25,0.23}{##1}}}
\expandafter\def\csname PY@tok@nv\endcsname{\def\PY@tc##1{\textcolor[rgb]{0.10,0.09,0.49}{##1}}}
\expandafter\def\csname PY@tok@no\endcsname{\def\PY@tc##1{\textcolor[rgb]{0.53,0.00,0.00}{##1}}}
\expandafter\def\csname PY@tok@nl\endcsname{\def\PY@tc##1{\textcolor[rgb]{0.63,0.63,0.00}{##1}}}
\expandafter\def\csname PY@tok@ni\endcsname{\let\PY@bf=\textbf\def\PY@tc##1{\textcolor[rgb]{0.60,0.60,0.60}{##1}}}
\expandafter\def\csname PY@tok@na\endcsname{\def\PY@tc##1{\textcolor[rgb]{0.49,0.56,0.16}{##1}}}
\expandafter\def\csname PY@tok@nt\endcsname{\let\PY@bf=\textbf\def\PY@tc##1{\textcolor[rgb]{0.00,0.50,0.00}{##1}}}
\expandafter\def\csname PY@tok@nd\endcsname{\def\PY@tc##1{\textcolor[rgb]{0.67,0.13,1.00}{##1}}}
\expandafter\def\csname PY@tok@s\endcsname{\def\PY@tc##1{\textcolor[rgb]{0.73,0.13,0.13}{##1}}}
\expandafter\def\csname PY@tok@sd\endcsname{\let\PY@it=\textit\def\PY@tc##1{\textcolor[rgb]{0.73,0.13,0.13}{##1}}}
\expandafter\def\csname PY@tok@si\endcsname{\let\PY@bf=\textbf\def\PY@tc##1{\textcolor[rgb]{0.73,0.40,0.53}{##1}}}
\expandafter\def\csname PY@tok@se\endcsname{\let\PY@bf=\textbf\def\PY@tc##1{\textcolor[rgb]{0.73,0.40,0.13}{##1}}}
\expandafter\def\csname PY@tok@sr\endcsname{\def\PY@tc##1{\textcolor[rgb]{0.73,0.40,0.53}{##1}}}
\expandafter\def\csname PY@tok@ss\endcsname{\def\PY@tc##1{\textcolor[rgb]{0.10,0.09,0.49}{##1}}}
\expandafter\def\csname PY@tok@sx\endcsname{\def\PY@tc##1{\textcolor[rgb]{0.00,0.50,0.00}{##1}}}
\expandafter\def\csname PY@tok@m\endcsname{\def\PY@tc##1{\textcolor[rgb]{0.40,0.40,0.40}{##1}}}
\expandafter\def\csname PY@tok@gh\endcsname{\let\PY@bf=\textbf\def\PY@tc##1{\textcolor[rgb]{0.00,0.00,0.50}{##1}}}
\expandafter\def\csname PY@tok@gu\endcsname{\let\PY@bf=\textbf\def\PY@tc##1{\textcolor[rgb]{0.50,0.00,0.50}{##1}}}
\expandafter\def\csname PY@tok@gd\endcsname{\def\PY@tc##1{\textcolor[rgb]{0.63,0.00,0.00}{##1}}}
\expandafter\def\csname PY@tok@gi\endcsname{\def\PY@tc##1{\textcolor[rgb]{0.00,0.63,0.00}{##1}}}
\expandafter\def\csname PY@tok@gr\endcsname{\def\PY@tc##1{\textcolor[rgb]{1.00,0.00,0.00}{##1}}}
\expandafter\def\csname PY@tok@ge\endcsname{\let\PY@it=\textit}
\expandafter\def\csname PY@tok@gs\endcsname{\let\PY@bf=\textbf}
\expandafter\def\csname PY@tok@gp\endcsname{\let\PY@bf=\textbf\def\PY@tc##1{\textcolor[rgb]{0.00,0.00,0.50}{##1}}}
\expandafter\def\csname PY@tok@go\endcsname{\def\PY@tc##1{\textcolor[rgb]{0.53,0.53,0.53}{##1}}}
\expandafter\def\csname PY@tok@gt\endcsname{\def\PY@tc##1{\textcolor[rgb]{0.00,0.27,0.87}{##1}}}
\expandafter\def\csname PY@tok@err\endcsname{\def\PY@bc##1{\setlength{\fboxsep}{0pt}\fcolorbox[rgb]{1.00,0.00,0.00}{1,1,1}{\strut ##1}}}
\expandafter\def\csname PY@tok@kc\endcsname{\let\PY@bf=\textbf\def\PY@tc##1{\textcolor[rgb]{0.00,0.50,0.00}{##1}}}
\expandafter\def\csname PY@tok@kd\endcsname{\let\PY@bf=\textbf\def\PY@tc##1{\textcolor[rgb]{0.00,0.50,0.00}{##1}}}
\expandafter\def\csname PY@tok@kn\endcsname{\let\PY@bf=\textbf\def\PY@tc##1{\textcolor[rgb]{0.00,0.50,0.00}{##1}}}
\expandafter\def\csname PY@tok@kr\endcsname{\let\PY@bf=\textbf\def\PY@tc##1{\textcolor[rgb]{0.00,0.50,0.00}{##1}}}
\expandafter\def\csname PY@tok@bp\endcsname{\def\PY@tc##1{\textcolor[rgb]{0.00,0.50,0.00}{##1}}}
\expandafter\def\csname PY@tok@fm\endcsname{\def\PY@tc##1{\textcolor[rgb]{0.00,0.00,1.00}{##1}}}
\expandafter\def\csname PY@tok@vc\endcsname{\def\PY@tc##1{\textcolor[rgb]{0.10,0.09,0.49}{##1}}}
\expandafter\def\csname PY@tok@vg\endcsname{\def\PY@tc##1{\textcolor[rgb]{0.10,0.09,0.49}{##1}}}
\expandafter\def\csname PY@tok@vi\endcsname{\def\PY@tc##1{\textcolor[rgb]{0.10,0.09,0.49}{##1}}}
\expandafter\def\csname PY@tok@vm\endcsname{\def\PY@tc##1{\textcolor[rgb]{0.10,0.09,0.49}{##1}}}
\expandafter\def\csname PY@tok@sa\endcsname{\def\PY@tc##1{\textcolor[rgb]{0.73,0.13,0.13}{##1}}}
\expandafter\def\csname PY@tok@sb\endcsname{\def\PY@tc##1{\textcolor[rgb]{0.73,0.13,0.13}{##1}}}
\expandafter\def\csname PY@tok@sc\endcsname{\def\PY@tc##1{\textcolor[rgb]{0.73,0.13,0.13}{##1}}}
\expandafter\def\csname PY@tok@dl\endcsname{\def\PY@tc##1{\textcolor[rgb]{0.73,0.13,0.13}{##1}}}
\expandafter\def\csname PY@tok@s2\endcsname{\def\PY@tc##1{\textcolor[rgb]{0.73,0.13,0.13}{##1}}}
\expandafter\def\csname PY@tok@sh\endcsname{\def\PY@tc##1{\textcolor[rgb]{0.73,0.13,0.13}{##1}}}
\expandafter\def\csname PY@tok@s1\endcsname{\def\PY@tc##1{\textcolor[rgb]{0.73,0.13,0.13}{##1}}}
\expandafter\def\csname PY@tok@mb\endcsname{\def\PY@tc##1{\textcolor[rgb]{0.40,0.40,0.40}{##1}}}
\expandafter\def\csname PY@tok@mf\endcsname{\def\PY@tc##1{\textcolor[rgb]{0.40,0.40,0.40}{##1}}}
\expandafter\def\csname PY@tok@mh\endcsname{\def\PY@tc##1{\textcolor[rgb]{0.40,0.40,0.40}{##1}}}
\expandafter\def\csname PY@tok@mi\endcsname{\def\PY@tc##1{\textcolor[rgb]{0.40,0.40,0.40}{##1}}}
\expandafter\def\csname PY@tok@il\endcsname{\def\PY@tc##1{\textcolor[rgb]{0.40,0.40,0.40}{##1}}}
\expandafter\def\csname PY@tok@mo\endcsname{\def\PY@tc##1{\textcolor[rgb]{0.40,0.40,0.40}{##1}}}
\expandafter\def\csname PY@tok@ch\endcsname{\let\PY@it=\textit\def\PY@tc##1{\textcolor[rgb]{0.25,0.50,0.50}{##1}}}
\expandafter\def\csname PY@tok@cm\endcsname{\let\PY@it=\textit\def\PY@tc##1{\textcolor[rgb]{0.25,0.50,0.50}{##1}}}
\expandafter\def\csname PY@tok@cpf\endcsname{\let\PY@it=\textit\def\PY@tc##1{\textcolor[rgb]{0.25,0.50,0.50}{##1}}}
\expandafter\def\csname PY@tok@c1\endcsname{\let\PY@it=\textit\def\PY@tc##1{\textcolor[rgb]{0.25,0.50,0.50}{##1}}}
\expandafter\def\csname PY@tok@cs\endcsname{\let\PY@it=\textit\def\PY@tc##1{\textcolor[rgb]{0.25,0.50,0.50}{##1}}}

\def\PYZbs{\char`\\}
\def\PYZus{\char`\_}
\def\PYZob{\char`\{}
\def\PYZcb{\char`\}}
\def\PYZca{\char`\^}
\def\PYZam{\char`\&}
\def\PYZlt{\char`\<}
\def\PYZgt{\char`\>}
\def\PYZsh{\char`\#}
\def\PYZpc{\char`\%}
\def\PYZdl{\char`\$}
\def\PYZhy{\char`\-}
\def\PYZsq{\char`\'}
\def\PYZdq{\char`\"}
\def\PYZti{\char`\~}
% for compatibility with earlier versions
\def\PYZat{@}
\def\PYZlb{[}
\def\PYZrb{]}
\makeatother


    % For linebreaks inside Verbatim environment from package fancyvrb. 
    \makeatletter
        \newbox\Wrappedcontinuationbox 
        \newbox\Wrappedvisiblespacebox 
        \newcommand*\Wrappedvisiblespace {\textcolor{red}{\textvisiblespace}} 
        \newcommand*\Wrappedcontinuationsymbol {\textcolor{red}{\llap{\tiny$\m@th\hookrightarrow$}}} 
        \newcommand*\Wrappedcontinuationindent {3ex } 
        \newcommand*\Wrappedafterbreak {\kern\Wrappedcontinuationindent\copy\Wrappedcontinuationbox} 
        % Take advantage of the already applied Pygments mark-up to insert 
        % potential linebreaks for TeX processing. 
        %        {, <, #, %, $, ' and ": go to next line. 
        %        _, }, ^, &, >, - and ~: stay at end of broken line. 
        % Use of \textquotesingle for straight quote. 
        \newcommand*\Wrappedbreaksatspecials {% 
            \def\PYGZus{\discretionary{\char`\_}{\Wrappedafterbreak}{\char`\_}}% 
            \def\PYGZob{\discretionary{}{\Wrappedafterbreak\char`\{}{\char`\{}}% 
            \def\PYGZcb{\discretionary{\char`\}}{\Wrappedafterbreak}{\char`\}}}% 
            \def\PYGZca{\discretionary{\char`\^}{\Wrappedafterbreak}{\char`\^}}% 
            \def\PYGZam{\discretionary{\char`\&}{\Wrappedafterbreak}{\char`\&}}% 
            \def\PYGZlt{\discretionary{}{\Wrappedafterbreak\char`\<}{\char`\<}}% 
            \def\PYGZgt{\discretionary{\char`\>}{\Wrappedafterbreak}{\char`\>}}% 
            \def\PYGZsh{\discretionary{}{\Wrappedafterbreak\char`\#}{\char`\#}}% 
            \def\PYGZpc{\discretionary{}{\Wrappedafterbreak\char`\%}{\char`\%}}% 
            \def\PYGZdl{\discretionary{}{\Wrappedafterbreak\char`\$}{\char`\$}}% 
            \def\PYGZhy{\discretionary{\char`\-}{\Wrappedafterbreak}{\char`\-}}% 
            \def\PYGZsq{\discretionary{}{\Wrappedafterbreak\textquotesingle}{\textquotesingle}}% 
            \def\PYGZdq{\discretionary{}{\Wrappedafterbreak\char`\"}{\char`\"}}% 
            \def\PYGZti{\discretionary{\char`\~}{\Wrappedafterbreak}{\char`\~}}% 
        } 
        % Some characters . , ; ? ! / are not pygmentized. 
        % This macro makes them "active" and they will insert potential linebreaks 
        \newcommand*\Wrappedbreaksatpunct {% 
            \lccode`\~`\.\lowercase{\def~}{\discretionary{\hbox{\char`\.}}{\Wrappedafterbreak}{\hbox{\char`\.}}}% 
            \lccode`\~`\,\lowercase{\def~}{\discretionary{\hbox{\char`\,}}{\Wrappedafterbreak}{\hbox{\char`\,}}}% 
            \lccode`\~`\;\lowercase{\def~}{\discretionary{\hbox{\char`\;}}{\Wrappedafterbreak}{\hbox{\char`\;}}}% 
            \lccode`\~`\:\lowercase{\def~}{\discretionary{\hbox{\char`\:}}{\Wrappedafterbreak}{\hbox{\char`\:}}}% 
            \lccode`\~`\?\lowercase{\def~}{\discretionary{\hbox{\char`\?}}{\Wrappedafterbreak}{\hbox{\char`\?}}}% 
            \lccode`\~`\!\lowercase{\def~}{\discretionary{\hbox{\char`\!}}{\Wrappedafterbreak}{\hbox{\char`\!}}}% 
            \lccode`\~`\/\lowercase{\def~}{\discretionary{\hbox{\char`\/}}{\Wrappedafterbreak}{\hbox{\char`\/}}}% 
            \catcode`\.\active
            \catcode`\,\active 
            \catcode`\;\active
            \catcode`\:\active
            \catcode`\?\active
            \catcode`\!\active
            \catcode`\/\active 
            \lccode`\~`\~ 	
        }
    \makeatother

    \let\OriginalVerbatim=\Verbatim
    \makeatletter
    \renewcommand{\Verbatim}[1][1]{%
        %\parskip\z@skip
        \sbox\Wrappedcontinuationbox {\Wrappedcontinuationsymbol}%
        \sbox\Wrappedvisiblespacebox {\FV@SetupFont\Wrappedvisiblespace}%
        \def\FancyVerbFormatLine ##1{\hsize\linewidth
            \vtop{\raggedright\hyphenpenalty\z@\exhyphenpenalty\z@
                \doublehyphendemerits\z@\finalhyphendemerits\z@
                \strut ##1\strut}%
        }%
        % If the linebreak is at a space, the latter will be displayed as visible
        % space at end of first line, and a continuation symbol starts next line.
        % Stretch/shrink are however usually zero for typewriter font.
        \def\FV@Space {%
            \nobreak\hskip\z@ plus\fontdimen3\font minus\fontdimen4\font
            \discretionary{\copy\Wrappedvisiblespacebox}{\Wrappedafterbreak}
            {\kern\fontdimen2\font}%
        }%
        
        % Allow breaks at special characters using \PYG... macros.
        \Wrappedbreaksatspecials
        % Breaks at punctuation characters . , ; ? ! and / need catcode=\active 	
        \OriginalVerbatim[#1,codes*=\Wrappedbreaksatpunct]%
    }
    \makeatother

    % Exact colors from NB
    \definecolor{incolor}{HTML}{303F9F}
    \definecolor{outcolor}{HTML}{D84315}
    \definecolor{cellborder}{HTML}{CFCFCF}
    \definecolor{cellbackground}{HTML}{F7F7F7}
    
    % prompt
    \makeatletter
    \newcommand{\boxspacing}{\kern\kvtcb@left@rule\kern\kvtcb@boxsep}
    \makeatother
    \newcommand{\prompt}[4]{
        \ttfamily\llap{{\color{#2}[#3]:\hspace{3pt}#4}}\vspace{-\baselineskip}
    }
    

    
    % Prevent overflowing lines due to hard-to-break entities
    \sloppy 
    % Setup hyperref package
    \hypersetup{
      breaklinks=true,  % so long urls are correctly broken across lines
      colorlinks=true,
      urlcolor=urlcolor,
      linkcolor=linkcolor,
      citecolor=citecolor,
      }
    % Slightly bigger margins than the latex defaults
    
    \geometry{verbose,tmargin=1in,bmargin=1in,lmargin=1in,rmargin=1in}
    
    

\begin{document}
    
    \maketitle
    
    

    
    \hypertarget{building-your-recurrent-neural-network---step-by-step}{%
\section{Building your Recurrent Neural Network - Step by
Step}\label{building-your-recurrent-neural-network---step-by-step}}

Welcome to Course 5's first assignment, where you'll be implementing key
components of a Recurrent Neural Network, or RNN, in NumPy!

By the end of this assignment, you'll be able to:

\begin{itemize}
\tightlist
\item
  Define notation for building sequence models
\item
  Describe the architecture of a basic RNN
\item
  Identify the main components of an LSTM
\item
  Implement backpropagation through time for a basic RNN and an LSTM
\item
  Give examples of several types of RNN
\end{itemize}

Recurrent Neural Networks (RNN) are very effective for Natural Language
Processing and other sequence tasks because they have ``memory.'' They
can read inputs \(x^{\langle t \rangle}\) (such as words) one at a time,
and remember some contextual information through the hidden layer
activations that get passed from one time step to the next. This allows
a unidirectional (one-way) RNN to take information from the past to
process later inputs. A bidirectional (two-way) RNN can take context
from both the past and the future, much like Marty McFly.

\textbf{Notation}: - Superscript \([l]\) denotes an object associated
with the \(l^{th}\) layer.

\begin{itemize}
\item
  Superscript \((i)\) denotes an object associated with the \(i^{th}\)
  example.
\item
  Superscript \(\langle t \rangle\) denotes an object at the \(t^{th}\)
  time step.
\item
  Subscript \(i\) denotes the \(i^{th}\) entry of a vector.
\end{itemize}

\textbf{Example}:\\
- \(a^{(2)[3]<4>}_5\) denotes the activation of the 2nd training example
(2), 3rd layer {[}3{]}, 4th time step \textless4\textgreater, and 5th
entry in the vector.

\hypertarget{pre-requisites}{%
\paragraph{Pre-requisites}\label{pre-requisites}}

\begin{itemize}
\tightlist
\item
  You should already be familiar with \texttt{numpy}
\item
  To refresh your knowledge of numpy, you can review course 1 of the
  specialization ``Neural Networks and Deep Learning'':

  \begin{itemize}
  \tightlist
  \item
    Specifically, review the week 2's practice assignment
    \href{https://www.coursera.org/learn/neural-networks-deep-learning/programming/isoAV/python-basics-with-numpy}{``Python
    Basics with Numpy (optional assignment)''}
  \end{itemize}
\end{itemize}

\hypertarget{be-careful-when-modifying-the-starter-code}{%
\paragraph{Be careful when modifying the starter
code!}\label{be-careful-when-modifying-the-starter-code}}

\begin{itemize}
\tightlist
\item
  When working on graded functions, please remember to only modify the
  code that is between:
\end{itemize}

\begin{Shaded}
\begin{Highlighting}[]
\CommentTok{\#\#\#\# START CODE HERE}
\end{Highlighting}
\end{Shaded}

and:

\begin{Shaded}
\begin{Highlighting}[]
\CommentTok{\#\#\#\# }\RegionMarkerTok{END}\CommentTok{ CODE HERE}
\end{Highlighting}
\end{Shaded}

\begin{itemize}
\tightlist
\item
  In particular, avoid modifying the first line of graded routines.
  These start with:
\end{itemize}

\begin{Shaded}
\begin{Highlighting}[]
\CommentTok{\# GRADED FUNCTION: routine\_name}
\end{Highlighting}
\end{Shaded}

The automatic grader (autograder) needs these to locate the function -
so even a change in spacing will cause issues with the autograder,
returning `failed' if any of these are modified or missing. Now, let's
get started!

\hypertarget{important-note-on-submission-to-the-autograder}{%
\subsection{Important Note on Submission to the
AutoGrader}\label{important-note-on-submission-to-the-autograder}}

Before submitting your assignment to the AutoGrader, please make sure
you are not doing the following:

\begin{enumerate}
\def\labelenumi{\arabic{enumi}.}
\tightlist
\item
  You have not added any \emph{extra} \texttt{print} statement(s) in the
  assignment.
\item
  You have not added any \emph{extra} code cell(s) in the assignment.
\item
  You have not changed any of the function parameters.
\item
  You are not using any global variables inside your graded exercises.
  Unless specifically instructed to do so, please refrain from it and
  use the local variables instead.
\item
  You are not changing the assignment code where it is not required,
  like creating \emph{extra} variables.
\end{enumerate}

If you do any of the following, you will get something like,
\texttt{Grader\ Error:\ Grader\ feedback\ not\ found} (or similarly
unexpected) error upon submitting your assignment. Before asking for
help/debugging the errors in your assignment, check for these first. If
this is the case, and you don't remember the changes you have made, you
can get a fresh copy of the assignment by following these
\href{https://www.coursera.org/learn/nlp-sequence-models/supplement/qHIve/h-ow-to-refresh-your-workspace}{instructions}.

    \hypertarget{table-of-content}{%
\subsection{Table of Content}\label{table-of-content}}

\begin{itemize}
\tightlist
\item
  Section \ref{0}
\item
  Section \ref{1}

  \begin{itemize}
  \tightlist
  \item
    Section \ref{1-1}

    \begin{itemize}
    \tightlist
    \item
      Section \ref{ex-1}
    \end{itemize}
  \item
    Section \ref{1-2}

    \begin{itemize}
    \tightlist
    \item
      Section \ref{ex-2}
    \end{itemize}
  \end{itemize}
\item
  Section \ref{2}

  \begin{itemize}
  \tightlist
  \item
    Section \ref{2-1}

    \begin{itemize}
    \tightlist
    \item
      Section \ref{ex-3}
    \end{itemize}
  \item
    Section \ref{2-2}

    \begin{itemize}
    \tightlist
    \item
      Section \ref{ex-4}
    \end{itemize}
  \end{itemize}
\item
  Section \ref{3}

  \begin{itemize}
  \tightlist
  \item
    Section \ref{3-1}

    \begin{itemize}
    \tightlist
    \item
      Section \ref{ex-5}
    \item
      Section \ref{ex-6}
    \end{itemize}
  \item
    Section \ref{3-2}

    \begin{itemize}
    \tightlist
    \item
      Section \ref{ex-7}
    \end{itemize}
  \item
    Section \ref{3-3}

    \begin{itemize}
    \tightlist
    \item
      Section \ref{ex-8}
    \end{itemize}
  \end{itemize}
\end{itemize}

    \#\# Packages

    \begin{tcolorbox}[breakable, size=fbox, boxrule=1pt, pad at break*=1mm,colback=cellbackground, colframe=cellborder]
\prompt{In}{incolor}{3}{\boxspacing}
\begin{Verbatim}[commandchars=\\\{\}]
\PY{c+c1}{\PYZsh{}\PYZsh{}\PYZsh{} v2.2}
\end{Verbatim}
\end{tcolorbox}

    \begin{tcolorbox}[breakable, size=fbox, boxrule=1pt, pad at break*=1mm,colback=cellbackground, colframe=cellborder]
\prompt{In}{incolor}{4}{\boxspacing}
\begin{Verbatim}[commandchars=\\\{\}]
\PY{k+kn}{import} \PY{n+nn}{numpy} \PY{k}{as} \PY{n+nn}{np}
\PY{k+kn}{from} \PY{n+nn}{rnn\PYZus{}utils} \PY{k+kn}{import} \PY{o}{*}
\PY{k+kn}{from} \PY{n+nn}{public\PYZus{}tests} \PY{k+kn}{import} \PY{o}{*}
\end{Verbatim}
\end{tcolorbox}

    \#\# 1 - Forward Propagation for the Basic Recurrent Neural Network

Later this week, you'll get a chance to generate music using an RNN! The
basic RNN that you'll implement has the following structure:

In this example, \(T_x = T_y\).

    Figure 1: Basic RNN model

    \hypertarget{dimensions-of-input-x}{%
\subsubsection{\texorpdfstring{Dimensions of input
\(x\)}{Dimensions of input x}}\label{dimensions-of-input-x}}

\hypertarget{input-with-n_x-number-of-units}{%
\paragraph{\texorpdfstring{Input with \(n_x\) number of
units}{Input with n\_x number of units}}\label{input-with-n_x-number-of-units}}

\begin{itemize}
\tightlist
\item
  For a single time step of a single input example,
  \(x^{(i) \langle t \rangle }\) is a one-dimensional input vector
\item
  Using language as an example, a language with a 5000-word vocabulary
  could be one-hot encoded into a vector that has 5000 units. So
  \(x^{(i)\langle t \rangle}\) would have the shape (5000,)\\
\item
  The notation \(n_x\) is used here to denote the number of units in a
  single time step of a single training example
\end{itemize}

    \hypertarget{time-steps-of-size-t_x}{%
\paragraph{\texorpdfstring{Time steps of size
\(T_{x}\)}{Time steps of size T\_\{x\}}}\label{time-steps-of-size-t_x}}

\begin{itemize}
\tightlist
\item
  A recurrent neural network has multiple time steps, which you'll index
  with \(t\).
\item
  In the lessons, you saw a single training example \(x^{(i)}\)
  consisting of multiple time steps \(T_x\). In this notebook, \(T_{x}\)
  will denote the number of timesteps in the longest sequence.
\end{itemize}

    \hypertarget{batches-of-size-m}{%
\paragraph{\texorpdfstring{Batches of size
\(m\)}{Batches of size m}}\label{batches-of-size-m}}

\begin{itemize}
\tightlist
\item
  Let's say we have mini-batches, each with 20 training examples\\
\item
  To benefit from vectorization, you'll stack 20 columns of \(x^{(i)}\)
  examples
\item
  For example, this tensor has the shape (5000,20,10)
\item
  You'll use \(m\) to denote the number of training examples\\
\item
  So, the shape of a mini-batch is \((n_x,m,T_x)\)
\end{itemize}

    \hypertarget{d-tensor-of-shape-n_xmt_x}{%
\paragraph{\texorpdfstring{3D Tensor of shape
\((n_{x},m,T_{x})\)}{3D Tensor of shape (n\_\{x\},m,T\_\{x\})}}\label{d-tensor-of-shape-n_xmt_x}}

\begin{itemize}
\tightlist
\item
  The 3-dimensional tensor \(x\) of shape \((n_x,m,T_x)\) represents the
  input \(x\) that is fed into the RNN
\end{itemize}

\hypertarget{taking-a-2d-slice-for-each-time-step-xlangle-t-rangle}{%
\paragraph{\texorpdfstring{Taking a 2D slice for each time step:
\(x^{\langle t \rangle}\)}{Taking a 2D slice for each time step: x\^{}\{\textbackslash langle t \textbackslash rangle\}}}\label{taking-a-2d-slice-for-each-time-step-xlangle-t-rangle}}

\begin{itemize}
\tightlist
\item
  At each time step, you'll use a mini-batch of training examples (not
  just a single example)
\item
  So, for each time step \(t\), you'll use a 2D slice of shape
  \((n_x,m)\)
\item
  This 2D slice is referred to as \(x^{\langle t \rangle}\). The
  variable name in the code is \texttt{xt}.
\end{itemize}

    \hypertarget{definition-of-hidden-state-a}{%
\subsubsection{\texorpdfstring{Definition of hidden state
\(a\)}{Definition of hidden state a}}\label{definition-of-hidden-state-a}}

\begin{itemize}
\tightlist
\item
  The activation \(a^{\langle t \rangle}\) that is passed to the RNN
  from one time step to another is called a ``hidden state.''
\end{itemize}

\hypertarget{dimensions-of-hidden-state-a}{%
\subsubsection{\texorpdfstring{Dimensions of hidden state
\(a\)}{Dimensions of hidden state a}}\label{dimensions-of-hidden-state-a}}

\begin{itemize}
\tightlist
\item
  Similar to the input tensor \(x\), the hidden state for a single
  training example is a vector of length \(n_{a}\)
\item
  If you include a mini-batch of \(m\) training examples, the shape of a
  mini-batch is \((n_{a},m)\)
\item
  When you include the time step dimension, the shape of the hidden
  state is \((n_{a}, m, T_x)\)
\item
  You'll loop through the time steps with index \(t\), and work with a
  2D slice of the 3D tensor\\
\item
  This 2D slice is referred to as \(a^{\langle t \rangle}\)
\item
  In the code, the variable names used are either \texttt{a\_prev} or
  \texttt{a\_next}, depending on the function being implemented
\item
  The shape of this 2D slice is \((n_{a}, m)\)
\end{itemize}

    \hypertarget{dimensions-of-prediction-haty}{%
\subsubsection{\texorpdfstring{Dimensions of prediction
\(\hat{y}\)}{Dimensions of prediction \textbackslash hat\{y\}}}\label{dimensions-of-prediction-haty}}

\begin{itemize}
\tightlist
\item
  Similar to the inputs and hidden states, \(\hat{y}\) is a 3D tensor of
  shape \((n_{y}, m, T_{y})\)

  \begin{itemize}
  \tightlist
  \item
    \(n_{y}\): number of units in the vector representing the prediction
  \item
    \(m\): number of examples in a mini-batch
  \item
    \(T_{y}\): number of time steps in the prediction
  \end{itemize}
\item
  For a single time step \(t\), a 2D slice
  \(\hat{y}^{\langle t \rangle}\) has shape \((n_{y}, m)\)
\item
  In the code, the variable names are:

  \begin{itemize}
  \tightlist
  \item
    \texttt{y\_pred}: \(\hat{y}\)
  \item
    \texttt{yt\_pred}: \(\hat{y}^{\langle t \rangle}\)
  \end{itemize}
\end{itemize}

    Here's how you can implement an RNN:

\hypertarget{steps}{%
\subsubsection{Steps:}\label{steps}}

\begin{enumerate}
\def\labelenumi{\arabic{enumi}.}
\tightlist
\item
  Implement the calculations needed for one time step of the RNN.
\item
  Implement a loop over \(T_x\) time steps in order to process all the
  inputs, one at a time.
\end{enumerate}

    \#\#\# 1.1 - RNN Cell

You can think of the recurrent neural network as the repeated use of a
single cell. First, you'll implement the computations for a single time
step. The following figure describes the operations for a single time
step of an RNN cell:

Figure 2: Basic RNN cell. Takes as input \(x^{\langle t \rangle}\)
(current input) and \(a^{\langle t - 1\rangle}\) (previous hidden state
containing information from the past), and outputs
\(a^{\langle t \rangle}\) which is given to the next RNN cell and also
used to predict \(\hat{y}^{\langle t \rangle}\)

\textbf{\texttt{RNN\ cell} versus \texttt{RNN\_cell\_forward}}: * Note
that an RNN cell outputs the hidden state \(a^{\langle t \rangle}\).\\
* \texttt{RNN\ cell} is shown in the figure as the inner box with solid
lines\\
* The function that you'll implement, \texttt{rnn\_cell\_forward}, also
calculates the prediction \(\hat{y}^{\langle t \rangle}\) *
\texttt{RNN\_cell\_forward} is shown in the figure as the outer box with
dashed lines

    \#\#\# Exercise 1 - rnn\_cell\_forward

Implement the RNN cell described in Figure 2.

\textbf{Instructions}: 1. Compute the hidden state with tanh activation:
\(a^{\langle t \rangle} = \tanh(W_{aa} a^{\langle t-1 \rangle} + W_{ax} x^{\langle t \rangle} + b_a)\)
2. Using your new hidden state \(a^{\langle t \rangle}\), compute the
prediction
\(\hat{y}^{\langle t \rangle} = softmax(W_{ya} a^{\langle t \rangle} + b_y)\).
(The function \texttt{softmax} is provided) 3. Store
\((a^{\langle t \rangle}, a^{\langle t-1 \rangle}, x^{\langle t \rangle}, parameters)\)
in a \texttt{cache} 4. Return \(a^{\langle t \rangle}\) ,
\(\hat{y}^{\langle t \rangle}\) and \texttt{cache}

\hypertarget{additional-hints}{%
\paragraph{Additional Hints}\label{additional-hints}}

\begin{itemize}
\tightlist
\item
  A little more information on
  \href{https://numpy.org/devdocs/reference/generated/numpy.tanh.html}{numpy.tanh}
\item
  In this assignment, there's an existing \texttt{softmax} function for
  you to use. It's located in the file `rnn\_utils.py' and has already
  been imported.
\item
  For matrix multiplication, use
  \href{https://docs.scipy.org/doc/numpy/reference/generated/numpy.dot.html}{numpy.dot}
\end{itemize}

    \begin{tcolorbox}[breakable, size=fbox, boxrule=1pt, pad at break*=1mm,colback=cellbackground, colframe=cellborder]
\prompt{In}{incolor}{13}{\boxspacing}
\begin{Verbatim}[commandchars=\\\{\}]
\PY{c+c1}{\PYZsh{} UNQ\PYZus{}C1 (UNIQUE CELL IDENTIFIER, DO NOT EDIT)}
\PY{c+c1}{\PYZsh{} GRADED FUNCTION: rnn\PYZus{}cell\PYZus{}forward}

\PY{k}{def} \PY{n+nf}{rnn\PYZus{}cell\PYZus{}forward}\PY{p}{(}\PY{n}{xt}\PY{p}{,} \PY{n}{a\PYZus{}prev}\PY{p}{,} \PY{n}{parameters}\PY{p}{)}\PY{p}{:}
    \PY{l+s+sd}{\PYZdq{}\PYZdq{}\PYZdq{}}
\PY{l+s+sd}{    Implements a single forward step of the RNN\PYZhy{}cell as described in Figure (2)}

\PY{l+s+sd}{    Arguments:}
\PY{l+s+sd}{    xt \PYZhy{}\PYZhy{} your input data at timestep \PYZdq{}t\PYZdq{}, numpy array of shape (n\PYZus{}x, m).}
\PY{l+s+sd}{    a\PYZus{}prev \PYZhy{}\PYZhy{} Hidden state at timestep \PYZdq{}t\PYZhy{}1\PYZdq{}, numpy array of shape (n\PYZus{}a, m)}
\PY{l+s+sd}{    parameters \PYZhy{}\PYZhy{} python dictionary containing:}
\PY{l+s+sd}{                        Wax \PYZhy{}\PYZhy{} Weight matrix multiplying the input, numpy array of shape (n\PYZus{}a, n\PYZus{}x)}
\PY{l+s+sd}{                        Waa \PYZhy{}\PYZhy{} Weight matrix multiplying the hidden state, numpy array of shape (n\PYZus{}a, n\PYZus{}a)}
\PY{l+s+sd}{                        Wya \PYZhy{}\PYZhy{} Weight matrix relating the hidden\PYZhy{}state to the output, numpy array of shape (n\PYZus{}y, n\PYZus{}a)}
\PY{l+s+sd}{                        ba \PYZhy{}\PYZhy{}  Bias, numpy array of shape (n\PYZus{}a, 1)}
\PY{l+s+sd}{                        by \PYZhy{}\PYZhy{} Bias relating the hidden\PYZhy{}state to the output, numpy array of shape (n\PYZus{}y, 1)}
\PY{l+s+sd}{    Returns:}
\PY{l+s+sd}{    a\PYZus{}next \PYZhy{}\PYZhy{} next hidden state, of shape (n\PYZus{}a, m)}
\PY{l+s+sd}{    yt\PYZus{}pred \PYZhy{}\PYZhy{} prediction at timestep \PYZdq{}t\PYZdq{}, numpy array of shape (n\PYZus{}y, m)}
\PY{l+s+sd}{    cache \PYZhy{}\PYZhy{} tuple of values needed for the backward pass, contains (a\PYZus{}next, a\PYZus{}prev, xt, parameters)}
\PY{l+s+sd}{    \PYZdq{}\PYZdq{}\PYZdq{}}
    
    \PY{c+c1}{\PYZsh{} Retrieve parameters from \PYZdq{}parameters\PYZdq{}}
    \PY{n}{Wax} \PY{o}{=} \PY{n}{parameters}\PY{p}{[}\PY{l+s+s2}{\PYZdq{}}\PY{l+s+s2}{Wax}\PY{l+s+s2}{\PYZdq{}}\PY{p}{]}
    \PY{n}{Waa} \PY{o}{=} \PY{n}{parameters}\PY{p}{[}\PY{l+s+s2}{\PYZdq{}}\PY{l+s+s2}{Waa}\PY{l+s+s2}{\PYZdq{}}\PY{p}{]}
    \PY{n}{Wya} \PY{o}{=} \PY{n}{parameters}\PY{p}{[}\PY{l+s+s2}{\PYZdq{}}\PY{l+s+s2}{Wya}\PY{l+s+s2}{\PYZdq{}}\PY{p}{]}
    \PY{n}{ba} \PY{o}{=} \PY{n}{parameters}\PY{p}{[}\PY{l+s+s2}{\PYZdq{}}\PY{l+s+s2}{ba}\PY{l+s+s2}{\PYZdq{}}\PY{p}{]}
    \PY{n}{by} \PY{o}{=} \PY{n}{parameters}\PY{p}{[}\PY{l+s+s2}{\PYZdq{}}\PY{l+s+s2}{by}\PY{l+s+s2}{\PYZdq{}}\PY{p}{]}
    
    \PY{c+c1}{\PYZsh{}\PYZsh{}\PYZsh{} START CODE HERE \PYZsh{}\PYZsh{}\PYZsh{} (≈2 lines)}
    \PY{c+c1}{\PYZsh{} compute next activation state using the formula given above}
    \PY{n}{a\PYZus{}next} \PY{o}{=} \PY{n}{np}\PY{o}{.}\PY{n}{tanh}\PY{p}{(}\PY{n}{np}\PY{o}{.}\PY{n}{matmul}\PY{p}{(}\PY{n}{Wax}\PY{p}{,} \PY{n}{xt}\PY{p}{)} \PY{o}{+} \PY{n}{np}\PY{o}{.}\PY{n}{matmul}\PY{p}{(}\PY{n}{Waa}\PY{p}{,} \PY{n}{a\PYZus{}prev}\PY{p}{)} \PY{o}{+} \PY{n}{ba}\PY{p}{)}
    \PY{c+c1}{\PYZsh{} compute output of the current cell using the formula given above}
    \PY{n}{yt\PYZus{}pred} \PY{o}{=} \PY{n}{softmax}\PY{p}{(}\PY{n}{np}\PY{o}{.}\PY{n}{matmul}\PY{p}{(}\PY{n}{Wya}\PY{p}{,} \PY{n}{a\PYZus{}next}\PY{p}{)} \PY{o}{+} \PY{n}{by}\PY{p}{)}
    \PY{c+c1}{\PYZsh{}\PYZsh{}\PYZsh{} END CODE HERE \PYZsh{}\PYZsh{}\PYZsh{}}
    
    \PY{c+c1}{\PYZsh{} store values you need for backward propagation in cache}
    \PY{n}{cache} \PY{o}{=} \PY{p}{(}\PY{n}{a\PYZus{}next}\PY{p}{,} \PY{n}{a\PYZus{}prev}\PY{p}{,} \PY{n}{xt}\PY{p}{,} \PY{n}{parameters}\PY{p}{)}
    
    \PY{k}{return} \PY{n}{a\PYZus{}next}\PY{p}{,} \PY{n}{yt\PYZus{}pred}\PY{p}{,} \PY{n}{cache}
\end{Verbatim}
\end{tcolorbox}

    \begin{tcolorbox}[breakable, size=fbox, boxrule=1pt, pad at break*=1mm,colback=cellbackground, colframe=cellborder]
\prompt{In}{incolor}{14}{\boxspacing}
\begin{Verbatim}[commandchars=\\\{\}]
\PY{n}{np}\PY{o}{.}\PY{n}{random}\PY{o}{.}\PY{n}{seed}\PY{p}{(}\PY{l+m+mi}{1}\PY{p}{)}
\PY{n}{xt\PYZus{}tmp} \PY{o}{=} \PY{n}{np}\PY{o}{.}\PY{n}{random}\PY{o}{.}\PY{n}{randn}\PY{p}{(}\PY{l+m+mi}{3}\PY{p}{,} \PY{l+m+mi}{10}\PY{p}{)}
\PY{n}{a\PYZus{}prev\PYZus{}tmp} \PY{o}{=} \PY{n}{np}\PY{o}{.}\PY{n}{random}\PY{o}{.}\PY{n}{randn}\PY{p}{(}\PY{l+m+mi}{5}\PY{p}{,} \PY{l+m+mi}{10}\PY{p}{)}
\PY{n}{parameters\PYZus{}tmp} \PY{o}{=} \PY{p}{\PYZob{}}\PY{p}{\PYZcb{}}
\PY{n}{parameters\PYZus{}tmp}\PY{p}{[}\PY{l+s+s1}{\PYZsq{}}\PY{l+s+s1}{Waa}\PY{l+s+s1}{\PYZsq{}}\PY{p}{]} \PY{o}{=} \PY{n}{np}\PY{o}{.}\PY{n}{random}\PY{o}{.}\PY{n}{randn}\PY{p}{(}\PY{l+m+mi}{5}\PY{p}{,} \PY{l+m+mi}{5}\PY{p}{)}
\PY{n}{parameters\PYZus{}tmp}\PY{p}{[}\PY{l+s+s1}{\PYZsq{}}\PY{l+s+s1}{Wax}\PY{l+s+s1}{\PYZsq{}}\PY{p}{]} \PY{o}{=} \PY{n}{np}\PY{o}{.}\PY{n}{random}\PY{o}{.}\PY{n}{randn}\PY{p}{(}\PY{l+m+mi}{5}\PY{p}{,} \PY{l+m+mi}{3}\PY{p}{)}
\PY{n}{parameters\PYZus{}tmp}\PY{p}{[}\PY{l+s+s1}{\PYZsq{}}\PY{l+s+s1}{Wya}\PY{l+s+s1}{\PYZsq{}}\PY{p}{]} \PY{o}{=} \PY{n}{np}\PY{o}{.}\PY{n}{random}\PY{o}{.}\PY{n}{randn}\PY{p}{(}\PY{l+m+mi}{2}\PY{p}{,} \PY{l+m+mi}{5}\PY{p}{)}
\PY{n}{parameters\PYZus{}tmp}\PY{p}{[}\PY{l+s+s1}{\PYZsq{}}\PY{l+s+s1}{ba}\PY{l+s+s1}{\PYZsq{}}\PY{p}{]} \PY{o}{=} \PY{n}{np}\PY{o}{.}\PY{n}{random}\PY{o}{.}\PY{n}{randn}\PY{p}{(}\PY{l+m+mi}{5}\PY{p}{,} \PY{l+m+mi}{1}\PY{p}{)}
\PY{n}{parameters\PYZus{}tmp}\PY{p}{[}\PY{l+s+s1}{\PYZsq{}}\PY{l+s+s1}{by}\PY{l+s+s1}{\PYZsq{}}\PY{p}{]} \PY{o}{=} \PY{n}{np}\PY{o}{.}\PY{n}{random}\PY{o}{.}\PY{n}{randn}\PY{p}{(}\PY{l+m+mi}{2}\PY{p}{,} \PY{l+m+mi}{1}\PY{p}{)}

\PY{n}{a\PYZus{}next\PYZus{}tmp}\PY{p}{,} \PY{n}{yt\PYZus{}pred\PYZus{}tmp}\PY{p}{,} \PY{n}{cache\PYZus{}tmp} \PY{o}{=} \PY{n}{rnn\PYZus{}cell\PYZus{}forward}\PY{p}{(}\PY{n}{xt\PYZus{}tmp}\PY{p}{,} \PY{n}{a\PYZus{}prev\PYZus{}tmp}\PY{p}{,} \PY{n}{parameters\PYZus{}tmp}\PY{p}{)}
\PY{n+nb}{print}\PY{p}{(}\PY{l+s+s2}{\PYZdq{}}\PY{l+s+s2}{a\PYZus{}next[4] = }\PY{l+s+se}{\PYZbs{}n}\PY{l+s+s2}{\PYZdq{}}\PY{p}{,} \PY{n}{a\PYZus{}next\PYZus{}tmp}\PY{p}{[}\PY{l+m+mi}{4}\PY{p}{]}\PY{p}{)}
\PY{n+nb}{print}\PY{p}{(}\PY{l+s+s2}{\PYZdq{}}\PY{l+s+s2}{a\PYZus{}next.shape = }\PY{l+s+se}{\PYZbs{}n}\PY{l+s+s2}{\PYZdq{}}\PY{p}{,} \PY{n}{a\PYZus{}next\PYZus{}tmp}\PY{o}{.}\PY{n}{shape}\PY{p}{)}
\PY{n+nb}{print}\PY{p}{(}\PY{l+s+s2}{\PYZdq{}}\PY{l+s+s2}{yt\PYZus{}pred[1] =}\PY{l+s+se}{\PYZbs{}n}\PY{l+s+s2}{\PYZdq{}}\PY{p}{,} \PY{n}{yt\PYZus{}pred\PYZus{}tmp}\PY{p}{[}\PY{l+m+mi}{1}\PY{p}{]}\PY{p}{)}
\PY{n+nb}{print}\PY{p}{(}\PY{l+s+s2}{\PYZdq{}}\PY{l+s+s2}{yt\PYZus{}pred.shape = }\PY{l+s+se}{\PYZbs{}n}\PY{l+s+s2}{\PYZdq{}}\PY{p}{,} \PY{n}{yt\PYZus{}pred\PYZus{}tmp}\PY{o}{.}\PY{n}{shape}\PY{p}{)}

\PY{c+c1}{\PYZsh{} UNIT TESTS}
\PY{n}{rnn\PYZus{}cell\PYZus{}forward\PYZus{}tests}\PY{p}{(}\PY{n}{rnn\PYZus{}cell\PYZus{}forward}\PY{p}{)}
\end{Verbatim}
\end{tcolorbox}

    \begin{Verbatim}[commandchars=\\\{\}]
a\_next[4] =
 [ 0.59584544  0.18141802  0.61311866  0.99808218  0.85016201  0.99980978
 -0.18887155  0.99815551  0.6531151   0.82872037]
a\_next.shape =
 (5, 10)
yt\_pred[1] =
 [0.9888161  0.01682021 0.21140899 0.36817467 0.98988387 0.88945212
 0.36920224 0.9966312  0.9982559  0.17746526]
yt\_pred.shape =
 (2, 10)
\textcolor{ansi-green-intense}{All tests passed}
    \end{Verbatim}

    \textbf{Expected Output}:

\begin{Shaded}
\begin{Highlighting}[]
\NormalTok{a\_next[}\DecValTok{4}\NormalTok{] }\OperatorTok{=} 
\NormalTok{ [ }\FloatTok{0.59584544}  \FloatTok{0.18141802}  \FloatTok{0.61311866}  \FloatTok{0.99808218}  \FloatTok{0.85016201}  \FloatTok{0.99980978}
 \OperatorTok{{-}}\FloatTok{0.18887155}  \FloatTok{0.99815551}  \FloatTok{0.6531151}   \FloatTok{0.82872037}\NormalTok{]}
\NormalTok{a\_next.shape }\OperatorTok{=} 
\NormalTok{ (}\DecValTok{5}\NormalTok{, }\DecValTok{10}\NormalTok{)}
\NormalTok{yt\_pred[}\DecValTok{1}\NormalTok{] }\OperatorTok{=}
\NormalTok{ [ }\FloatTok{0.9888161}   \FloatTok{0.01682021}  \FloatTok{0.21140899}  \FloatTok{0.36817467}  \FloatTok{0.98988387}  \FloatTok{0.88945212}
  \FloatTok{0.36920224}  \FloatTok{0.9966312}   \FloatTok{0.9982559}   \FloatTok{0.17746526}\NormalTok{]}
\NormalTok{yt\_pred.shape }\OperatorTok{=} 
\NormalTok{ (}\DecValTok{2}\NormalTok{, }\DecValTok{10}\NormalTok{)}
\end{Highlighting}
\end{Shaded}

    \#\#\# 1.2 - RNN Forward Pass

\begin{itemize}
\tightlist
\item
  A recurrent neural network (RNN) is a repetition of the RNN cell that
  you've just built.

  \begin{itemize}
  \tightlist
  \item
    If your input sequence of data is 10 time steps long, then you will
    re-use the RNN cell 10 times
  \end{itemize}
\item
  Each cell takes two inputs at each time step:

  \begin{itemize}
  \tightlist
  \item
    \(a^{\langle t-1 \rangle}\): The hidden state from the previous cell
  \item
    \(x^{\langle t \rangle}\): The current time step's input data
  \end{itemize}
\item
  It has two outputs at each time step:

  \begin{itemize}
  \tightlist
  \item
    A hidden state (\(a^{\langle t \rangle}\))
  \item
    A prediction (\(y^{\langle t \rangle}\))
  \end{itemize}
\item
  The weights and biases \((W_{aa}, W_{ax}, b_{a}, W_{ay}, b_{y})\) are
  re-used each time step

  \begin{itemize}
  \tightlist
  \item
    They are maintained between calls to \texttt{rnn\_cell\_forward} in
    the `parameters' dictionary
  \end{itemize}
\end{itemize}

Figure 3: Basic RNN. The input sequence
\(x = (x^{\langle 1 \rangle}, x^{\langle 2 \rangle}, ..., x^{\langle T_x \rangle})\)
is carried over \(T_x\) time steps. The network outputs
\(y = (y^{\langle 1 \rangle}, y^{\langle 2 \rangle}, ..., y^{\langle T_x \rangle})\).

    \#\#\# Exercise 2 - rnn\_forward

Implement the forward propagation of the RNN described in Figure 3.

\textbf{Instructions}: * Create a 3D array of zeros, \(a\) of shape
\((n_{a}, m, T_{x})\) that will store all the hidden states computed by
the RNN * Create a 3D array of zeros, \(\hat{y}\), of shape
\((n_{y}, m, T_{x})\) that will store the predictions\\
- Note that in this case, \(T_{y} = T_{x}\) (the prediction and input
have the same number of time steps) * Initialize the 2D hidden state
\texttt{a\_next} by setting it equal to the initial hidden state,
\(a_{0}\) * At each time step \(t\): - Get \(x^{\langle t \rangle}\),
which is a 2D slice of \(x\) for a single time step \(t\) -
\(x^{\langle t \rangle}\) has shape \((n_{x}, m)\) - \(x\) has shape
\((n_{x}, m, T_{x})\) - Update the 2D hidden state
\(a^{\langle t \rangle}\) (variable name \texttt{a\_next}), the
prediction \(\hat{y}^{\langle t \rangle}\) and the cache by running
\texttt{rnn\_cell\_forward} - \(a^{\langle t \rangle}\) has shape
\((n_{a}, m)\) - Store the 2D hidden state in the 3D tensor \(a\), at
the \(t^{th}\) position - \(a\) has shape \((n_{a}, m, T_{x})\) - Store
the 2D \(\hat{y}^{\langle t \rangle}\) prediction (variable name
\texttt{yt\_pred}) in the 3D tensor \(\hat{y}_{pred}\) at the \(t^{th}\)
position - \(\hat{y}^{\langle t \rangle}\) has shape \((n_{y}, m)\) -
\(\hat{y}\) has shape \((n_{y}, m, T_x)\) - Append the cache to the list
of caches * Return the 3D tensor \(a\) and \(\hat{y}\), as well as the
list of caches

\hypertarget{additional-hints}{%
\paragraph{Additional Hints}\label{additional-hints}}

\begin{itemize}
\tightlist
\item
  Some helpful documentation on
  \href{https://docs.scipy.org/doc/numpy/reference/generated/numpy.zeros.html}{np.zeros}
\item
  If you have a 3 dimensional numpy array and are indexing by its third
  dimension, you can use array slicing like this:
  \texttt{var\_name{[}:,:,i{]}}
\end{itemize}

    \begin{tcolorbox}[breakable, size=fbox, boxrule=1pt, pad at break*=1mm,colback=cellbackground, colframe=cellborder]
\prompt{In}{incolor}{17}{\boxspacing}
\begin{Verbatim}[commandchars=\\\{\}]
\PY{c+c1}{\PYZsh{} UNQ\PYZus{}C2 (UNIQUE CELL IDENTIFIER, DO NOT EDIT)}
\PY{c+c1}{\PYZsh{} GRADED FUNCTION: rnn\PYZus{}forward}

\PY{k}{def} \PY{n+nf}{rnn\PYZus{}forward}\PY{p}{(}\PY{n}{x}\PY{p}{,} \PY{n}{a0}\PY{p}{,} \PY{n}{parameters}\PY{p}{)}\PY{p}{:}
    \PY{l+s+sd}{\PYZdq{}\PYZdq{}\PYZdq{}}
\PY{l+s+sd}{    Implement the forward propagation of the recurrent neural network described in Figure (3).}

\PY{l+s+sd}{    Arguments:}
\PY{l+s+sd}{    x \PYZhy{}\PYZhy{} Input data for every time\PYZhy{}step, of shape (n\PYZus{}x, m, T\PYZus{}x).}
\PY{l+s+sd}{    a0 \PYZhy{}\PYZhy{} Initial hidden state, of shape (n\PYZus{}a, m)}
\PY{l+s+sd}{    parameters \PYZhy{}\PYZhy{} python dictionary containing:}
\PY{l+s+sd}{                        Waa \PYZhy{}\PYZhy{} Weight matrix multiplying the hidden state, numpy array of shape (n\PYZus{}a, n\PYZus{}a)}
\PY{l+s+sd}{                        Wax \PYZhy{}\PYZhy{} Weight matrix multiplying the input, numpy array of shape (n\PYZus{}a, n\PYZus{}x)}
\PY{l+s+sd}{                        Wya \PYZhy{}\PYZhy{} Weight matrix relating the hidden\PYZhy{}state to the output, numpy array of shape (n\PYZus{}y, n\PYZus{}a)}
\PY{l+s+sd}{                        ba \PYZhy{}\PYZhy{}  Bias numpy array of shape (n\PYZus{}a, 1)}
\PY{l+s+sd}{                        by \PYZhy{}\PYZhy{} Bias relating the hidden\PYZhy{}state to the output, numpy array of shape (n\PYZus{}y, 1)}

\PY{l+s+sd}{    Returns:}
\PY{l+s+sd}{    a \PYZhy{}\PYZhy{} Hidden states for every time\PYZhy{}step, numpy array of shape (n\PYZus{}a, m, T\PYZus{}x)}
\PY{l+s+sd}{    y\PYZus{}pred \PYZhy{}\PYZhy{} Predictions for every time\PYZhy{}step, numpy array of shape (n\PYZus{}y, m, T\PYZus{}x)}
\PY{l+s+sd}{    caches \PYZhy{}\PYZhy{} tuple of values needed for the backward pass, contains (list of caches, x)}
\PY{l+s+sd}{    \PYZdq{}\PYZdq{}\PYZdq{}}
    
    \PY{c+c1}{\PYZsh{} Initialize \PYZdq{}caches\PYZdq{} which will contain the list of all caches}
    \PY{n}{caches} \PY{o}{=} \PY{p}{[}\PY{p}{]}
    
    \PY{c+c1}{\PYZsh{} Retrieve dimensions from shapes of x and parameters[\PYZdq{}Wya\PYZdq{}]}
    \PY{n}{n\PYZus{}x}\PY{p}{,} \PY{n}{m}\PY{p}{,} \PY{n}{T\PYZus{}x} \PY{o}{=} \PY{n}{x}\PY{o}{.}\PY{n}{shape}
    \PY{n}{n\PYZus{}y}\PY{p}{,} \PY{n}{n\PYZus{}a} \PY{o}{=} \PY{n}{parameters}\PY{p}{[}\PY{l+s+s2}{\PYZdq{}}\PY{l+s+s2}{Wya}\PY{l+s+s2}{\PYZdq{}}\PY{p}{]}\PY{o}{.}\PY{n}{shape}
    
    \PY{c+c1}{\PYZsh{}\PYZsh{}\PYZsh{} START CODE HERE \PYZsh{}\PYZsh{}\PYZsh{}}
    
    \PY{c+c1}{\PYZsh{} initialize \PYZdq{}a\PYZdq{} and \PYZdq{}y\PYZus{}pred\PYZdq{} with zeros (≈2 lines)}
    \PY{n}{a} \PY{o}{=} \PY{n}{np}\PY{o}{.}\PY{n}{zeros}\PY{p}{(}\PY{p}{(}\PY{n}{n\PYZus{}a}\PY{p}{,} \PY{n}{m}\PY{p}{,} \PY{n}{T\PYZus{}x}\PY{p}{)}\PY{p}{)}
    \PY{n}{y\PYZus{}pred} \PY{o}{=} \PY{n}{np}\PY{o}{.}\PY{n}{zeros}\PY{p}{(}\PY{p}{(}\PY{n}{n\PYZus{}y}\PY{p}{,} \PY{n}{m}\PY{p}{,} \PY{n}{T\PYZus{}x}\PY{p}{)}\PY{p}{)}
    
    \PY{c+c1}{\PYZsh{} Initialize a\PYZus{}next (≈1 line)}
    \PY{n}{a\PYZus{}next} \PY{o}{=} \PY{n}{a0}
    
    \PY{c+c1}{\PYZsh{} loop over all time\PYZhy{}steps}
    \PY{k}{for} \PY{n}{t} \PY{o+ow}{in} \PY{n+nb}{range}\PY{p}{(}\PY{n}{T\PYZus{}x}\PY{p}{)}\PY{p}{:}
        \PY{c+c1}{\PYZsh{} Update next hidden state, compute the prediction, get the cache (≈1 line)}
        \PY{n}{a\PYZus{}next}\PY{p}{,} \PY{n}{yt\PYZus{}pred}\PY{p}{,} \PY{n}{cache} \PY{o}{=} \PY{n}{rnn\PYZus{}cell\PYZus{}forward}\PY{p}{(}\PY{n}{x}\PY{p}{[}\PY{p}{:}\PY{p}{,}\PY{p}{:}\PY{p}{,}\PY{n}{t}\PY{p}{]}\PY{p}{,} \PY{n}{a\PYZus{}next}\PY{p}{,} \PY{n}{parameters}\PY{p}{)}
        \PY{c+c1}{\PYZsh{} Save the value of the new \PYZdq{}next\PYZdq{} hidden state in a (≈1 line)}
        \PY{n}{a}\PY{p}{[}\PY{p}{:}\PY{p}{,}\PY{p}{:}\PY{p}{,}\PY{n}{t}\PY{p}{]} \PY{o}{=} \PY{n}{a\PYZus{}next}
        \PY{c+c1}{\PYZsh{} Save the value of the prediction in y (≈1 line)}
        \PY{n}{y\PYZus{}pred}\PY{p}{[}\PY{p}{:}\PY{p}{,}\PY{p}{:}\PY{p}{,}\PY{n}{t}\PY{p}{]} \PY{o}{=} \PY{n}{yt\PYZus{}pred}
        \PY{c+c1}{\PYZsh{} Append \PYZdq{}cache\PYZdq{} to \PYZdq{}caches\PYZdq{} (≈1 line)}
        \PY{n}{caches}\PY{o}{.}\PY{n}{append}\PY{p}{(}\PY{n}{cache}\PY{p}{)}
    \PY{c+c1}{\PYZsh{}\PYZsh{}\PYZsh{} END CODE HERE \PYZsh{}\PYZsh{}\PYZsh{}}
    
    \PY{c+c1}{\PYZsh{} store values needed for backward propagation in cache}
    \PY{n}{caches} \PY{o}{=} \PY{p}{(}\PY{n}{caches}\PY{p}{,} \PY{n}{x}\PY{p}{)}
    
    \PY{k}{return} \PY{n}{a}\PY{p}{,} \PY{n}{y\PYZus{}pred}\PY{p}{,} \PY{n}{caches}
\end{Verbatim}
\end{tcolorbox}

    \begin{tcolorbox}[breakable, size=fbox, boxrule=1pt, pad at break*=1mm,colback=cellbackground, colframe=cellborder]
\prompt{In}{incolor}{18}{\boxspacing}
\begin{Verbatim}[commandchars=\\\{\}]
\PY{n}{np}\PY{o}{.}\PY{n}{random}\PY{o}{.}\PY{n}{seed}\PY{p}{(}\PY{l+m+mi}{1}\PY{p}{)}
\PY{n}{x\PYZus{}tmp} \PY{o}{=} \PY{n}{np}\PY{o}{.}\PY{n}{random}\PY{o}{.}\PY{n}{randn}\PY{p}{(}\PY{l+m+mi}{3}\PY{p}{,} \PY{l+m+mi}{10}\PY{p}{,} \PY{l+m+mi}{4}\PY{p}{)}
\PY{n}{a0\PYZus{}tmp} \PY{o}{=} \PY{n}{np}\PY{o}{.}\PY{n}{random}\PY{o}{.}\PY{n}{randn}\PY{p}{(}\PY{l+m+mi}{5}\PY{p}{,} \PY{l+m+mi}{10}\PY{p}{)}
\PY{n}{parameters\PYZus{}tmp} \PY{o}{=} \PY{p}{\PYZob{}}\PY{p}{\PYZcb{}}
\PY{n}{parameters\PYZus{}tmp}\PY{p}{[}\PY{l+s+s1}{\PYZsq{}}\PY{l+s+s1}{Waa}\PY{l+s+s1}{\PYZsq{}}\PY{p}{]} \PY{o}{=} \PY{n}{np}\PY{o}{.}\PY{n}{random}\PY{o}{.}\PY{n}{randn}\PY{p}{(}\PY{l+m+mi}{5}\PY{p}{,} \PY{l+m+mi}{5}\PY{p}{)}
\PY{n}{parameters\PYZus{}tmp}\PY{p}{[}\PY{l+s+s1}{\PYZsq{}}\PY{l+s+s1}{Wax}\PY{l+s+s1}{\PYZsq{}}\PY{p}{]} \PY{o}{=} \PY{n}{np}\PY{o}{.}\PY{n}{random}\PY{o}{.}\PY{n}{randn}\PY{p}{(}\PY{l+m+mi}{5}\PY{p}{,} \PY{l+m+mi}{3}\PY{p}{)}
\PY{n}{parameters\PYZus{}tmp}\PY{p}{[}\PY{l+s+s1}{\PYZsq{}}\PY{l+s+s1}{Wya}\PY{l+s+s1}{\PYZsq{}}\PY{p}{]} \PY{o}{=} \PY{n}{np}\PY{o}{.}\PY{n}{random}\PY{o}{.}\PY{n}{randn}\PY{p}{(}\PY{l+m+mi}{2}\PY{p}{,} \PY{l+m+mi}{5}\PY{p}{)}
\PY{n}{parameters\PYZus{}tmp}\PY{p}{[}\PY{l+s+s1}{\PYZsq{}}\PY{l+s+s1}{ba}\PY{l+s+s1}{\PYZsq{}}\PY{p}{]} \PY{o}{=} \PY{n}{np}\PY{o}{.}\PY{n}{random}\PY{o}{.}\PY{n}{randn}\PY{p}{(}\PY{l+m+mi}{5}\PY{p}{,} \PY{l+m+mi}{1}\PY{p}{)}
\PY{n}{parameters\PYZus{}tmp}\PY{p}{[}\PY{l+s+s1}{\PYZsq{}}\PY{l+s+s1}{by}\PY{l+s+s1}{\PYZsq{}}\PY{p}{]} \PY{o}{=} \PY{n}{np}\PY{o}{.}\PY{n}{random}\PY{o}{.}\PY{n}{randn}\PY{p}{(}\PY{l+m+mi}{2}\PY{p}{,} \PY{l+m+mi}{1}\PY{p}{)}

\PY{n}{a\PYZus{}tmp}\PY{p}{,} \PY{n}{y\PYZus{}pred\PYZus{}tmp}\PY{p}{,} \PY{n}{caches\PYZus{}tmp} \PY{o}{=} \PY{n}{rnn\PYZus{}forward}\PY{p}{(}\PY{n}{x\PYZus{}tmp}\PY{p}{,} \PY{n}{a0\PYZus{}tmp}\PY{p}{,} \PY{n}{parameters\PYZus{}tmp}\PY{p}{)}
\PY{n+nb}{print}\PY{p}{(}\PY{l+s+s2}{\PYZdq{}}\PY{l+s+s2}{a[4][1] = }\PY{l+s+se}{\PYZbs{}n}\PY{l+s+s2}{\PYZdq{}}\PY{p}{,} \PY{n}{a\PYZus{}tmp}\PY{p}{[}\PY{l+m+mi}{4}\PY{p}{]}\PY{p}{[}\PY{l+m+mi}{1}\PY{p}{]}\PY{p}{)}
\PY{n+nb}{print}\PY{p}{(}\PY{l+s+s2}{\PYZdq{}}\PY{l+s+s2}{a.shape = }\PY{l+s+se}{\PYZbs{}n}\PY{l+s+s2}{\PYZdq{}}\PY{p}{,} \PY{n}{a\PYZus{}tmp}\PY{o}{.}\PY{n}{shape}\PY{p}{)}
\PY{n+nb}{print}\PY{p}{(}\PY{l+s+s2}{\PYZdq{}}\PY{l+s+s2}{y\PYZus{}pred[1][3] =}\PY{l+s+se}{\PYZbs{}n}\PY{l+s+s2}{\PYZdq{}}\PY{p}{,} \PY{n}{y\PYZus{}pred\PYZus{}tmp}\PY{p}{[}\PY{l+m+mi}{1}\PY{p}{]}\PY{p}{[}\PY{l+m+mi}{3}\PY{p}{]}\PY{p}{)}
\PY{n+nb}{print}\PY{p}{(}\PY{l+s+s2}{\PYZdq{}}\PY{l+s+s2}{y\PYZus{}pred.shape = }\PY{l+s+se}{\PYZbs{}n}\PY{l+s+s2}{\PYZdq{}}\PY{p}{,} \PY{n}{y\PYZus{}pred\PYZus{}tmp}\PY{o}{.}\PY{n}{shape}\PY{p}{)}
\PY{n+nb}{print}\PY{p}{(}\PY{l+s+s2}{\PYZdq{}}\PY{l+s+s2}{caches[1][1][3] =}\PY{l+s+se}{\PYZbs{}n}\PY{l+s+s2}{\PYZdq{}}\PY{p}{,} \PY{n}{caches\PYZus{}tmp}\PY{p}{[}\PY{l+m+mi}{1}\PY{p}{]}\PY{p}{[}\PY{l+m+mi}{1}\PY{p}{]}\PY{p}{[}\PY{l+m+mi}{3}\PY{p}{]}\PY{p}{)}
\PY{n+nb}{print}\PY{p}{(}\PY{l+s+s2}{\PYZdq{}}\PY{l+s+s2}{len(caches) = }\PY{l+s+se}{\PYZbs{}n}\PY{l+s+s2}{\PYZdq{}}\PY{p}{,} \PY{n+nb}{len}\PY{p}{(}\PY{n}{caches\PYZus{}tmp}\PY{p}{)}\PY{p}{)}

\PY{c+c1}{\PYZsh{}UNIT TEST    }
\PY{n}{rnn\PYZus{}forward\PYZus{}test}\PY{p}{(}\PY{n}{rnn\PYZus{}forward}\PY{p}{)}
\end{Verbatim}
\end{tcolorbox}

    \begin{Verbatim}[commandchars=\\\{\}]
a[4][1] =
 [-0.99999375  0.77911235 -0.99861469 -0.99833267]
a.shape =
 (5, 10, 4)
y\_pred[1][3] =
 [0.79560373 0.86224861 0.11118257 0.81515947]
y\_pred.shape =
 (2, 10, 4)
caches[1][1][3] =
 [-1.1425182  -0.34934272 -0.20889423  0.58662319]
len(caches) =
 2
\textcolor{ansi-green-intense}{All tests passed}
    \end{Verbatim}

    \textbf{Expected Output}:

\begin{Shaded}
\begin{Highlighting}[]
\NormalTok{a[}\DecValTok{4}\NormalTok{][}\DecValTok{1}\NormalTok{] }\OperatorTok{=} 
\NormalTok{ [}\OperatorTok{{-}}\FloatTok{0.99999375}  \FloatTok{0.77911235} \OperatorTok{{-}}\FloatTok{0.99861469} \OperatorTok{{-}}\FloatTok{0.99833267}\NormalTok{]}
\NormalTok{a.shape }\OperatorTok{=} 
\NormalTok{ (}\DecValTok{5}\NormalTok{, }\DecValTok{10}\NormalTok{, }\DecValTok{4}\NormalTok{)}
\NormalTok{y\_pred[}\DecValTok{1}\NormalTok{][}\DecValTok{3}\NormalTok{] }\OperatorTok{=}
\NormalTok{ [ }\FloatTok{0.79560373}  \FloatTok{0.86224861}  \FloatTok{0.11118257}  \FloatTok{0.81515947}\NormalTok{]}
\NormalTok{y\_pred.shape }\OperatorTok{=} 
\NormalTok{ (}\DecValTok{2}\NormalTok{, }\DecValTok{10}\NormalTok{, }\DecValTok{4}\NormalTok{)}
\NormalTok{caches[}\DecValTok{1}\NormalTok{][}\DecValTok{1}\NormalTok{][}\DecValTok{3}\NormalTok{] }\OperatorTok{=}
\NormalTok{ [}\OperatorTok{{-}}\FloatTok{1.1425182}  \OperatorTok{{-}}\FloatTok{0.34934272} \OperatorTok{{-}}\FloatTok{0.20889423}  \FloatTok{0.58662319}\NormalTok{]}
\BuiltInTok{len}\NormalTok{(caches) }\OperatorTok{=} 
 \DecValTok{2}
\end{Highlighting}
\end{Shaded}

    \hypertarget{congratulations}{%
\subsubsection{Congratulations!}\label{congratulations}}

You've successfully built the forward propagation of a recurrent neural
network from scratch. Nice work!

\hypertarget{situations-when-this-rnn-will-perform-better}{%
\paragraph{Situations when this RNN will perform
better:}\label{situations-when-this-rnn-will-perform-better}}

\begin{itemize}
\tightlist
\item
  This will work well enough for some applications, but it suffers from
  vanishing gradients.
\item
  The RNN works best when each output \(\hat{y}^{\langle t \rangle}\)
  can be estimated using ``local'' context.\\
\item
  ``Local'' context refers to information that is close to the
  prediction's time step \(t\).
\item
  More formally, local context refers to inputs
  \(x^{\langle t' \rangle}\) and predictions
  \(\hat{y}^{\langle t \rangle}\) where \(t'\) is close to \(t\).
\end{itemize}

What you should remember: * The recurrent neural network, or RNN, is
essentially the repeated use of a single cell. * A basic RNN reads
inputs one at a time, and remembers information through the hidden layer
activations (hidden states) that are passed from one time step to the
next. * The time step dimension determines how many times to re-use the
RNN cell * Each cell takes two inputs at each time step: * The hidden
state from the previous cell * The current time step's input data * Each
cell has two outputs at each time step: * A hidden state * A prediction

In the next section, you'll build a more complex model, the LSTM, which
is better at addressing vanishing gradients. The LSTM is better able to
remember a piece of information and save it for many time steps.

    \#\# 2 - Long Short-Term Memory (LSTM) Network

The following figure shows the operations of an LSTM cell:

Figure 4: LSTM cell. This tracks and updates a ``cell state,'' or memory
variable \(c^{\langle t \rangle}\) at every time step, which can be
different from \(a^{\langle t \rangle}\).\\
Note, the \(softmax^{}\) includes a dense layer and softmax.

Similar to the RNN example above, you'll begin by implementing the LSTM
cell for a single time step. Then, you'll iteratively call it from
inside a ``for loop'' to have it process an input with \(T_x\) time
steps.

    \hypertarget{overview-of-gates-and-states}{%
\subsubsection{Overview of gates and
states}\label{overview-of-gates-and-states}}

\hypertarget{forget-gate-mathbfgamma_f}{%
\paragraph{\texorpdfstring{Forget gate
\(\mathbf{\Gamma}_{f}\)}{Forget gate \textbackslash mathbf\{\textbackslash Gamma\}\_\{f\}}}\label{forget-gate-mathbfgamma_f}}

\begin{itemize}
\tightlist
\item
  Let's assume you are reading words in a piece of text, and plan to use
  an LSTM to keep track of grammatical structures, such as whether the
  subject is singular (``puppy'') or plural (``puppies'').
\item
  If the subject changes its state (from a singular word to a plural
  word), the memory of the previous state becomes outdated, so you'll
  ``forget'' that outdated state.
\item
  The ``forget gate'' is a tensor containing values between 0 and 1.

  \begin{itemize}
  \tightlist
  \item
    If a unit in the forget gate has a value close to 0, the LSTM will
    ``forget'' the stored state in the corresponding unit of the
    previous cell state.
  \item
    If a unit in the forget gate has a value close to 1, the LSTM will
    mostly remember the corresponding value in the stored state.
  \end{itemize}
\end{itemize}

\hypertarget{equation}{%
\subparagraph{Equation}\label{equation}}

\[\mathbf{\Gamma}_f^{\langle t \rangle} = \sigma(\mathbf{W}_f[\mathbf{a}^{\langle t-1 \rangle}, \mathbf{x}^{\langle t \rangle}] + \mathbf{b}_f)\tag{1} \]

\hypertarget{explanation-of-the-equation}{%
\subparagraph{Explanation of the
equation:}\label{explanation-of-the-equation}}

\begin{itemize}
\tightlist
\item
  \(\mathbf{W_{f}}\) contains weights that govern the forget gate's
  behavior.
\item
  The previous time step's hidden state \([a^{\langle t-1 \rangle}\) and
  current time step's input \(x^{\langle t \rangle}]\) are concatenated
  together and multiplied by \(\mathbf{W_{f}}\).
\item
  A sigmoid function is used to make each of the gate tensor's values
  \(\mathbf{\Gamma}_f^{\langle t \rangle}\) range from 0 to 1.
\item
  The forget gate \(\mathbf{\Gamma}_f^{\langle t \rangle}\) has the same
  dimensions as the previous cell state \(c^{\langle t-1 \rangle}\).
\item
  This means that the two can be multiplied together, element-wise.
\item
  Multiplying the tensors
  \(\mathbf{\Gamma}_f^{\langle t \rangle} * \mathbf{c}^{\langle t-1 \rangle}\)
  is like applying a mask over the previous cell state.
\item
  If a single value in \(\mathbf{\Gamma}_f^{\langle t \rangle}\) is 0 or
  close to 0, then the product is close to 0.

  \begin{itemize}
  \tightlist
  \item
    This keeps the information stored in the corresponding unit in
    \(\mathbf{c}^{\langle t-1 \rangle}\) from being remembered for the
    next time step.
  \end{itemize}
\item
  Similarly, if one value is close to 1, the product is close to the
  original value in the previous cell state.

  \begin{itemize}
  \tightlist
  \item
    The LSTM will keep the information from the corresponding unit of
    \(\mathbf{c}^{\langle t-1 \rangle}\), to be used in the next time
    step.
  \end{itemize}
\end{itemize}

\hypertarget{variable-names-in-the-code}{%
\subparagraph{Variable names in the
code}\label{variable-names-in-the-code}}

The variable names in the code are similar to the equations, with slight
differences.\\
* \texttt{Wf}: forget gate weight \(\mathbf{W}_{f}\) * \texttt{bf}:
forget gate bias \(\mathbf{b}_{f}\) * \texttt{ft}: forget gate
\(\Gamma_f^{\langle t \rangle}\)

    \hypertarget{candidate-value-tildemathbfclangle-t-rangle}{%
\paragraph{\texorpdfstring{Candidate value
\(\tilde{\mathbf{c}}^{\langle t \rangle}\)}{Candidate value \textbackslash tilde\{\textbackslash mathbf\{c\}\}\^{}\{\textbackslash langle t \textbackslash rangle\}}}\label{candidate-value-tildemathbfclangle-t-rangle}}

\begin{itemize}
\tightlist
\item
  The candidate value is a tensor containing information from the
  current time step that \textbf{may} be stored in the current cell
  state \(\mathbf{c}^{\langle t \rangle}\).
\item
  The parts of the candidate value that get passed on depend on the
  update gate.
\item
  The candidate value is a tensor containing values that range from -1
  to 1.
\item
  The tilde ``\textasciitilde{}'' is used to differentiate the candidate
  \(\tilde{\mathbf{c}}^{\langle t \rangle}\) from the cell state
  \(\mathbf{c}^{\langle t \rangle}\).
\end{itemize}

\hypertarget{equation}{%
\subparagraph{Equation}\label{equation}}

\[\mathbf{\tilde{c}}^{\langle t \rangle} = \tanh\left( \mathbf{W}_{c} [\mathbf{a}^{\langle t - 1 \rangle}, \mathbf{x}^{\langle t \rangle}] + \mathbf{b}_{c} \right) \tag{3}\]

\hypertarget{explanation-of-the-equation}{%
\subparagraph{Explanation of the
equation}\label{explanation-of-the-equation}}

\begin{itemize}
\tightlist
\item
  The \emph{tanh} function produces values between -1 and 1.
\end{itemize}

\hypertarget{variable-names-in-the-code}{%
\subparagraph{Variable names in the
code}\label{variable-names-in-the-code}}

\begin{itemize}
\tightlist
\item
  \texttt{cct}: candidate value
  \(\mathbf{\tilde{c}}^{\langle t \rangle}\)
\end{itemize}

    \hypertarget{update-gate-mathbfgamma_i}{%
\paragraph{\texorpdfstring{Update gate
\(\mathbf{\Gamma}_{i}\)}{Update gate \textbackslash mathbf\{\textbackslash Gamma\}\_\{i\}}}\label{update-gate-mathbfgamma_i}}

\begin{itemize}
\tightlist
\item
  You use the update gate to decide what aspects of the candidate
  \(\tilde{\mathbf{c}}^{\langle t \rangle}\) to add to the cell state
  \(c^{\langle t \rangle}\).
\item
  The update gate decides what parts of a ``candidate'' tensor
  \(\tilde{\mathbf{c}}^{\langle t \rangle}\) are passed onto the cell
  state \(\mathbf{c}^{\langle t \rangle}\).
\item
  The update gate is a tensor containing values between 0 and 1.

  \begin{itemize}
  \tightlist
  \item
    When a unit in the update gate is close to 1, it allows the value of
    the candidate \(\tilde{\mathbf{c}}^{\langle t \rangle}\) to be
    passed onto the hidden state \(\mathbf{c}^{\langle t \rangle}\)
  \item
    When a unit in the update gate is close to 0, it prevents the
    corresponding value in the candidate from being passed onto the
    hidden state.
  \end{itemize}
\item
  Notice that the subscript ``i'' is used and not ``u'', to follow the
  convention used in the literature.
\end{itemize}

\hypertarget{equation}{%
\subparagraph{Equation}\label{equation}}

\[\mathbf{\Gamma}_i^{\langle t \rangle} = \sigma(\mathbf{W}_i[a^{\langle t-1 \rangle}, \mathbf{x}^{\langle t \rangle}] + \mathbf{b}_i)\tag{2} \]

\hypertarget{explanation-of-the-equation}{%
\subparagraph{Explanation of the
equation}\label{explanation-of-the-equation}}

\begin{itemize}
\tightlist
\item
  Similar to the forget gate, here
  \(\mathbf{\Gamma}_i^{\langle t \rangle}\), the sigmoid produces values
  between 0 and 1.
\item
  The update gate is multiplied element-wise with the candidate, and
  this product
  (\(\mathbf{\Gamma}_{i}^{\langle t \rangle} * \tilde{c}^{\langle t \rangle}\))
  is used in determining the cell state
  \(\mathbf{c}^{\langle t \rangle}\).
\end{itemize}

\hypertarget{variable-names-in-code-please-note-that-theyre-different-than-the-equations}{%
\subparagraph{Variable names in code (Please note that they're different
than the
equations)}\label{variable-names-in-code-please-note-that-theyre-different-than-the-equations}}

In the code, you'll use the variable names found in the academic
literature. These variables don't use ``u'' to denote ``update''. *
\texttt{Wi} is the update gate weight \(\mathbf{W}_i\) (not ``Wu'') *
\texttt{bi} is the update gate bias \(\mathbf{b}_i\) (not ``bu'') *
\texttt{it} is the update gate \(\mathbf{\Gamma}_i^{\langle t \rangle}\)
(not ``ut'')

    \hypertarget{cell-state-mathbfclangle-t-rangle}{%
\paragraph{\texorpdfstring{Cell state
\(\mathbf{c}^{\langle t \rangle}\)}{Cell state \textbackslash mathbf\{c\}\^{}\{\textbackslash langle t \textbackslash rangle\}}}\label{cell-state-mathbfclangle-t-rangle}}

\begin{itemize}
\tightlist
\item
  The cell state is the ``memory'' that gets passed onto future time
  steps.
\item
  The new cell state \(\mathbf{c}^{\langle t \rangle}\) is a combination
  of the previous cell state and the candidate value.
\end{itemize}

\hypertarget{equation}{%
\subparagraph{Equation}\label{equation}}

\[ \mathbf{c}^{\langle t \rangle} = \mathbf{\Gamma}_f^{\langle t \rangle}* \mathbf{c}^{\langle t-1 \rangle} + \mathbf{\Gamma}_{i}^{\langle t \rangle} *\mathbf{\tilde{c}}^{\langle t \rangle} \tag{4} \]

\hypertarget{explanation-of-equation}{%
\subparagraph{Explanation of equation}\label{explanation-of-equation}}

\begin{itemize}
\tightlist
\item
  The previous cell state \(\mathbf{c}^{\langle t-1 \rangle}\) is
  adjusted (weighted) by the forget gate
  \(\mathbf{\Gamma}_{f}^{\langle t \rangle}\)
\item
  and the candidate value \(\tilde{\mathbf{c}}^{\langle t \rangle}\),
  adjusted (weighted) by the update gate
  \(\mathbf{\Gamma}_{i}^{\langle t \rangle}\)
\end{itemize}

\hypertarget{variable-names-and-shapes-in-the-code}{%
\subparagraph{Variable names and shapes in the
code}\label{variable-names-and-shapes-in-the-code}}

\begin{itemize}
\tightlist
\item
  \texttt{c}: cell state, including all time steps, \(\mathbf{c}\) shape
  \((n_{a}, m, T_x)\)
\item
  \texttt{c\_next}: new (next) cell state,
  \(\mathbf{c}^{\langle t \rangle}\) shape \((n_{a}, m)\)
\item
  \texttt{c\_prev}: previous cell state,
  \(\mathbf{c}^{\langle t-1 \rangle}\), shape \((n_{a}, m)\)
\end{itemize}

    \hypertarget{output-gate-mathbfgamma_o}{%
\paragraph{\texorpdfstring{Output gate
\(\mathbf{\Gamma}_{o}\)}{Output gate \textbackslash mathbf\{\textbackslash Gamma\}\_\{o\}}}\label{output-gate-mathbfgamma_o}}

\begin{itemize}
\tightlist
\item
  The output gate decides what gets sent as the prediction (output) of
  the time step.
\item
  The output gate is like the other gates, in that it contains values
  that range from 0 to 1.
\end{itemize}

\hypertarget{equation}{%
\subparagraph{Equation}\label{equation}}

\[ \mathbf{\Gamma}_o^{\langle t \rangle}=  \sigma(\mathbf{W}_o[\mathbf{a}^{\langle t-1 \rangle}, \mathbf{x}^{\langle t \rangle}] + \mathbf{b}_{o})\tag{5}\]

\hypertarget{explanation-of-the-equation}{%
\subparagraph{Explanation of the
equation}\label{explanation-of-the-equation}}

\begin{itemize}
\tightlist
\item
  The output gate is determined by the previous hidden state
  \(\mathbf{a}^{\langle t-1 \rangle}\) and the current input
  \(\mathbf{x}^{\langle t \rangle}\)
\item
  The sigmoid makes the gate range from 0 to 1.
\end{itemize}

\hypertarget{variable-names-in-the-code}{%
\subparagraph{Variable names in the
code}\label{variable-names-in-the-code}}

\begin{itemize}
\tightlist
\item
  \texttt{Wo}: output gate weight, \(\mathbf{W_o}\)
\item
  \texttt{bo}: output gate bias, \(\mathbf{b_o}\)
\item
  \texttt{ot}: output gate, \(\mathbf{\Gamma}_{o}^{\langle t \rangle}\)
\end{itemize}

    \hypertarget{hidden-state-mathbfalangle-t-rangle}{%
\paragraph{\texorpdfstring{Hidden state
\(\mathbf{a}^{\langle t \rangle}\)}{Hidden state \textbackslash mathbf\{a\}\^{}\{\textbackslash langle t \textbackslash rangle\}}}\label{hidden-state-mathbfalangle-t-rangle}}

\begin{itemize}
\tightlist
\item
  The hidden state gets passed to the LSTM cell's next time step.
\item
  It is used to determine the three gates
  (\(\mathbf{\Gamma}_{f}, \mathbf{\Gamma}_{u}, \mathbf{\Gamma}_{o}\)) of
  the next time step.
\item
  The hidden state is also used for the prediction
  \(y^{\langle t \rangle}\).
\end{itemize}

\hypertarget{equation}{%
\subparagraph{Equation}\label{equation}}

\[ \mathbf{a}^{\langle t \rangle} = \mathbf{\Gamma}_o^{\langle t \rangle} * \tanh(\mathbf{c}^{\langle t \rangle})\tag{6} \]

\hypertarget{explanation-of-equation}{%
\subparagraph{Explanation of equation}\label{explanation-of-equation}}

\begin{itemize}
\tightlist
\item
  The hidden state \(\mathbf{a}^{\langle t \rangle}\) is determined by
  the cell state \(\mathbf{c}^{\langle t \rangle}\) in combination with
  the output gate \(\mathbf{\Gamma}_{o}\).
\item
  The cell state state is passed through the \texttt{tanh} function to
  rescale values between -1 and 1.
\item
  The output gate acts like a ``mask'' that either preserves the values
  of \(\tanh(\mathbf{c}^{\langle t \rangle})\) or keeps those values
  from being included in the hidden state
  \(\mathbf{a}^{\langle t \rangle}\)
\end{itemize}

\hypertarget{variable-names-and-shapes-in-the-code}{%
\subparagraph{Variable names and shapes in the
code}\label{variable-names-and-shapes-in-the-code}}

\begin{itemize}
\tightlist
\item
  \texttt{a}: hidden state, including time steps. \(\mathbf{a}\) has
  shape \((n_{a}, m, T_{x})\)
\item
  \texttt{a\_prev}: hidden state from previous time step.
  \(\mathbf{a}^{\langle t-1 \rangle}\) has shape \((n_{a}, m)\)
\item
  \texttt{a\_next}: hidden state for next time step.
  \(\mathbf{a}^{\langle t \rangle}\) has shape \((n_{a}, m)\)
\end{itemize}

    \hypertarget{prediction-mathbfylangle-t-rangle_pred}{%
\paragraph{\texorpdfstring{Prediction
\(\mathbf{y}^{\langle t \rangle}_{pred}\)}{Prediction \textbackslash mathbf\{y\}\^{}\{\textbackslash langle t \textbackslash rangle\}\_\{pred\}}}\label{prediction-mathbfylangle-t-rangle_pred}}

\begin{itemize}
\tightlist
\item
  The prediction in this use case is a classification, so you'll use a
  softmax.
\end{itemize}

The equation is:
\[\mathbf{y}^{\langle t \rangle}_{pred} = \textrm{softmax}(\mathbf{W}_{y} \mathbf{a}^{\langle t \rangle} + \mathbf{b}_{y})\]

\hypertarget{variable-names-and-shapes-in-the-code}{%
\subparagraph{Variable names and shapes in the
code}\label{variable-names-and-shapes-in-the-code}}

\begin{itemize}
\tightlist
\item
  \texttt{y\_pred}: prediction, including all time steps.
  \(\mathbf{y}_{pred}\) has shape \((n_{y}, m, T_{x})\). Note that
  \((T_{y} = T_{x})\) for this example.
\item
  \texttt{yt\_pred}: prediction for the current time step \(t\).
  \(\mathbf{y}^{\langle t \rangle}_{pred}\) has shape \((n_{y}, m)\)
\end{itemize}

    \#\#\# 2.1 - LSTM Cell

\#\#\# Exercise 3 - lstm\_cell\_forward

Implement the LSTM cell described in Figure 4.

\textbf{Instructions}: 1. Concatenate the hidden state
\(a^{\langle t-1 \rangle}\) and input \(x^{\langle t \rangle}\) into a
single matrix:

\[concat = \begin{bmatrix} a^{\langle t-1 \rangle} \\ x^{\langle t \rangle} \end{bmatrix}\]

\begin{enumerate}
\def\labelenumi{\arabic{enumi}.}
\setcounter{enumi}{1}
\tightlist
\item
  Compute all formulas (1 through 6) for the gates, hidden state, and
  cell state.
\item
  Compute the prediction \(y^{\langle t \rangle}\).
\end{enumerate}

    \hypertarget{additional-hints}{%
\paragraph{Additional Hints}\label{additional-hints}}

\begin{itemize}
\tightlist
\item
  You can use
  \href{https://docs.scipy.org/doc/numpy/reference/generated/numpy.concatenate.html}{numpy.concatenate}.
  Check which value to use for the \texttt{axis} parameter.
\item
  The functions \texttt{sigmoid()} and \texttt{softmax} are imported
  from \texttt{rnn\_utils.py}.
\item
  Some docs for
  \href{https://docs.scipy.org/doc/numpy/reference/generated/numpy.tanh.html}{numpy.tanh}
\item
  Use
  \href{https://docs.scipy.org/doc/numpy/reference/generated/numpy.dot.html}{numpy.dot}
  for matrix multiplication.
\item
  Notice that the variable names \texttt{Wi}, \texttt{bi} refer to the
  weights and biases of the \textbf{update} gate. There are no variables
  named ``Wu'' or ``bu'' in this function.
\end{itemize}

    \begin{tcolorbox}[breakable, size=fbox, boxrule=1pt, pad at break*=1mm,colback=cellbackground, colframe=cellborder]
\prompt{In}{incolor}{29}{\boxspacing}
\begin{Verbatim}[commandchars=\\\{\}]
\PY{c+c1}{\PYZsh{} UNQ\PYZus{}C3 (UNIQUE CELL IDENTIFIER, DO NOT EDIT)}
\PY{c+c1}{\PYZsh{} GRADED FUNCTION: lstm\PYZus{}cell\PYZus{}forward}

\PY{k}{def} \PY{n+nf}{lstm\PYZus{}cell\PYZus{}forward}\PY{p}{(}\PY{n}{xt}\PY{p}{,} \PY{n}{a\PYZus{}prev}\PY{p}{,} \PY{n}{c\PYZus{}prev}\PY{p}{,} \PY{n}{parameters}\PY{p}{)}\PY{p}{:}
    \PY{l+s+sd}{\PYZdq{}\PYZdq{}\PYZdq{}}
\PY{l+s+sd}{    Implement a single forward step of the LSTM\PYZhy{}cell as described in Figure (4)}

\PY{l+s+sd}{    Arguments:}
\PY{l+s+sd}{    xt \PYZhy{}\PYZhy{} your input data at timestep \PYZdq{}t\PYZdq{}, numpy array of shape (n\PYZus{}x, m).}
\PY{l+s+sd}{    a\PYZus{}prev \PYZhy{}\PYZhy{} Hidden state at timestep \PYZdq{}t\PYZhy{}1\PYZdq{}, numpy array of shape (n\PYZus{}a, m)}
\PY{l+s+sd}{    c\PYZus{}prev \PYZhy{}\PYZhy{} Memory state at timestep \PYZdq{}t\PYZhy{}1\PYZdq{}, numpy array of shape (n\PYZus{}a, m)}
\PY{l+s+sd}{    parameters \PYZhy{}\PYZhy{} python dictionary containing:}
\PY{l+s+sd}{                        Wf \PYZhy{}\PYZhy{} Weight matrix of the forget gate, numpy array of shape (n\PYZus{}a, n\PYZus{}a + n\PYZus{}x)}
\PY{l+s+sd}{                        bf \PYZhy{}\PYZhy{} Bias of the forget gate, numpy array of shape (n\PYZus{}a, 1)}
\PY{l+s+sd}{                        Wi \PYZhy{}\PYZhy{} Weight matrix of the update gate, numpy array of shape (n\PYZus{}a, n\PYZus{}a + n\PYZus{}x)}
\PY{l+s+sd}{                        bi \PYZhy{}\PYZhy{} Bias of the update gate, numpy array of shape (n\PYZus{}a, 1)}
\PY{l+s+sd}{                        Wc \PYZhy{}\PYZhy{} Weight matrix of the first \PYZdq{}tanh\PYZdq{}, numpy array of shape (n\PYZus{}a, n\PYZus{}a + n\PYZus{}x)}
\PY{l+s+sd}{                        bc \PYZhy{}\PYZhy{}  Bias of the first \PYZdq{}tanh\PYZdq{}, numpy array of shape (n\PYZus{}a, 1)}
\PY{l+s+sd}{                        Wo \PYZhy{}\PYZhy{} Weight matrix of the output gate, numpy array of shape (n\PYZus{}a, n\PYZus{}a + n\PYZus{}x)}
\PY{l+s+sd}{                        bo \PYZhy{}\PYZhy{}  Bias of the output gate, numpy array of shape (n\PYZus{}a, 1)}
\PY{l+s+sd}{                        Wy \PYZhy{}\PYZhy{} Weight matrix relating the hidden\PYZhy{}state to the output, numpy array of shape (n\PYZus{}y, n\PYZus{}a)}
\PY{l+s+sd}{                        by \PYZhy{}\PYZhy{} Bias relating the hidden\PYZhy{}state to the output, numpy array of shape (n\PYZus{}y, 1)}
\PY{l+s+sd}{                        }
\PY{l+s+sd}{    Returns:}
\PY{l+s+sd}{    a\PYZus{}next \PYZhy{}\PYZhy{} next hidden state, of shape (n\PYZus{}a, m)}
\PY{l+s+sd}{    c\PYZus{}next \PYZhy{}\PYZhy{} next memory state, of shape (n\PYZus{}a, m)}
\PY{l+s+sd}{    yt\PYZus{}pred \PYZhy{}\PYZhy{} prediction at timestep \PYZdq{}t\PYZdq{}, numpy array of shape (n\PYZus{}y, m)}
\PY{l+s+sd}{    cache \PYZhy{}\PYZhy{} tuple of values needed for the backward pass, contains (a\PYZus{}next, c\PYZus{}next, a\PYZus{}prev, c\PYZus{}prev, xt, parameters)}
\PY{l+s+sd}{    }
\PY{l+s+sd}{    Note: ft/it/ot stand for the forget/update/output gates, cct stands for the candidate value (c tilde),}
\PY{l+s+sd}{          c stands for the cell state (memory)}
\PY{l+s+sd}{    \PYZdq{}\PYZdq{}\PYZdq{}}

    \PY{c+c1}{\PYZsh{} Retrieve parameters from \PYZdq{}parameters\PYZdq{}}
    \PY{n}{Wf} \PY{o}{=} \PY{n}{parameters}\PY{p}{[}\PY{l+s+s2}{\PYZdq{}}\PY{l+s+s2}{Wf}\PY{l+s+s2}{\PYZdq{}}\PY{p}{]} \PY{c+c1}{\PYZsh{} forget gate weight}
    \PY{n}{bf} \PY{o}{=} \PY{n}{parameters}\PY{p}{[}\PY{l+s+s2}{\PYZdq{}}\PY{l+s+s2}{bf}\PY{l+s+s2}{\PYZdq{}}\PY{p}{]}
    \PY{n}{Wi} \PY{o}{=} \PY{n}{parameters}\PY{p}{[}\PY{l+s+s2}{\PYZdq{}}\PY{l+s+s2}{Wi}\PY{l+s+s2}{\PYZdq{}}\PY{p}{]} \PY{c+c1}{\PYZsh{} update gate weight (notice the variable name)}
    \PY{n}{bi} \PY{o}{=} \PY{n}{parameters}\PY{p}{[}\PY{l+s+s2}{\PYZdq{}}\PY{l+s+s2}{bi}\PY{l+s+s2}{\PYZdq{}}\PY{p}{]} \PY{c+c1}{\PYZsh{} (notice the variable name)}
    \PY{n}{Wc} \PY{o}{=} \PY{n}{parameters}\PY{p}{[}\PY{l+s+s2}{\PYZdq{}}\PY{l+s+s2}{Wc}\PY{l+s+s2}{\PYZdq{}}\PY{p}{]} \PY{c+c1}{\PYZsh{} candidate value weight}
    \PY{n}{bc} \PY{o}{=} \PY{n}{parameters}\PY{p}{[}\PY{l+s+s2}{\PYZdq{}}\PY{l+s+s2}{bc}\PY{l+s+s2}{\PYZdq{}}\PY{p}{]}
    \PY{n}{Wo} \PY{o}{=} \PY{n}{parameters}\PY{p}{[}\PY{l+s+s2}{\PYZdq{}}\PY{l+s+s2}{Wo}\PY{l+s+s2}{\PYZdq{}}\PY{p}{]} \PY{c+c1}{\PYZsh{} output gate weight}
    \PY{n}{bo} \PY{o}{=} \PY{n}{parameters}\PY{p}{[}\PY{l+s+s2}{\PYZdq{}}\PY{l+s+s2}{bo}\PY{l+s+s2}{\PYZdq{}}\PY{p}{]}
    \PY{n}{Wy} \PY{o}{=} \PY{n}{parameters}\PY{p}{[}\PY{l+s+s2}{\PYZdq{}}\PY{l+s+s2}{Wy}\PY{l+s+s2}{\PYZdq{}}\PY{p}{]} \PY{c+c1}{\PYZsh{} prediction weight}
    \PY{n}{by} \PY{o}{=} \PY{n}{parameters}\PY{p}{[}\PY{l+s+s2}{\PYZdq{}}\PY{l+s+s2}{by}\PY{l+s+s2}{\PYZdq{}}\PY{p}{]}
    
    \PY{c+c1}{\PYZsh{} Retrieve dimensions from shapes of xt and Wy}
    \PY{n}{n\PYZus{}x}\PY{p}{,} \PY{n}{m} \PY{o}{=} \PY{n}{xt}\PY{o}{.}\PY{n}{shape}
    \PY{n}{n\PYZus{}y}\PY{p}{,} \PY{n}{n\PYZus{}a} \PY{o}{=} \PY{n}{Wy}\PY{o}{.}\PY{n}{shape}

    \PY{c+c1}{\PYZsh{}\PYZsh{}\PYZsh{} START CODE HERE \PYZsh{}\PYZsh{}\PYZsh{}}
    \PY{c+c1}{\PYZsh{} Concatenate a\PYZus{}prev and xt (≈1 line)}
    \PY{n}{concat} \PY{o}{=} \PY{n}{np}\PY{o}{.}\PY{n}{concatenate}\PY{p}{(}\PY{p}{(}\PY{n}{a\PYZus{}prev}\PY{p}{,} \PY{n}{xt}\PY{p}{)}\PY{p}{,} \PY{n}{axis}\PY{o}{=}\PY{l+m+mi}{0}\PY{p}{)}

    \PY{c+c1}{\PYZsh{} Compute values for ft, it, cct, c\PYZus{}next, ot, a\PYZus{}next using the formulas given figure (4) (≈6 lines)}
    \PY{n}{ft} \PY{o}{=} \PY{n}{sigmoid}\PY{p}{(}\PY{n}{np}\PY{o}{.}\PY{n}{dot}\PY{p}{(}\PY{n}{Wf}\PY{p}{,} \PY{n}{concat}\PY{p}{)} \PY{o}{+} \PY{n}{bf}\PY{p}{)}
    \PY{n}{it} \PY{o}{=} \PY{n}{sigmoid}\PY{p}{(}\PY{n}{np}\PY{o}{.}\PY{n}{dot}\PY{p}{(}\PY{n}{Wi}\PY{p}{,} \PY{n}{concat}\PY{p}{)} \PY{o}{+} \PY{n}{bi}\PY{p}{)}
    \PY{n}{cct} \PY{o}{=} \PY{n}{np}\PY{o}{.}\PY{n}{tanh}\PY{p}{(}\PY{n}{np}\PY{o}{.}\PY{n}{dot}\PY{p}{(}\PY{n}{Wc}\PY{p}{,} \PY{n}{concat}\PY{p}{)} \PY{o}{+} \PY{n}{bc}\PY{p}{)}
    \PY{n}{c\PYZus{}next} \PY{o}{=} \PY{n}{ft} \PY{o}{*} \PY{n}{c\PYZus{}prev} \PY{o}{+} \PY{n}{it} \PY{o}{*} \PY{n}{cct}
    \PY{n}{ot} \PY{o}{=} \PY{n}{sigmoid}\PY{p}{(}\PY{n}{np}\PY{o}{.}\PY{n}{matmul}\PY{p}{(}\PY{n}{Wo}\PY{p}{,} \PY{n}{concat}\PY{p}{)} \PY{o}{+} \PY{n}{bo}\PY{p}{)}
    \PY{n}{a\PYZus{}next} \PY{o}{=} \PY{n}{ot} \PY{o}{*} \PY{n}{np}\PY{o}{.}\PY{n}{tanh}\PY{p}{(}\PY{n}{c\PYZus{}next}\PY{p}{)}
    
    \PY{c+c1}{\PYZsh{} Compute prediction of the LSTM cell (≈1 line)}
    \PY{n}{yt\PYZus{}pred} \PY{o}{=} \PY{n}{softmax}\PY{p}{(}\PY{n}{np}\PY{o}{.}\PY{n}{matmul}\PY{p}{(}\PY{n}{Wy}\PY{p}{,} \PY{n}{a\PYZus{}next}\PY{p}{)} \PY{o}{+} \PY{n}{by}\PY{p}{)}
    \PY{c+c1}{\PYZsh{}\PYZsh{}\PYZsh{} END CODE HERE \PYZsh{}\PYZsh{}\PYZsh{}}

    \PY{c+c1}{\PYZsh{} store values needed for backward propagation in cache}
    \PY{n}{cache} \PY{o}{=} \PY{p}{(}\PY{n}{a\PYZus{}next}\PY{p}{,} \PY{n}{c\PYZus{}next}\PY{p}{,} \PY{n}{a\PYZus{}prev}\PY{p}{,} \PY{n}{c\PYZus{}prev}\PY{p}{,} \PY{n}{ft}\PY{p}{,} \PY{n}{it}\PY{p}{,} \PY{n}{cct}\PY{p}{,} \PY{n}{ot}\PY{p}{,} \PY{n}{xt}\PY{p}{,} \PY{n}{parameters}\PY{p}{)}

    \PY{k}{return} \PY{n}{a\PYZus{}next}\PY{p}{,} \PY{n}{c\PYZus{}next}\PY{p}{,} \PY{n}{yt\PYZus{}pred}\PY{p}{,} \PY{n}{cache}
\end{Verbatim}
\end{tcolorbox}

    \begin{tcolorbox}[breakable, size=fbox, boxrule=1pt, pad at break*=1mm,colback=cellbackground, colframe=cellborder]
\prompt{In}{incolor}{30}{\boxspacing}
\begin{Verbatim}[commandchars=\\\{\}]
\PY{n}{np}\PY{o}{.}\PY{n}{random}\PY{o}{.}\PY{n}{seed}\PY{p}{(}\PY{l+m+mi}{1}\PY{p}{)}
\PY{n}{xt\PYZus{}tmp} \PY{o}{=} \PY{n}{np}\PY{o}{.}\PY{n}{random}\PY{o}{.}\PY{n}{randn}\PY{p}{(}\PY{l+m+mi}{3}\PY{p}{,} \PY{l+m+mi}{10}\PY{p}{)}
\PY{n}{a\PYZus{}prev\PYZus{}tmp} \PY{o}{=} \PY{n}{np}\PY{o}{.}\PY{n}{random}\PY{o}{.}\PY{n}{randn}\PY{p}{(}\PY{l+m+mi}{5}\PY{p}{,} \PY{l+m+mi}{10}\PY{p}{)}
\PY{n}{c\PYZus{}prev\PYZus{}tmp} \PY{o}{=} \PY{n}{np}\PY{o}{.}\PY{n}{random}\PY{o}{.}\PY{n}{randn}\PY{p}{(}\PY{l+m+mi}{5}\PY{p}{,} \PY{l+m+mi}{10}\PY{p}{)}
\PY{n}{parameters\PYZus{}tmp} \PY{o}{=} \PY{p}{\PYZob{}}\PY{p}{\PYZcb{}}
\PY{n}{parameters\PYZus{}tmp}\PY{p}{[}\PY{l+s+s1}{\PYZsq{}}\PY{l+s+s1}{Wf}\PY{l+s+s1}{\PYZsq{}}\PY{p}{]} \PY{o}{=} \PY{n}{np}\PY{o}{.}\PY{n}{random}\PY{o}{.}\PY{n}{randn}\PY{p}{(}\PY{l+m+mi}{5}\PY{p}{,} \PY{l+m+mi}{5} \PY{o}{+} \PY{l+m+mi}{3}\PY{p}{)}
\PY{n}{parameters\PYZus{}tmp}\PY{p}{[}\PY{l+s+s1}{\PYZsq{}}\PY{l+s+s1}{bf}\PY{l+s+s1}{\PYZsq{}}\PY{p}{]} \PY{o}{=} \PY{n}{np}\PY{o}{.}\PY{n}{random}\PY{o}{.}\PY{n}{randn}\PY{p}{(}\PY{l+m+mi}{5}\PY{p}{,} \PY{l+m+mi}{1}\PY{p}{)}
\PY{n}{parameters\PYZus{}tmp}\PY{p}{[}\PY{l+s+s1}{\PYZsq{}}\PY{l+s+s1}{Wi}\PY{l+s+s1}{\PYZsq{}}\PY{p}{]} \PY{o}{=} \PY{n}{np}\PY{o}{.}\PY{n}{random}\PY{o}{.}\PY{n}{randn}\PY{p}{(}\PY{l+m+mi}{5}\PY{p}{,} \PY{l+m+mi}{5} \PY{o}{+} \PY{l+m+mi}{3}\PY{p}{)}
\PY{n}{parameters\PYZus{}tmp}\PY{p}{[}\PY{l+s+s1}{\PYZsq{}}\PY{l+s+s1}{bi}\PY{l+s+s1}{\PYZsq{}}\PY{p}{]} \PY{o}{=} \PY{n}{np}\PY{o}{.}\PY{n}{random}\PY{o}{.}\PY{n}{randn}\PY{p}{(}\PY{l+m+mi}{5}\PY{p}{,} \PY{l+m+mi}{1}\PY{p}{)}
\PY{n}{parameters\PYZus{}tmp}\PY{p}{[}\PY{l+s+s1}{\PYZsq{}}\PY{l+s+s1}{Wo}\PY{l+s+s1}{\PYZsq{}}\PY{p}{]} \PY{o}{=} \PY{n}{np}\PY{o}{.}\PY{n}{random}\PY{o}{.}\PY{n}{randn}\PY{p}{(}\PY{l+m+mi}{5}\PY{p}{,} \PY{l+m+mi}{5} \PY{o}{+} \PY{l+m+mi}{3}\PY{p}{)}
\PY{n}{parameters\PYZus{}tmp}\PY{p}{[}\PY{l+s+s1}{\PYZsq{}}\PY{l+s+s1}{bo}\PY{l+s+s1}{\PYZsq{}}\PY{p}{]} \PY{o}{=} \PY{n}{np}\PY{o}{.}\PY{n}{random}\PY{o}{.}\PY{n}{randn}\PY{p}{(}\PY{l+m+mi}{5}\PY{p}{,} \PY{l+m+mi}{1}\PY{p}{)}
\PY{n}{parameters\PYZus{}tmp}\PY{p}{[}\PY{l+s+s1}{\PYZsq{}}\PY{l+s+s1}{Wc}\PY{l+s+s1}{\PYZsq{}}\PY{p}{]} \PY{o}{=} \PY{n}{np}\PY{o}{.}\PY{n}{random}\PY{o}{.}\PY{n}{randn}\PY{p}{(}\PY{l+m+mi}{5}\PY{p}{,} \PY{l+m+mi}{5} \PY{o}{+} \PY{l+m+mi}{3}\PY{p}{)}
\PY{n}{parameters\PYZus{}tmp}\PY{p}{[}\PY{l+s+s1}{\PYZsq{}}\PY{l+s+s1}{bc}\PY{l+s+s1}{\PYZsq{}}\PY{p}{]} \PY{o}{=} \PY{n}{np}\PY{o}{.}\PY{n}{random}\PY{o}{.}\PY{n}{randn}\PY{p}{(}\PY{l+m+mi}{5}\PY{p}{,} \PY{l+m+mi}{1}\PY{p}{)}
\PY{n}{parameters\PYZus{}tmp}\PY{p}{[}\PY{l+s+s1}{\PYZsq{}}\PY{l+s+s1}{Wy}\PY{l+s+s1}{\PYZsq{}}\PY{p}{]} \PY{o}{=} \PY{n}{np}\PY{o}{.}\PY{n}{random}\PY{o}{.}\PY{n}{randn}\PY{p}{(}\PY{l+m+mi}{2}\PY{p}{,} \PY{l+m+mi}{5}\PY{p}{)}
\PY{n}{parameters\PYZus{}tmp}\PY{p}{[}\PY{l+s+s1}{\PYZsq{}}\PY{l+s+s1}{by}\PY{l+s+s1}{\PYZsq{}}\PY{p}{]} \PY{o}{=} \PY{n}{np}\PY{o}{.}\PY{n}{random}\PY{o}{.}\PY{n}{randn}\PY{p}{(}\PY{l+m+mi}{2}\PY{p}{,} \PY{l+m+mi}{1}\PY{p}{)}

\PY{n}{a\PYZus{}next\PYZus{}tmp}\PY{p}{,} \PY{n}{c\PYZus{}next\PYZus{}tmp}\PY{p}{,} \PY{n}{yt\PYZus{}tmp}\PY{p}{,} \PY{n}{cache\PYZus{}tmp} \PY{o}{=} \PY{n}{lstm\PYZus{}cell\PYZus{}forward}\PY{p}{(}\PY{n}{xt\PYZus{}tmp}\PY{p}{,} \PY{n}{a\PYZus{}prev\PYZus{}tmp}\PY{p}{,} \PY{n}{c\PYZus{}prev\PYZus{}tmp}\PY{p}{,} \PY{n}{parameters\PYZus{}tmp}\PY{p}{)}

\PY{n+nb}{print}\PY{p}{(}\PY{l+s+s2}{\PYZdq{}}\PY{l+s+s2}{a\PYZus{}next[4] = }\PY{l+s+se}{\PYZbs{}n}\PY{l+s+s2}{\PYZdq{}}\PY{p}{,} \PY{n}{a\PYZus{}next\PYZus{}tmp}\PY{p}{[}\PY{l+m+mi}{4}\PY{p}{]}\PY{p}{)}
\PY{n+nb}{print}\PY{p}{(}\PY{l+s+s2}{\PYZdq{}}\PY{l+s+s2}{a\PYZus{}next.shape = }\PY{l+s+s2}{\PYZdq{}}\PY{p}{,} \PY{n}{a\PYZus{}next\PYZus{}tmp}\PY{o}{.}\PY{n}{shape}\PY{p}{)}
\PY{n+nb}{print}\PY{p}{(}\PY{l+s+s2}{\PYZdq{}}\PY{l+s+s2}{c\PYZus{}next[2] = }\PY{l+s+se}{\PYZbs{}n}\PY{l+s+s2}{\PYZdq{}}\PY{p}{,} \PY{n}{c\PYZus{}next\PYZus{}tmp}\PY{p}{[}\PY{l+m+mi}{2}\PY{p}{]}\PY{p}{)}
\PY{n+nb}{print}\PY{p}{(}\PY{l+s+s2}{\PYZdq{}}\PY{l+s+s2}{c\PYZus{}next.shape = }\PY{l+s+s2}{\PYZdq{}}\PY{p}{,} \PY{n}{c\PYZus{}next\PYZus{}tmp}\PY{o}{.}\PY{n}{shape}\PY{p}{)}
\PY{n+nb}{print}\PY{p}{(}\PY{l+s+s2}{\PYZdq{}}\PY{l+s+s2}{yt[1] =}\PY{l+s+s2}{\PYZdq{}}\PY{p}{,} \PY{n}{yt\PYZus{}tmp}\PY{p}{[}\PY{l+m+mi}{1}\PY{p}{]}\PY{p}{)}
\PY{n+nb}{print}\PY{p}{(}\PY{l+s+s2}{\PYZdq{}}\PY{l+s+s2}{yt.shape = }\PY{l+s+s2}{\PYZdq{}}\PY{p}{,} \PY{n}{yt\PYZus{}tmp}\PY{o}{.}\PY{n}{shape}\PY{p}{)}
\PY{n+nb}{print}\PY{p}{(}\PY{l+s+s2}{\PYZdq{}}\PY{l+s+s2}{cache[1][3] =}\PY{l+s+se}{\PYZbs{}n}\PY{l+s+s2}{\PYZdq{}}\PY{p}{,} \PY{n}{cache\PYZus{}tmp}\PY{p}{[}\PY{l+m+mi}{1}\PY{p}{]}\PY{p}{[}\PY{l+m+mi}{3}\PY{p}{]}\PY{p}{)}
\PY{n+nb}{print}\PY{p}{(}\PY{l+s+s2}{\PYZdq{}}\PY{l+s+s2}{len(cache) = }\PY{l+s+s2}{\PYZdq{}}\PY{p}{,} \PY{n+nb}{len}\PY{p}{(}\PY{n}{cache\PYZus{}tmp}\PY{p}{)}\PY{p}{)}

\PY{c+c1}{\PYZsh{} UNIT TEST}
\PY{n}{lstm\PYZus{}cell\PYZus{}forward\PYZus{}test}\PY{p}{(}\PY{n}{lstm\PYZus{}cell\PYZus{}forward}\PY{p}{)}
\end{Verbatim}
\end{tcolorbox}

    \begin{Verbatim}[commandchars=\\\{\}]
a\_next[4] =
 [-0.66408471  0.0036921   0.02088357  0.22834167 -0.85575339  0.00138482
  0.76566531  0.34631421 -0.00215674  0.43827275]
a\_next.shape =  (5, 10)
c\_next[2] =
 [ 0.63267805  1.00570849  0.35504474  0.20690913 -1.64566718  0.11832942
  0.76449811 -0.0981561  -0.74348425 -0.26810932]
c\_next.shape =  (5, 10)
yt[1] = [0.79913913 0.15986619 0.22412122 0.15606108 0.97057211 0.31146381
 0.00943007 0.12666353 0.39380172 0.07828381]
yt.shape =  (2, 10)
cache[1][3] =
 [-0.16263996  1.03729328  0.72938082 -0.54101719  0.02752074 -0.30821874
  0.07651101 -1.03752894  1.41219977 -0.37647422]
len(cache) =  10
\textcolor{ansi-green-intense}{All tests passed}
    \end{Verbatim}

    \textbf{Expected Output}:

\begin{Shaded}
\begin{Highlighting}[]
\NormalTok{a\_next[}\DecValTok{4}\NormalTok{] }\OperatorTok{=} 
\NormalTok{ [}\OperatorTok{{-}}\FloatTok{0.66408471}  \FloatTok{0.0036921}   \FloatTok{0.02088357}  \FloatTok{0.22834167} \OperatorTok{{-}}\FloatTok{0.85575339}  \FloatTok{0.00138482}
  \FloatTok{0.76566531}  \FloatTok{0.34631421} \OperatorTok{{-}}\FloatTok{0.00215674}  \FloatTok{0.43827275}\NormalTok{]}
\NormalTok{a\_next.shape }\OperatorTok{=}\NormalTok{  (}\DecValTok{5}\NormalTok{, }\DecValTok{10}\NormalTok{)}
\NormalTok{c\_next[}\DecValTok{2}\NormalTok{] }\OperatorTok{=} 
\NormalTok{ [ }\FloatTok{0.63267805}  \FloatTok{1.00570849}  \FloatTok{0.35504474}  \FloatTok{0.20690913} \OperatorTok{{-}}\FloatTok{1.64566718}  \FloatTok{0.11832942}
  \FloatTok{0.76449811} \OperatorTok{{-}}\FloatTok{0.0981561}  \OperatorTok{{-}}\FloatTok{0.74348425} \OperatorTok{{-}}\FloatTok{0.26810932}\NormalTok{]}
\NormalTok{c\_next.shape }\OperatorTok{=}\NormalTok{  (}\DecValTok{5}\NormalTok{, }\DecValTok{10}\NormalTok{)}
\NormalTok{yt[}\DecValTok{1}\NormalTok{] }\OperatorTok{=}\NormalTok{ [ }\FloatTok{0.79913913}  \FloatTok{0.15986619}  \FloatTok{0.22412122}  \FloatTok{0.15606108}  \FloatTok{0.97057211}  \FloatTok{0.31146381}
  \FloatTok{0.00943007}  \FloatTok{0.12666353}  \FloatTok{0.39380172}  \FloatTok{0.07828381}\NormalTok{]}
\NormalTok{yt.shape }\OperatorTok{=}\NormalTok{  (}\DecValTok{2}\NormalTok{, }\DecValTok{10}\NormalTok{)}
\NormalTok{cache[}\DecValTok{1}\NormalTok{][}\DecValTok{3}\NormalTok{] }\OperatorTok{=}
\NormalTok{ [}\OperatorTok{{-}}\FloatTok{0.16263996}  \FloatTok{1.03729328}  \FloatTok{0.72938082} \OperatorTok{{-}}\FloatTok{0.54101719}  \FloatTok{0.02752074} \OperatorTok{{-}}\FloatTok{0.30821874}
  \FloatTok{0.07651101} \OperatorTok{{-}}\FloatTok{1.03752894}  \FloatTok{1.41219977} \OperatorTok{{-}}\FloatTok{0.37647422}\NormalTok{]}
\BuiltInTok{len}\NormalTok{(cache) }\OperatorTok{=}  \DecValTok{10}
\end{Highlighting}
\end{Shaded}

    \#\#\# 2.2 - Forward Pass for LSTM

Now that you have implemented one step of an LSTM, you can iterate this
over it using a for loop to process a sequence of \(T_x\) inputs.

Figure 5: LSTM over multiple time steps.

~\\
\#\#\# Exercise 4 - lstm\_forward

Implement \texttt{lstm\_forward()} to run an LSTM over \(T_x\) time
steps.

\textbf{Instructions} * Get the dimensions \(n_x, n_a, n_y, m, T_x\)
from the shape of the variables: \texttt{x} and \texttt{parameters} *
Initialize the 3D tensors \(a\), \(c\) and \(y\) - \(a\): hidden state,
shape \((n_{a}, m, T_{x})\) - \(c\): cell state, shape
\((n_{a}, m, T_{x})\) - \(y\): prediction, shape \((n_{y}, m, T_{x})\)
(Note that \(T_{y} = T_{x}\) in this example) - \textbf{Note} Setting
one variable equal to the other is a ``copy by reference''. In other
words, don't do
\texttt{c\ =\ a\textquotesingle{},\ otherwise\ both\ these\ variables\ point\ to\ the\ same\ underlying\ variable.\ *\ Initialize\ the\ 2D\ tensor\ \$a\^{}\{\textbackslash{}langle\ t\ \textbackslash{}rangle\}\$\ \ \ \ \ \ -\ \$a\^{}\{\textbackslash{}langle\ t\ \textbackslash{}rangle\}\$\ stores\ the\ hidden\ state\ for\ time\ step\ \$t\$.\ \ The\ variable\ name\ is}a\_next\texttt{.\ \ \ \ \ -\ \$a\^{}\{\textbackslash{}langle\ 0\ \textbackslash{}rangle\}\$,\ the\ initial\ hidden\ state\ at\ time\ step\ 0,\ is\ passed\ in\ when\ calling\ the\ function.\ The\ variable\ name\ is}a0\texttt{.\ \ \ \ \ -\ \$a\^{}\{\textbackslash{}langle\ t\ \textbackslash{}rangle\}\$\ and\ \$a\^{}\{\textbackslash{}langle\ 0\ \textbackslash{}rangle\}\$\ represent\ a\ single\ time\ step,\ so\ they\ both\ have\ the\ shape\ \ \$(n\_\{a\},\ m)\$\ \ \ \ \ \ -\ Initialize\ \$a\^{}\{\textbackslash{}langle\ t\ \textbackslash{}rangle\}\$\ by\ setting\ it\ to\ the\ initial\ hidden\ state\ (\$a\^{}\{\textbackslash{}langle\ 0\ \textbackslash{}rangle\}\$)\ that\ is\ passed\ into\ the\ function.\ *\ Initialize\ \$c\^{}\{\textbackslash{}langle\ t\ \textbackslash{}rangle\}\$\ with\ zeros.\ \ \ \ \ \ -\ The\ variable\ name\ is}c\_next\texttt{-\ \$c\^{}\{\textbackslash{}langle\ t\ \textbackslash{}rangle\}\$\ represents\ a\ single\ time\ step,\ so\ its\ shape\ is\ \$(n\_\{a\},\ m)\$\ \ \ \ \ -\ **Note**:\ create}c\_next\texttt{as\ its\ own\ variable\ with\ its\ own\ location\ in\ memory.\ \ Do\ not\ initialize\ it\ as\ a\ slice\ of\ the\ 3D\ tensor\ \$c\$.\ \ In\ other\ words,\ **don\textquotesingle{}t**\ do}c\_next
=
c{[}:,:,0{]}\texttt{.\ *\ For\ each\ time\ step,\ do\ the\ following:\ \ \ \ \ -\ From\ the\ 3D\ tensor\ \$x\$,\ get\ a\ 2D\ slice\ \$x\^{}\{\textbackslash{}langle\ t\ \textbackslash{}rangle\}\$\ at\ time\ step\ \$t\$\ \ \ \ \ -\ Call\ the}lstm\_cell\_forward`
function that you defined previously, to get the hidden state, cell
state, prediction, and cache - Store the hidden state, cell state and
prediction (the 2D tensors) inside the 3D tensors - Append the cache to
the list of caches

    \begin{tcolorbox}[breakable, size=fbox, boxrule=1pt, pad at break*=1mm,colback=cellbackground, colframe=cellborder]
\prompt{In}{incolor}{39}{\boxspacing}
\begin{Verbatim}[commandchars=\\\{\}]
\PY{c+c1}{\PYZsh{} UNQ\PYZus{}C4 (UNIQUE CELL IDENTIFIER, DO NOT EDIT)}
\PY{c+c1}{\PYZsh{} GRADED FUNCTION: lstm\PYZus{}forward}

\PY{k}{def} \PY{n+nf}{lstm\PYZus{}forward}\PY{p}{(}\PY{n}{x}\PY{p}{,} \PY{n}{a0}\PY{p}{,} \PY{n}{parameters}\PY{p}{)}\PY{p}{:}
    \PY{l+s+sd}{\PYZdq{}\PYZdq{}\PYZdq{}}
\PY{l+s+sd}{    Implement the forward propagation of the recurrent neural network using an LSTM\PYZhy{}cell described in Figure (4).}

\PY{l+s+sd}{    Arguments:}
\PY{l+s+sd}{    x \PYZhy{}\PYZhy{} Input data for every time\PYZhy{}step, of shape (n\PYZus{}x, m, T\PYZus{}x).}
\PY{l+s+sd}{    a0 \PYZhy{}\PYZhy{} Initial hidden state, of shape (n\PYZus{}a, m)}
\PY{l+s+sd}{    parameters \PYZhy{}\PYZhy{} python dictionary containing:}
\PY{l+s+sd}{                        Wf \PYZhy{}\PYZhy{} Weight matrix of the forget gate, numpy array of shape (n\PYZus{}a, n\PYZus{}a + n\PYZus{}x)}
\PY{l+s+sd}{                        bf \PYZhy{}\PYZhy{} Bias of the forget gate, numpy array of shape (n\PYZus{}a, 1)}
\PY{l+s+sd}{                        Wi \PYZhy{}\PYZhy{} Weight matrix of the update gate, numpy array of shape (n\PYZus{}a, n\PYZus{}a + n\PYZus{}x)}
\PY{l+s+sd}{                        bi \PYZhy{}\PYZhy{} Bias of the update gate, numpy array of shape (n\PYZus{}a, 1)}
\PY{l+s+sd}{                        Wc \PYZhy{}\PYZhy{} Weight matrix of the first \PYZdq{}tanh\PYZdq{}, numpy array of shape (n\PYZus{}a, n\PYZus{}a + n\PYZus{}x)}
\PY{l+s+sd}{                        bc \PYZhy{}\PYZhy{} Bias of the first \PYZdq{}tanh\PYZdq{}, numpy array of shape (n\PYZus{}a, 1)}
\PY{l+s+sd}{                        Wo \PYZhy{}\PYZhy{} Weight matrix of the output gate, numpy array of shape (n\PYZus{}a, n\PYZus{}a + n\PYZus{}x)}
\PY{l+s+sd}{                        bo \PYZhy{}\PYZhy{} Bias of the output gate, numpy array of shape (n\PYZus{}a, 1)}
\PY{l+s+sd}{                        Wy \PYZhy{}\PYZhy{} Weight matrix relating the hidden\PYZhy{}state to the output, numpy array of shape (n\PYZus{}y, n\PYZus{}a)}
\PY{l+s+sd}{                        by \PYZhy{}\PYZhy{} Bias relating the hidden\PYZhy{}state to the output, numpy array of shape (n\PYZus{}y, 1)}
\PY{l+s+sd}{                        }
\PY{l+s+sd}{    Returns:}
\PY{l+s+sd}{    a \PYZhy{}\PYZhy{} Hidden states for every time\PYZhy{}step, numpy array of shape (n\PYZus{}a, m, T\PYZus{}x)}
\PY{l+s+sd}{    y \PYZhy{}\PYZhy{} Predictions for every time\PYZhy{}step, numpy array of shape (n\PYZus{}y, m, T\PYZus{}x)}
\PY{l+s+sd}{    c \PYZhy{}\PYZhy{} The value of the cell state, numpy array of shape (n\PYZus{}a, m, T\PYZus{}x)}
\PY{l+s+sd}{    caches \PYZhy{}\PYZhy{} tuple of values needed for the backward pass, contains (list of all the caches, x)}
\PY{l+s+sd}{    \PYZdq{}\PYZdq{}\PYZdq{}}

    \PY{c+c1}{\PYZsh{} Initialize \PYZdq{}caches\PYZdq{}, which will track the list of all the caches}
    \PY{n}{caches} \PY{o}{=} \PY{p}{[}\PY{p}{]}
    
    \PY{c+c1}{\PYZsh{}\PYZsh{}\PYZsh{} START CODE HERE \PYZsh{}\PYZsh{}\PYZsh{}}
    \PY{n}{Wy} \PY{o}{=} \PY{n}{parameters}\PY{p}{[}\PY{l+s+s1}{\PYZsq{}}\PY{l+s+s1}{Wy}\PY{l+s+s1}{\PYZsq{}}\PY{p}{]} \PY{c+c1}{\PYZsh{} saving parameters[\PYZsq{}Wy\PYZsq{}] in a local variable in case students use Wy instead of parameters[\PYZsq{}Wy\PYZsq{}]}
    \PY{c+c1}{\PYZsh{} Retrieve dimensions from shapes of x and parameters[\PYZsq{}Wy\PYZsq{}] (≈2 lines)}
    \PY{n}{n\PYZus{}x}\PY{p}{,} \PY{n}{m}\PY{p}{,} \PY{n}{T\PYZus{}x} \PY{o}{=} \PY{n}{x}\PY{o}{.}\PY{n}{shape}
    \PY{n}{n\PYZus{}y}\PY{p}{,} \PY{n}{n\PYZus{}a} \PY{o}{=} \PY{n}{Wy}\PY{o}{.}\PY{n}{shape}
    
    \PY{c+c1}{\PYZsh{} initialize \PYZdq{}a\PYZdq{}, \PYZdq{}c\PYZdq{} and \PYZdq{}y\PYZdq{} with zeros (≈3 lines)}
    \PY{n}{a} \PY{o}{=} \PY{n}{np}\PY{o}{.}\PY{n}{zeros}\PY{p}{(}\PY{p}{(}\PY{n}{n\PYZus{}a}\PY{p}{,} \PY{n}{m}\PY{p}{,} \PY{n}{T\PYZus{}x}\PY{p}{)}\PY{p}{)}
    \PY{n}{c} \PY{o}{=} \PY{n}{np}\PY{o}{.}\PY{n}{zeros}\PY{p}{(}\PY{p}{(}\PY{n}{n\PYZus{}a}\PY{p}{,} \PY{n}{m}\PY{p}{,} \PY{n}{T\PYZus{}x}\PY{p}{)}\PY{p}{)}
    \PY{n}{y} \PY{o}{=} \PY{n}{np}\PY{o}{.}\PY{n}{zeros}\PY{p}{(}\PY{p}{(}\PY{n}{n\PYZus{}y}\PY{p}{,} \PY{n}{m}\PY{p}{,} \PY{n}{T\PYZus{}x}\PY{p}{)}\PY{p}{)}
    
    \PY{c+c1}{\PYZsh{} Initialize a\PYZus{}next and c\PYZus{}next (≈2 lines)}
    \PY{n}{a\PYZus{}next} \PY{o}{=} \PY{n}{a0}
    \PY{n}{c\PYZus{}next} \PY{o}{=} \PY{n}{np}\PY{o}{.}\PY{n}{zeros}\PY{p}{(}\PY{p}{(}\PY{n}{n\PYZus{}a}\PY{p}{,} \PY{n}{m}\PY{p}{)}\PY{p}{)}
    
    \PY{c+c1}{\PYZsh{} loop over all time\PYZhy{}steps}
    \PY{k}{for} \PY{n}{t} \PY{o+ow}{in} \PY{n+nb}{range}\PY{p}{(}\PY{n}{T\PYZus{}x}\PY{p}{)}\PY{p}{:}
        \PY{c+c1}{\PYZsh{} Get the 2D slice \PYZsq{}xt\PYZsq{} from the 3D input \PYZsq{}x\PYZsq{} at time step \PYZsq{}t\PYZsq{}}
        \PY{n}{xt} \PY{o}{=} \PY{n}{x}\PY{p}{[}\PY{p}{:}\PY{p}{,}\PY{p}{:}\PY{p}{,}\PY{n}{t}\PY{p}{]}
        \PY{c+c1}{\PYZsh{} Update next hidden state, next memory state, compute the prediction, get the cache (≈1 line)}
        \PY{n}{a\PYZus{}next}\PY{p}{,} \PY{n}{c\PYZus{}next}\PY{p}{,} \PY{n}{yt}\PY{p}{,} \PY{n}{cache} \PY{o}{=} \PY{n}{lstm\PYZus{}cell\PYZus{}forward}\PY{p}{(}\PY{n}{xt}\PY{p}{,} \PY{n}{a\PYZus{}next}\PY{p}{,} \PY{n}{c\PYZus{}next}\PY{p}{,} \PY{n}{parameters}\PY{p}{)}
        \PY{c+c1}{\PYZsh{} Save the value of the new \PYZdq{}next\PYZdq{} hidden state in a (≈1 line)}
        \PY{n}{a}\PY{p}{[}\PY{p}{:}\PY{p}{,}\PY{p}{:}\PY{p}{,}\PY{n}{t}\PY{p}{]} \PY{o}{=} \PY{n}{a\PYZus{}next}
        \PY{c+c1}{\PYZsh{} Save the value of the next cell state (≈1 line)}
        \PY{n}{c}\PY{p}{[}\PY{p}{:}\PY{p}{,}\PY{p}{:}\PY{p}{,}\PY{n}{t}\PY{p}{]}  \PY{o}{=} \PY{n}{c\PYZus{}next}
        \PY{c+c1}{\PYZsh{} Save the value of the prediction in y (≈1 line)}
        \PY{n}{y}\PY{p}{[}\PY{p}{:}\PY{p}{,}\PY{p}{:}\PY{p}{,}\PY{n}{t}\PY{p}{]} \PY{o}{=} \PY{n}{yt}
        \PY{c+c1}{\PYZsh{} Append the cache into caches (≈1 line)}
        \PY{n}{caches}\PY{o}{.}\PY{n}{append}\PY{p}{(}\PY{n}{cache}\PY{p}{)}
        
    \PY{c+c1}{\PYZsh{}\PYZsh{}\PYZsh{} END CODE HERE \PYZsh{}\PYZsh{}\PYZsh{}}
    
    \PY{c+c1}{\PYZsh{} store values needed for backward propagation in cache}
    \PY{n}{caches} \PY{o}{=} \PY{p}{(}\PY{n}{caches}\PY{p}{,} \PY{n}{x}\PY{p}{)}

    \PY{k}{return} \PY{n}{a}\PY{p}{,} \PY{n}{y}\PY{p}{,} \PY{n}{c}\PY{p}{,} \PY{n}{caches}
\end{Verbatim}
\end{tcolorbox}

    \begin{tcolorbox}[breakable, size=fbox, boxrule=1pt, pad at break*=1mm,colback=cellbackground, colframe=cellborder]
\prompt{In}{incolor}{40}{\boxspacing}
\begin{Verbatim}[commandchars=\\\{\}]
\PY{n}{np}\PY{o}{.}\PY{n}{random}\PY{o}{.}\PY{n}{seed}\PY{p}{(}\PY{l+m+mi}{1}\PY{p}{)}
\PY{n}{x\PYZus{}tmp} \PY{o}{=} \PY{n}{np}\PY{o}{.}\PY{n}{random}\PY{o}{.}\PY{n}{randn}\PY{p}{(}\PY{l+m+mi}{3}\PY{p}{,} \PY{l+m+mi}{10}\PY{p}{,} \PY{l+m+mi}{7}\PY{p}{)}
\PY{n}{a0\PYZus{}tmp} \PY{o}{=} \PY{n}{np}\PY{o}{.}\PY{n}{random}\PY{o}{.}\PY{n}{randn}\PY{p}{(}\PY{l+m+mi}{5}\PY{p}{,} \PY{l+m+mi}{10}\PY{p}{)}
\PY{n}{parameters\PYZus{}tmp} \PY{o}{=} \PY{p}{\PYZob{}}\PY{p}{\PYZcb{}}
\PY{n}{parameters\PYZus{}tmp}\PY{p}{[}\PY{l+s+s1}{\PYZsq{}}\PY{l+s+s1}{Wf}\PY{l+s+s1}{\PYZsq{}}\PY{p}{]} \PY{o}{=} \PY{n}{np}\PY{o}{.}\PY{n}{random}\PY{o}{.}\PY{n}{randn}\PY{p}{(}\PY{l+m+mi}{5}\PY{p}{,} \PY{l+m+mi}{5} \PY{o}{+} \PY{l+m+mi}{3}\PY{p}{)}
\PY{n}{parameters\PYZus{}tmp}\PY{p}{[}\PY{l+s+s1}{\PYZsq{}}\PY{l+s+s1}{bf}\PY{l+s+s1}{\PYZsq{}}\PY{p}{]} \PY{o}{=} \PY{n}{np}\PY{o}{.}\PY{n}{random}\PY{o}{.}\PY{n}{randn}\PY{p}{(}\PY{l+m+mi}{5}\PY{p}{,} \PY{l+m+mi}{1}\PY{p}{)}
\PY{n}{parameters\PYZus{}tmp}\PY{p}{[}\PY{l+s+s1}{\PYZsq{}}\PY{l+s+s1}{Wi}\PY{l+s+s1}{\PYZsq{}}\PY{p}{]} \PY{o}{=} \PY{n}{np}\PY{o}{.}\PY{n}{random}\PY{o}{.}\PY{n}{randn}\PY{p}{(}\PY{l+m+mi}{5}\PY{p}{,} \PY{l+m+mi}{5} \PY{o}{+} \PY{l+m+mi}{3}\PY{p}{)}
\PY{n}{parameters\PYZus{}tmp}\PY{p}{[}\PY{l+s+s1}{\PYZsq{}}\PY{l+s+s1}{bi}\PY{l+s+s1}{\PYZsq{}}\PY{p}{]}\PY{o}{=} \PY{n}{np}\PY{o}{.}\PY{n}{random}\PY{o}{.}\PY{n}{randn}\PY{p}{(}\PY{l+m+mi}{5}\PY{p}{,} \PY{l+m+mi}{1}\PY{p}{)}
\PY{n}{parameters\PYZus{}tmp}\PY{p}{[}\PY{l+s+s1}{\PYZsq{}}\PY{l+s+s1}{Wo}\PY{l+s+s1}{\PYZsq{}}\PY{p}{]} \PY{o}{=} \PY{n}{np}\PY{o}{.}\PY{n}{random}\PY{o}{.}\PY{n}{randn}\PY{p}{(}\PY{l+m+mi}{5}\PY{p}{,} \PY{l+m+mi}{5} \PY{o}{+} \PY{l+m+mi}{3}\PY{p}{)}
\PY{n}{parameters\PYZus{}tmp}\PY{p}{[}\PY{l+s+s1}{\PYZsq{}}\PY{l+s+s1}{bo}\PY{l+s+s1}{\PYZsq{}}\PY{p}{]} \PY{o}{=} \PY{n}{np}\PY{o}{.}\PY{n}{random}\PY{o}{.}\PY{n}{randn}\PY{p}{(}\PY{l+m+mi}{5}\PY{p}{,} \PY{l+m+mi}{1}\PY{p}{)}
\PY{n}{parameters\PYZus{}tmp}\PY{p}{[}\PY{l+s+s1}{\PYZsq{}}\PY{l+s+s1}{Wc}\PY{l+s+s1}{\PYZsq{}}\PY{p}{]} \PY{o}{=} \PY{n}{np}\PY{o}{.}\PY{n}{random}\PY{o}{.}\PY{n}{randn}\PY{p}{(}\PY{l+m+mi}{5}\PY{p}{,} \PY{l+m+mi}{5} \PY{o}{+} \PY{l+m+mi}{3}\PY{p}{)}
\PY{n}{parameters\PYZus{}tmp}\PY{p}{[}\PY{l+s+s1}{\PYZsq{}}\PY{l+s+s1}{bc}\PY{l+s+s1}{\PYZsq{}}\PY{p}{]} \PY{o}{=} \PY{n}{np}\PY{o}{.}\PY{n}{random}\PY{o}{.}\PY{n}{randn}\PY{p}{(}\PY{l+m+mi}{5}\PY{p}{,} \PY{l+m+mi}{1}\PY{p}{)}
\PY{n}{parameters\PYZus{}tmp}\PY{p}{[}\PY{l+s+s1}{\PYZsq{}}\PY{l+s+s1}{Wy}\PY{l+s+s1}{\PYZsq{}}\PY{p}{]} \PY{o}{=} \PY{n}{np}\PY{o}{.}\PY{n}{random}\PY{o}{.}\PY{n}{randn}\PY{p}{(}\PY{l+m+mi}{2}\PY{p}{,} \PY{l+m+mi}{5}\PY{p}{)}
\PY{n}{parameters\PYZus{}tmp}\PY{p}{[}\PY{l+s+s1}{\PYZsq{}}\PY{l+s+s1}{by}\PY{l+s+s1}{\PYZsq{}}\PY{p}{]} \PY{o}{=} \PY{n}{np}\PY{o}{.}\PY{n}{random}\PY{o}{.}\PY{n}{randn}\PY{p}{(}\PY{l+m+mi}{2}\PY{p}{,} \PY{l+m+mi}{1}\PY{p}{)}

\PY{n}{a\PYZus{}tmp}\PY{p}{,} \PY{n}{y\PYZus{}tmp}\PY{p}{,} \PY{n}{c\PYZus{}tmp}\PY{p}{,} \PY{n}{caches\PYZus{}tmp} \PY{o}{=} \PY{n}{lstm\PYZus{}forward}\PY{p}{(}\PY{n}{x\PYZus{}tmp}\PY{p}{,} \PY{n}{a0\PYZus{}tmp}\PY{p}{,} \PY{n}{parameters\PYZus{}tmp}\PY{p}{)}
\PY{n+nb}{print}\PY{p}{(}\PY{l+s+s2}{\PYZdq{}}\PY{l+s+s2}{a[4][3][6] = }\PY{l+s+s2}{\PYZdq{}}\PY{p}{,} \PY{n}{a\PYZus{}tmp}\PY{p}{[}\PY{l+m+mi}{4}\PY{p}{]}\PY{p}{[}\PY{l+m+mi}{3}\PY{p}{]}\PY{p}{[}\PY{l+m+mi}{6}\PY{p}{]}\PY{p}{)}
\PY{n+nb}{print}\PY{p}{(}\PY{l+s+s2}{\PYZdq{}}\PY{l+s+s2}{a.shape = }\PY{l+s+s2}{\PYZdq{}}\PY{p}{,} \PY{n}{a\PYZus{}tmp}\PY{o}{.}\PY{n}{shape}\PY{p}{)}
\PY{n+nb}{print}\PY{p}{(}\PY{l+s+s2}{\PYZdq{}}\PY{l+s+s2}{y[1][4][3] =}\PY{l+s+s2}{\PYZdq{}}\PY{p}{,} \PY{n}{y\PYZus{}tmp}\PY{p}{[}\PY{l+m+mi}{1}\PY{p}{]}\PY{p}{[}\PY{l+m+mi}{4}\PY{p}{]}\PY{p}{[}\PY{l+m+mi}{3}\PY{p}{]}\PY{p}{)}
\PY{n+nb}{print}\PY{p}{(}\PY{l+s+s2}{\PYZdq{}}\PY{l+s+s2}{y.shape = }\PY{l+s+s2}{\PYZdq{}}\PY{p}{,} \PY{n}{y\PYZus{}tmp}\PY{o}{.}\PY{n}{shape}\PY{p}{)}
\PY{n+nb}{print}\PY{p}{(}\PY{l+s+s2}{\PYZdq{}}\PY{l+s+s2}{caches[1][1][1] =}\PY{l+s+se}{\PYZbs{}n}\PY{l+s+s2}{\PYZdq{}}\PY{p}{,} \PY{n}{caches\PYZus{}tmp}\PY{p}{[}\PY{l+m+mi}{1}\PY{p}{]}\PY{p}{[}\PY{l+m+mi}{1}\PY{p}{]}\PY{p}{[}\PY{l+m+mi}{1}\PY{p}{]}\PY{p}{)}
\PY{n+nb}{print}\PY{p}{(}\PY{l+s+s2}{\PYZdq{}}\PY{l+s+s2}{c[1][2][1]}\PY{l+s+s2}{\PYZdq{}}\PY{p}{,} \PY{n}{c\PYZus{}tmp}\PY{p}{[}\PY{l+m+mi}{1}\PY{p}{]}\PY{p}{[}\PY{l+m+mi}{2}\PY{p}{]}\PY{p}{[}\PY{l+m+mi}{1}\PY{p}{]}\PY{p}{)}
\PY{n+nb}{print}\PY{p}{(}\PY{l+s+s2}{\PYZdq{}}\PY{l+s+s2}{len(caches) = }\PY{l+s+s2}{\PYZdq{}}\PY{p}{,} \PY{n+nb}{len}\PY{p}{(}\PY{n}{caches\PYZus{}tmp}\PY{p}{)}\PY{p}{)}

\PY{c+c1}{\PYZsh{} UNIT TEST    }
\PY{n}{lstm\PYZus{}forward\PYZus{}test}\PY{p}{(}\PY{n}{lstm\PYZus{}forward}\PY{p}{)}
\end{Verbatim}
\end{tcolorbox}

    \begin{Verbatim}[commandchars=\\\{\}]
a[4][3][6] =  0.17211776753291672
a.shape =  (5, 10, 7)
y[1][4][3] = 0.9508734618501101
y.shape =  (2, 10, 7)
caches[1][1][1] =
 [ 0.82797464  0.23009474  0.76201118 -0.22232814 -0.20075807  0.18656139
  0.41005165]
c[1][2][1] -0.8555449167181981
len(caches) =  2
\textcolor{ansi-green-intense}{All tests passed}
    \end{Verbatim}

    \textbf{Expected Output}:

\begin{Shaded}
\begin{Highlighting}[]
\NormalTok{a[}\DecValTok{4}\NormalTok{][}\DecValTok{3}\NormalTok{][}\DecValTok{6}\NormalTok{] }\OperatorTok{=}  \FloatTok{0.172117767533}
\NormalTok{a.shape }\OperatorTok{=}\NormalTok{  (}\DecValTok{5}\NormalTok{, }\DecValTok{10}\NormalTok{, }\DecValTok{7}\NormalTok{)}
\NormalTok{y[}\DecValTok{1}\NormalTok{][}\DecValTok{4}\NormalTok{][}\DecValTok{3}\NormalTok{] }\OperatorTok{=} \FloatTok{0.95087346185}
\NormalTok{y.shape }\OperatorTok{=}\NormalTok{  (}\DecValTok{2}\NormalTok{, }\DecValTok{10}\NormalTok{, }\DecValTok{7}\NormalTok{)}
\NormalTok{caches[}\DecValTok{1}\NormalTok{][}\DecValTok{1}\NormalTok{][}\DecValTok{1}\NormalTok{] }\OperatorTok{=}
\NormalTok{ [ }\FloatTok{0.82797464}  \FloatTok{0.23009474}  \FloatTok{0.76201118} \OperatorTok{{-}}\FloatTok{0.22232814} \OperatorTok{{-}}\FloatTok{0.20075807}  \FloatTok{0.18656139}
  \FloatTok{0.41005165}\NormalTok{]}
\NormalTok{c[}\DecValTok{1}\NormalTok{][}\DecValTok{2}\NormalTok{][}\DecValTok{1}\NormalTok{] }\OperatorTok{{-}}\FloatTok{0.855544916718}
\BuiltInTok{len}\NormalTok{(caches) }\OperatorTok{=}  \DecValTok{2}
\end{Highlighting}
\end{Shaded}

    \hypertarget{congratulations}{%
\subsubsection{Congratulations!}\label{congratulations}}

You have now implemented the forward passes for both the basic RNN and
the LSTM. When using a deep learning framework, implementing the forward
pass is sufficient to build systems that achieve great performance. The
framework will take care of the rest.

What you should remember:

\begin{itemize}
\tightlist
\item
  An LSTM is similar to an RNN in that they both use hidden states to
  pass along information, but an LSTM also uses a cell state, which is
  like a long-term memory, to help deal with the issue of vanishing
  gradients
\item
  An LSTM cell consists of a cell state, or long-term memory, a hidden
  state, or short-term memory, along with 3 gates that constantly update
  the relevancy of its inputs:

  \begin{itemize}
  \tightlist
  \item
    A forget gate, which decides which input units should be remembered
    and passed along. It's a tensor with values between 0 and 1.

    \begin{itemize}
    \tightlist
    \item
      If a unit has a value close to 0, the LSTM will ``forget'' the
      stored state in the previous cell state.
    \item
      If it has a value close to 1, the LSTM will mostly remember the
      corresponding value.
    \end{itemize}
  \item
    An update gate, again a tensor containing values between 0 and 1. It
    decides on what information to throw away, and what new information
    to add.

    \begin{itemize}
    \tightlist
    \item
      When a unit in the update gate is close to 1, the value of its
      candidate is passed on to the hidden state.
    \item
      When a unit in the update gate is close to 0, it's prevented from
      being passed onto the hidden state.
    \end{itemize}
  \item
    And an output gate, which decides what gets sent as the output of
    the time step
  \end{itemize}
\end{itemize}

Let's recap all you've accomplished so far. You have:

\begin{itemize}
\tightlist
\item
  Used notation for building sequence models
\item
  Become familiar with the architecture of a basic RNN and an LSTM, and
  can describe their components
\end{itemize}

The rest of this notebook is optional, and will not be graded, but as
always, you are encouraged to push your own understanding! Good luck and
have fun.

    ~\\
\#\# 3 - Backpropagation in Recurrent Neural Networks (OPTIONAL /
UNGRADED)

In modern deep learning frameworks, you only have to implement the
forward pass, and the framework takes care of the backward pass, so most
deep learning engineers do not need to bother with the details of the
backward pass. If, however, you are an expert in calculus (or are just
curious) and want to see the details of backprop in RNNs, you can work
through this optional portion of the notebook.

When in an earlier
\href{https://www.coursera.org/learn/neural-networks-deep-learning/lecture/0VSHe/derivatives-with-a-computation-graph}{course}
you implemented a simple (fully connected) neural network, you used
backpropagation to compute the derivatives with respect to the cost to
update the parameters. Similarly, in recurrent neural networks you can
calculate the derivatives with respect to the cost in order to update
the parameters. The backprop equations are quite complicated, and so
were not derived in lecture. However, they're briefly presented for your
viewing pleasure below.

    Note that this notebook does not implement the backward path from the
Loss `J' backwards to `a'. This would have included the dense layer and
softmax, which are a part of the forward path. This is assumed to be
calculated elsewhere and the result passed to \texttt{rnn\_backward} in
`da'. It is further assumed that loss has been adjusted for batch size
(m) and division by the number of examples is not required here.

    This section is optional and ungraded, because it's more difficult and
has fewer details regarding its implementation. Note that this section
only implements key elements of the full path!

Onward, brave one:

    ~\\
\#\#\# 3.1 - Basic RNN Backward Pass

Begin by computing the backward pass for the basic RNN cell. Then, in
the following sections, iterate through the cells.

Figure 6: The RNN cell's backward pass. Just like in a fully-connected
neural network, the derivative of the cost function \(J\) backpropagates
through the time steps of the RNN by following the chain rule from
calculus. Internal to the cell, the chain rule is also used to calculate
\((\frac{\partial J}{\partial W_{ax}},\frac{\partial J}{\partial W_{aa}},\frac{\partial J}{\partial b})\)
to update the parameters \((W_{ax}, W_{aa}, b_a)\). The operation can
utilize the cached results from the forward path.

    Recall from lecture that the shorthand for the partial derivative of
cost relative to a variable is \texttt{dVariable}. For example,
\(\frac{\partial J}{\partial W_{ax}}\) is \(dW_{ax}\). This will be used
throughout the remaining sections.

    Figure 7: This implementation of \texttt{rnn\_cell\_backward} does
\textbf{not} include the output dense layer and softmax which are
included in \texttt{rnn\_cell\_forward}.

\(da_{next}\) is \(\frac{\partial{J}}{\partial a^{\langle t \rangle}}\)
and includes loss from previous stages and current stage output logic.
The addition shown in green will be part of your implementation of
\texttt{rnn\_backward}.

    \hypertarget{equations}{%
\subparagraph{Equations}\label{equations}}

To compute \texttt{rnn\_cell\_backward}, you can use the following
equations. It's a good exercise to derive them by hand. Here, \(*\)
denotes element-wise multiplication while the absence of a symbol
indicates matrix multiplication.

\begin{align}
\displaystyle a^{\langle t \rangle} &= \tanh(W_{ax} x^{\langle t \rangle} + W_{aa} a^{\langle t-1 \rangle} + b_{a})\tag{-} \\[8pt]
\displaystyle \frac{\partial \tanh(x)} {\partial x} &= 1 - \tanh^2(x) \tag{-} \\[8pt]
\displaystyle {dtanh} &= da_{next} * ( 1 - \tanh^2(W_{ax}x^{\langle t \rangle}+W_{aa} a^{\langle t-1 \rangle} + b_{a})) \tag{0} \\[8pt]
\displaystyle  {dW_{ax}} &= dtanh \cdot x^{\langle t \rangle T}\tag{1} \\[8pt]
\displaystyle dW_{aa} &= dtanh \cdot a^{\langle t-1 \rangle T}\tag{2} \\[8pt]
\displaystyle db_a& = \sum_{batch}dtanh\tag{3} \\[8pt]
\displaystyle dx^{\langle t \rangle} &= { W_{ax}}^T \cdot dtanh\tag{4} \\[8pt]
\displaystyle da_{prev} &= { W_{aa}}^T \cdot dtanh\tag{5}
\end{align}

    \#\#\# Exercise 5 - rnn\_cell\_backward

Implementing \texttt{rnn\_cell\_backward}.

The results can be computed directly by implementing the equations
above. However, you have an option to simplify them by computing `dz'
and using the chain rule.\\
This can be further simplified by noting that
\(\tanh(W_{ax}x^{\langle t \rangle}+W_{aa} a^{\langle t-1 \rangle} + b_{a})\)
was computed and saved as \texttt{a\_next} in the forward pass.

To calculate \texttt{dba}, the `batch' above is a sum across all `m'
examples (axis= 1). Note that you should use the
\texttt{keepdims\ =\ True} option.

It may be worthwhile to review Course 1
\href{https://www.coursera.org/learn/neural-networks-deep-learning/lecture/0VSHe/derivatives-with-a-computation-graph}{Derivatives
with a Computation Graph} through
\href{https://www.coursera.org/learn/neural-networks-deep-learning/lecture/6dDj7/backpropagation-intuition-optional}{Backpropagation
Intuition}, which decompose the calculation into steps using the chain
rule.\\
Matrix vector derivatives are described
\href{http://cs231n.stanford.edu/vecDerivs.pdf}{here}, though the
equations above incorporate the required transformations.

\textbf{Note}: \texttt{rnn\_cell\_backward} does \textbf{not} include
the calculation of loss from \(y \langle t \rangle\). This is
incorporated into the incoming \texttt{da\_next}. This is a slight
mismatch with \texttt{rnn\_cell\_forward}, which includes a dense layer
and softmax.

\textbf{Note on the code}:

\(\displaystyle dx^{\langle t \rangle}\) is represented by dxt,\\
\(\displaystyle d W_{ax}\) is represented by dWax,\\
\(\displaystyle da_{prev}\) is represented by da\_prev,\\
\(\displaystyle dW_{aa}\) is represented by dWaa,\\
\(\displaystyle db_{a}\) is represented by dba,\\
\texttt{dz} is not derived above but can optionally be derived by
students to simplify the repeated calculations.

    \begin{tcolorbox}[breakable, size=fbox, boxrule=1pt, pad at break*=1mm,colback=cellbackground, colframe=cellborder]
\prompt{In}{incolor}{52}{\boxspacing}
\begin{Verbatim}[commandchars=\\\{\}]
\PY{c+c1}{\PYZsh{} UNGRADED FUNCTION: rnn\PYZus{}cell\PYZus{}backward}

\PY{k}{def} \PY{n+nf}{rnn\PYZus{}cell\PYZus{}backward}\PY{p}{(}\PY{n}{da\PYZus{}next}\PY{p}{,} \PY{n}{cache}\PY{p}{)}\PY{p}{:}
    \PY{l+s+sd}{\PYZdq{}\PYZdq{}\PYZdq{}}
\PY{l+s+sd}{    Implements the backward pass for the RNN\PYZhy{}cell (single time\PYZhy{}step).}

\PY{l+s+sd}{    Arguments:}
\PY{l+s+sd}{    da\PYZus{}next \PYZhy{}\PYZhy{} Gradient of loss with respect to next hidden state}
\PY{l+s+sd}{    cache \PYZhy{}\PYZhy{} python dictionary containing useful values (output of rnn\PYZus{}cell\PYZus{}forward())}

\PY{l+s+sd}{    Returns:}
\PY{l+s+sd}{    gradients \PYZhy{}\PYZhy{} python dictionary containing:}
\PY{l+s+sd}{                        dx \PYZhy{}\PYZhy{} Gradients of input data, of shape (n\PYZus{}x, m)}
\PY{l+s+sd}{                        da\PYZus{}prev \PYZhy{}\PYZhy{} Gradients of previous hidden state, of shape (n\PYZus{}a, m)}
\PY{l+s+sd}{                        dWax \PYZhy{}\PYZhy{} Gradients of input\PYZhy{}to\PYZhy{}hidden weights, of shape (n\PYZus{}a, n\PYZus{}x)}
\PY{l+s+sd}{                        dWaa \PYZhy{}\PYZhy{} Gradients of hidden\PYZhy{}to\PYZhy{}hidden weights, of shape (n\PYZus{}a, n\PYZus{}a)}
\PY{l+s+sd}{                        dba \PYZhy{}\PYZhy{} Gradients of bias vector, of shape (n\PYZus{}a, 1)}
\PY{l+s+sd}{    \PYZdq{}\PYZdq{}\PYZdq{}}
    
    \PY{c+c1}{\PYZsh{} Retrieve values from cache}
    \PY{p}{(}\PY{n}{a\PYZus{}next}\PY{p}{,} \PY{n}{a\PYZus{}prev}\PY{p}{,} \PY{n}{xt}\PY{p}{,} \PY{n}{parameters}\PY{p}{)} \PY{o}{=} \PY{n}{cache}
    
    \PY{c+c1}{\PYZsh{} Retrieve values from parameters}
    \PY{n}{Wax} \PY{o}{=} \PY{n}{parameters}\PY{p}{[}\PY{l+s+s2}{\PYZdq{}}\PY{l+s+s2}{Wax}\PY{l+s+s2}{\PYZdq{}}\PY{p}{]}
    \PY{n}{Waa} \PY{o}{=} \PY{n}{parameters}\PY{p}{[}\PY{l+s+s2}{\PYZdq{}}\PY{l+s+s2}{Waa}\PY{l+s+s2}{\PYZdq{}}\PY{p}{]}
    \PY{n}{Wya} \PY{o}{=} \PY{n}{parameters}\PY{p}{[}\PY{l+s+s2}{\PYZdq{}}\PY{l+s+s2}{Wya}\PY{l+s+s2}{\PYZdq{}}\PY{p}{]}
    \PY{n}{ba} \PY{o}{=} \PY{n}{parameters}\PY{p}{[}\PY{l+s+s2}{\PYZdq{}}\PY{l+s+s2}{ba}\PY{l+s+s2}{\PYZdq{}}\PY{p}{]}
    \PY{n}{by} \PY{o}{=} \PY{n}{parameters}\PY{p}{[}\PY{l+s+s2}{\PYZdq{}}\PY{l+s+s2}{by}\PY{l+s+s2}{\PYZdq{}}\PY{p}{]}

    \PY{c+c1}{\PYZsh{}\PYZsh{}\PYZsh{} START CODE HERE \PYZsh{}\PYZsh{}\PYZsh{}}
    \PY{c+c1}{\PYZsh{} compute the gradient of dtanh term using a\PYZus{}next and da\PYZus{}next (≈1 line)}
    \PY{n}{dtanh} \PY{o}{=} \PY{n}{da\PYZus{}next} \PY{o}{*} \PY{p}{(}\PY{l+m+mi}{1} \PY{o}{\PYZhy{}} \PY{n}{a\PYZus{}next} \PY{o}{*} \PY{n}{a\PYZus{}next}\PY{p}{)}

    \PY{c+c1}{\PYZsh{} compute the gradient of the loss with respect to Wax (≈2 lines)}
    \PY{n}{dxt} \PY{o}{=} \PY{n}{np}\PY{o}{.}\PY{n}{matmul}\PY{p}{(}\PY{n}{np}\PY{o}{.}\PY{n}{transpose}\PY{p}{(}\PY{n}{Wax}\PY{p}{)}\PY{p}{,} \PY{n}{dtanh}\PY{p}{)}
    \PY{n}{dWax} \PY{o}{=} \PY{n}{np}\PY{o}{.}\PY{n}{matmul}\PY{p}{(}\PY{n}{dtanh}\PY{p}{,} \PY{n}{np}\PY{o}{.}\PY{n}{transpose}\PY{p}{(}\PY{n}{xt}\PY{p}{)}\PY{p}{)}

    \PY{c+c1}{\PYZsh{} compute the gradient with respect to Waa (≈2 lines)}
    \PY{n}{da\PYZus{}prev} \PY{o}{=} \PY{n}{np}\PY{o}{.}\PY{n}{matmul}\PY{p}{(}\PY{n}{np}\PY{o}{.}\PY{n}{transpose}\PY{p}{(}\PY{n}{Waa}\PY{p}{)}\PY{p}{,} \PY{n}{dtanh}\PY{p}{)}
    \PY{n}{dWaa} \PY{o}{=} \PY{n}{np}\PY{o}{.}\PY{n}{matmul}\PY{p}{(}\PY{n}{dtanh}\PY{p}{,} \PY{n}{np}\PY{o}{.}\PY{n}{transpose}\PY{p}{(}\PY{n}{a\PYZus{}prev}\PY{p}{)}\PY{p}{)}

    \PY{c+c1}{\PYZsh{} compute the gradient with respect to b (≈1 line)}
    \PY{n}{dba} \PY{o}{=} \PY{n}{np}\PY{o}{.}\PY{n}{sum}\PY{p}{(}\PY{n}{dtanh}\PY{p}{,} \PY{n}{axis}\PY{o}{=}\PY{l+m+mi}{1}\PY{p}{,} \PY{n}{keepdims}\PY{o}{=}\PY{k+kc}{True}\PY{p}{)}

    \PY{c+c1}{\PYZsh{}\PYZsh{}\PYZsh{} END CODE HERE \PYZsh{}\PYZsh{}\PYZsh{}}
    
    \PY{c+c1}{\PYZsh{} Store the gradients in a python dictionary}
    \PY{n}{gradients} \PY{o}{=} \PY{p}{\PYZob{}}\PY{l+s+s2}{\PYZdq{}}\PY{l+s+s2}{dxt}\PY{l+s+s2}{\PYZdq{}}\PY{p}{:} \PY{n}{dxt}\PY{p}{,} \PY{l+s+s2}{\PYZdq{}}\PY{l+s+s2}{da\PYZus{}prev}\PY{l+s+s2}{\PYZdq{}}\PY{p}{:} \PY{n}{da\PYZus{}prev}\PY{p}{,} \PY{l+s+s2}{\PYZdq{}}\PY{l+s+s2}{dWax}\PY{l+s+s2}{\PYZdq{}}\PY{p}{:} \PY{n}{dWax}\PY{p}{,} \PY{l+s+s2}{\PYZdq{}}\PY{l+s+s2}{dWaa}\PY{l+s+s2}{\PYZdq{}}\PY{p}{:} \PY{n}{dWaa}\PY{p}{,} \PY{l+s+s2}{\PYZdq{}}\PY{l+s+s2}{dba}\PY{l+s+s2}{\PYZdq{}}\PY{p}{:} \PY{n}{dba}\PY{p}{\PYZcb{}}
    
    \PY{k}{return} \PY{n}{gradients}
\end{Verbatim}
\end{tcolorbox}

    \begin{tcolorbox}[breakable, size=fbox, boxrule=1pt, pad at break*=1mm,colback=cellbackground, colframe=cellborder]
\prompt{In}{incolor}{53}{\boxspacing}
\begin{Verbatim}[commandchars=\\\{\}]
\PY{n}{np}\PY{o}{.}\PY{n}{random}\PY{o}{.}\PY{n}{seed}\PY{p}{(}\PY{l+m+mi}{1}\PY{p}{)}
\PY{n}{xt\PYZus{}tmp} \PY{o}{=} \PY{n}{np}\PY{o}{.}\PY{n}{random}\PY{o}{.}\PY{n}{randn}\PY{p}{(}\PY{l+m+mi}{3}\PY{p}{,}\PY{l+m+mi}{10}\PY{p}{)}
\PY{n}{a\PYZus{}prev\PYZus{}tmp} \PY{o}{=} \PY{n}{np}\PY{o}{.}\PY{n}{random}\PY{o}{.}\PY{n}{randn}\PY{p}{(}\PY{l+m+mi}{5}\PY{p}{,}\PY{l+m+mi}{10}\PY{p}{)}
\PY{n}{parameters\PYZus{}tmp} \PY{o}{=} \PY{p}{\PYZob{}}\PY{p}{\PYZcb{}}
\PY{n}{parameters\PYZus{}tmp}\PY{p}{[}\PY{l+s+s1}{\PYZsq{}}\PY{l+s+s1}{Wax}\PY{l+s+s1}{\PYZsq{}}\PY{p}{]} \PY{o}{=} \PY{n}{np}\PY{o}{.}\PY{n}{random}\PY{o}{.}\PY{n}{randn}\PY{p}{(}\PY{l+m+mi}{5}\PY{p}{,}\PY{l+m+mi}{3}\PY{p}{)}
\PY{n}{parameters\PYZus{}tmp}\PY{p}{[}\PY{l+s+s1}{\PYZsq{}}\PY{l+s+s1}{Waa}\PY{l+s+s1}{\PYZsq{}}\PY{p}{]} \PY{o}{=} \PY{n}{np}\PY{o}{.}\PY{n}{random}\PY{o}{.}\PY{n}{randn}\PY{p}{(}\PY{l+m+mi}{5}\PY{p}{,}\PY{l+m+mi}{5}\PY{p}{)}
\PY{n}{parameters\PYZus{}tmp}\PY{p}{[}\PY{l+s+s1}{\PYZsq{}}\PY{l+s+s1}{Wya}\PY{l+s+s1}{\PYZsq{}}\PY{p}{]} \PY{o}{=} \PY{n}{np}\PY{o}{.}\PY{n}{random}\PY{o}{.}\PY{n}{randn}\PY{p}{(}\PY{l+m+mi}{2}\PY{p}{,}\PY{l+m+mi}{5}\PY{p}{)}
\PY{n}{parameters\PYZus{}tmp}\PY{p}{[}\PY{l+s+s1}{\PYZsq{}}\PY{l+s+s1}{ba}\PY{l+s+s1}{\PYZsq{}}\PY{p}{]} \PY{o}{=} \PY{n}{np}\PY{o}{.}\PY{n}{random}\PY{o}{.}\PY{n}{randn}\PY{p}{(}\PY{l+m+mi}{5}\PY{p}{,}\PY{l+m+mi}{1}\PY{p}{)}
\PY{n}{parameters\PYZus{}tmp}\PY{p}{[}\PY{l+s+s1}{\PYZsq{}}\PY{l+s+s1}{by}\PY{l+s+s1}{\PYZsq{}}\PY{p}{]} \PY{o}{=} \PY{n}{np}\PY{o}{.}\PY{n}{random}\PY{o}{.}\PY{n}{randn}\PY{p}{(}\PY{l+m+mi}{2}\PY{p}{,}\PY{l+m+mi}{1}\PY{p}{)}

\PY{n}{a\PYZus{}next\PYZus{}tmp}\PY{p}{,} \PY{n}{yt\PYZus{}tmp}\PY{p}{,} \PY{n}{cache\PYZus{}tmp} \PY{o}{=} \PY{n}{rnn\PYZus{}cell\PYZus{}forward}\PY{p}{(}\PY{n}{xt\PYZus{}tmp}\PY{p}{,} \PY{n}{a\PYZus{}prev\PYZus{}tmp}\PY{p}{,} \PY{n}{parameters\PYZus{}tmp}\PY{p}{)}

\PY{n}{da\PYZus{}next\PYZus{}tmp} \PY{o}{=} \PY{n}{np}\PY{o}{.}\PY{n}{random}\PY{o}{.}\PY{n}{randn}\PY{p}{(}\PY{l+m+mi}{5}\PY{p}{,}\PY{l+m+mi}{10}\PY{p}{)}
\PY{n}{gradients\PYZus{}tmp} \PY{o}{=} \PY{n}{rnn\PYZus{}cell\PYZus{}backward}\PY{p}{(}\PY{n}{da\PYZus{}next\PYZus{}tmp}\PY{p}{,} \PY{n}{cache\PYZus{}tmp}\PY{p}{)}
\PY{n+nb}{print}\PY{p}{(}\PY{l+s+s2}{\PYZdq{}}\PY{l+s+s2}{gradients[}\PY{l+s+se}{\PYZbs{}\PYZdq{}}\PY{l+s+s2}{dxt}\PY{l+s+se}{\PYZbs{}\PYZdq{}}\PY{l+s+s2}{][1][2] =}\PY{l+s+s2}{\PYZdq{}}\PY{p}{,} \PY{n}{gradients\PYZus{}tmp}\PY{p}{[}\PY{l+s+s2}{\PYZdq{}}\PY{l+s+s2}{dxt}\PY{l+s+s2}{\PYZdq{}}\PY{p}{]}\PY{p}{[}\PY{l+m+mi}{1}\PY{p}{]}\PY{p}{[}\PY{l+m+mi}{2}\PY{p}{]}\PY{p}{)}
\PY{n+nb}{print}\PY{p}{(}\PY{l+s+s2}{\PYZdq{}}\PY{l+s+s2}{gradients[}\PY{l+s+se}{\PYZbs{}\PYZdq{}}\PY{l+s+s2}{dxt}\PY{l+s+se}{\PYZbs{}\PYZdq{}}\PY{l+s+s2}{].shape =}\PY{l+s+s2}{\PYZdq{}}\PY{p}{,} \PY{n}{gradients\PYZus{}tmp}\PY{p}{[}\PY{l+s+s2}{\PYZdq{}}\PY{l+s+s2}{dxt}\PY{l+s+s2}{\PYZdq{}}\PY{p}{]}\PY{o}{.}\PY{n}{shape}\PY{p}{)}
\PY{n+nb}{print}\PY{p}{(}\PY{l+s+s2}{\PYZdq{}}\PY{l+s+s2}{gradients[}\PY{l+s+se}{\PYZbs{}\PYZdq{}}\PY{l+s+s2}{da\PYZus{}prev}\PY{l+s+se}{\PYZbs{}\PYZdq{}}\PY{l+s+s2}{][2][3] =}\PY{l+s+s2}{\PYZdq{}}\PY{p}{,} \PY{n}{gradients\PYZus{}tmp}\PY{p}{[}\PY{l+s+s2}{\PYZdq{}}\PY{l+s+s2}{da\PYZus{}prev}\PY{l+s+s2}{\PYZdq{}}\PY{p}{]}\PY{p}{[}\PY{l+m+mi}{2}\PY{p}{]}\PY{p}{[}\PY{l+m+mi}{3}\PY{p}{]}\PY{p}{)}
\PY{n+nb}{print}\PY{p}{(}\PY{l+s+s2}{\PYZdq{}}\PY{l+s+s2}{gradients[}\PY{l+s+se}{\PYZbs{}\PYZdq{}}\PY{l+s+s2}{da\PYZus{}prev}\PY{l+s+se}{\PYZbs{}\PYZdq{}}\PY{l+s+s2}{].shape =}\PY{l+s+s2}{\PYZdq{}}\PY{p}{,} \PY{n}{gradients\PYZus{}tmp}\PY{p}{[}\PY{l+s+s2}{\PYZdq{}}\PY{l+s+s2}{da\PYZus{}prev}\PY{l+s+s2}{\PYZdq{}}\PY{p}{]}\PY{o}{.}\PY{n}{shape}\PY{p}{)}
\PY{n+nb}{print}\PY{p}{(}\PY{l+s+s2}{\PYZdq{}}\PY{l+s+s2}{gradients[}\PY{l+s+se}{\PYZbs{}\PYZdq{}}\PY{l+s+s2}{dWax}\PY{l+s+se}{\PYZbs{}\PYZdq{}}\PY{l+s+s2}{][3][1] =}\PY{l+s+s2}{\PYZdq{}}\PY{p}{,} \PY{n}{gradients\PYZus{}tmp}\PY{p}{[}\PY{l+s+s2}{\PYZdq{}}\PY{l+s+s2}{dWax}\PY{l+s+s2}{\PYZdq{}}\PY{p}{]}\PY{p}{[}\PY{l+m+mi}{3}\PY{p}{]}\PY{p}{[}\PY{l+m+mi}{1}\PY{p}{]}\PY{p}{)}
\PY{n+nb}{print}\PY{p}{(}\PY{l+s+s2}{\PYZdq{}}\PY{l+s+s2}{gradients[}\PY{l+s+se}{\PYZbs{}\PYZdq{}}\PY{l+s+s2}{dWax}\PY{l+s+se}{\PYZbs{}\PYZdq{}}\PY{l+s+s2}{].shape =}\PY{l+s+s2}{\PYZdq{}}\PY{p}{,} \PY{n}{gradients\PYZus{}tmp}\PY{p}{[}\PY{l+s+s2}{\PYZdq{}}\PY{l+s+s2}{dWax}\PY{l+s+s2}{\PYZdq{}}\PY{p}{]}\PY{o}{.}\PY{n}{shape}\PY{p}{)}
\PY{n+nb}{print}\PY{p}{(}\PY{l+s+s2}{\PYZdq{}}\PY{l+s+s2}{gradients[}\PY{l+s+se}{\PYZbs{}\PYZdq{}}\PY{l+s+s2}{dWaa}\PY{l+s+se}{\PYZbs{}\PYZdq{}}\PY{l+s+s2}{][1][2] =}\PY{l+s+s2}{\PYZdq{}}\PY{p}{,} \PY{n}{gradients\PYZus{}tmp}\PY{p}{[}\PY{l+s+s2}{\PYZdq{}}\PY{l+s+s2}{dWaa}\PY{l+s+s2}{\PYZdq{}}\PY{p}{]}\PY{p}{[}\PY{l+m+mi}{1}\PY{p}{]}\PY{p}{[}\PY{l+m+mi}{2}\PY{p}{]}\PY{p}{)}
\PY{n+nb}{print}\PY{p}{(}\PY{l+s+s2}{\PYZdq{}}\PY{l+s+s2}{gradients[}\PY{l+s+se}{\PYZbs{}\PYZdq{}}\PY{l+s+s2}{dWaa}\PY{l+s+se}{\PYZbs{}\PYZdq{}}\PY{l+s+s2}{].shape =}\PY{l+s+s2}{\PYZdq{}}\PY{p}{,} \PY{n}{gradients\PYZus{}tmp}\PY{p}{[}\PY{l+s+s2}{\PYZdq{}}\PY{l+s+s2}{dWaa}\PY{l+s+s2}{\PYZdq{}}\PY{p}{]}\PY{o}{.}\PY{n}{shape}\PY{p}{)}
\PY{n+nb}{print}\PY{p}{(}\PY{l+s+s2}{\PYZdq{}}\PY{l+s+s2}{gradients[}\PY{l+s+se}{\PYZbs{}\PYZdq{}}\PY{l+s+s2}{dba}\PY{l+s+se}{\PYZbs{}\PYZdq{}}\PY{l+s+s2}{][4] =}\PY{l+s+s2}{\PYZdq{}}\PY{p}{,} \PY{n}{gradients\PYZus{}tmp}\PY{p}{[}\PY{l+s+s2}{\PYZdq{}}\PY{l+s+s2}{dba}\PY{l+s+s2}{\PYZdq{}}\PY{p}{]}\PY{p}{[}\PY{l+m+mi}{4}\PY{p}{]}\PY{p}{)}
\PY{n+nb}{print}\PY{p}{(}\PY{l+s+s2}{\PYZdq{}}\PY{l+s+s2}{gradients[}\PY{l+s+se}{\PYZbs{}\PYZdq{}}\PY{l+s+s2}{dba}\PY{l+s+se}{\PYZbs{}\PYZdq{}}\PY{l+s+s2}{].shape =}\PY{l+s+s2}{\PYZdq{}}\PY{p}{,} \PY{n}{gradients\PYZus{}tmp}\PY{p}{[}\PY{l+s+s2}{\PYZdq{}}\PY{l+s+s2}{dba}\PY{l+s+s2}{\PYZdq{}}\PY{p}{]}\PY{o}{.}\PY{n}{shape}\PY{p}{)}
\end{Verbatim}
\end{tcolorbox}

    \begin{Verbatim}[commandchars=\\\{\}]
gradients["dxt"][1][2] = -1.3872130506020925
gradients["dxt"].shape = (3, 10)
gradients["da\_prev"][2][3] = -0.15239949377395495
gradients["da\_prev"].shape = (5, 10)
gradients["dWax"][3][1] = 0.4107728249354584
gradients["dWax"].shape = (5, 3)
gradients["dWaa"][1][2] = 1.1503450668497135
gradients["dWaa"].shape = (5, 5)
gradients["dba"][4] = [0.20023491]
gradients["dba"].shape = (5, 1)
    \end{Verbatim}

    \textbf{Expected Output}:

gradients{[}``dxt''{]}{[}1{]}{[}2{]} =

-1.3872130506

\begin{verbatim}
<tr>
    <td>
        <b>gradients["dxt"].shape</b> =
    </td>
    <td>
       (3, 10)
    </td>
</tr>
    <tr>
    <td>
        <b>gradients["da_prev"][2][3]</b> =
    </td>
    <td>
       -0.152399493774
    </td>
</tr>
    <tr>
    <td>
        <b>gradients["da_prev"].shape</b> =
    </td>
    <td>
       (5, 10)
    </td>
</tr>
    <tr>
    <td>
        <b>gradients["dWax"][3][1]</b> =
    </td>
    <td>
       0.410772824935
    </td>
</tr>
        <tr>
    <td>
        <b>gradients["dWax"].shape</b> =
    </td>
    <td>
       (5, 3)
    </td>
</tr>
    <tr>
    <td>
        <b>gradients["dWaa"][1][2]</b> = 
    </td>
    <td>
       1.15034506685
    </td>
</tr>
    <tr>
    <td>
        <b>gradients["dWaa"].shape</b> =
    </td>
    <td>
       (5, 5)
    </td>
</tr>
    <tr>
    <td>
        <b>gradients["dba"][4]</b> = 
    </td>
    <td>
       [ 0.20023491]
    </td>
</tr>
    <tr>
    <td>
        <b>gradients["dba"].shape</b> = 
    </td>
    <td>
       (5, 1)
    </td>
</tr>
\end{verbatim}

    \#\#\# Exercise 6 - rnn\_backward

Computing the gradients of the cost with respect to
\(a^{\langle t \rangle}\) at every time step \(t\) is useful because it
is what helps the gradient backpropagate to the previous RNN cell. To do
so, you need to iterate through all the time steps starting at the end,
and at each step, you increment the overall \(db_a\), \(dW_{aa}\),
\(dW_{ax}\) and you store \(dx\).

\textbf{Instructions}:

Implement the \texttt{rnn\_backward} function. Initialize the return
variables with zeros first, and then loop through all the time steps
while calling the \texttt{rnn\_cell\_backward} at each time time step,
updating the other variables accordingly.

    \begin{itemize}
\tightlist
\item
  Note that this notebook does not implement the backward path from the
  Loss `J' backwards to `a'.

  \begin{itemize}
  \tightlist
  \item
    This would have included the dense layer and softmax which are a
    part of the forward path.
  \item
    This is assumed to be calculated elsewhere and the result passed to
    \texttt{rnn\_backward} in `da'.
  \item
    You must combine this with the loss from the previous stages when
    calling \texttt{rnn\_cell\_backward} (see figure 7 above).
  \end{itemize}
\item
  It is further assumed that loss has been adjusted for batch size (m).

  \begin{itemize}
  \tightlist
  \item
    Therefore, division by the number of examples is not required here.
  \end{itemize}
\end{itemize}

    \begin{tcolorbox}[breakable, size=fbox, boxrule=1pt, pad at break*=1mm,colback=cellbackground, colframe=cellborder]
\prompt{In}{incolor}{58}{\boxspacing}
\begin{Verbatim}[commandchars=\\\{\}]
\PY{c+c1}{\PYZsh{} UNGRADED FUNCTION: rnn\PYZus{}backward}

\PY{k}{def} \PY{n+nf}{rnn\PYZus{}backward}\PY{p}{(}\PY{n}{da}\PY{p}{,} \PY{n}{caches}\PY{p}{)}\PY{p}{:}
    \PY{l+s+sd}{\PYZdq{}\PYZdq{}\PYZdq{}}
\PY{l+s+sd}{    Implement the backward pass for a RNN over an entire sequence of input data.}

\PY{l+s+sd}{    Arguments:}
\PY{l+s+sd}{    da \PYZhy{}\PYZhy{} Upstream gradients of all hidden states, of shape (n\PYZus{}a, m, T\PYZus{}x)}
\PY{l+s+sd}{    caches \PYZhy{}\PYZhy{} tuple containing information from the forward pass (rnn\PYZus{}forward)}
\PY{l+s+sd}{    }
\PY{l+s+sd}{    Returns:}
\PY{l+s+sd}{    gradients \PYZhy{}\PYZhy{} python dictionary containing:}
\PY{l+s+sd}{                        dx \PYZhy{}\PYZhy{} Gradient w.r.t. the input data, numpy\PYZhy{}array of shape (n\PYZus{}x, m, T\PYZus{}x)}
\PY{l+s+sd}{                        da0 \PYZhy{}\PYZhy{} Gradient w.r.t the initial hidden state, numpy\PYZhy{}array of shape (n\PYZus{}a, m)}
\PY{l+s+sd}{                        dWax \PYZhy{}\PYZhy{} Gradient w.r.t the input\PYZsq{}s weight matrix, numpy\PYZhy{}array of shape (n\PYZus{}a, n\PYZus{}x)}
\PY{l+s+sd}{                        dWaa \PYZhy{}\PYZhy{} Gradient w.r.t the hidden state\PYZsq{}s weight matrix, numpy\PYZhy{}arrayof shape (n\PYZus{}a, n\PYZus{}a)}
\PY{l+s+sd}{                        dba \PYZhy{}\PYZhy{} Gradient w.r.t the bias, of shape (n\PYZus{}a, 1)}
\PY{l+s+sd}{    \PYZdq{}\PYZdq{}\PYZdq{}}
        
    \PY{c+c1}{\PYZsh{}\PYZsh{}\PYZsh{} START CODE HERE \PYZsh{}\PYZsh{}\PYZsh{}}
    
    \PY{c+c1}{\PYZsh{} Retrieve values from the first cache (t=1) of caches (≈2 lines)}
    \PY{p}{(}\PY{n}{caches}\PY{p}{,} \PY{n}{x}\PY{p}{)} \PY{o}{=} \PY{n}{caches}
    \PY{p}{(}\PY{n}{a1}\PY{p}{,} \PY{n}{a0}\PY{p}{,} \PY{n}{x1}\PY{p}{,} \PY{n}{parameters}\PY{p}{)} \PY{o}{=} \PY{n}{caches}\PY{p}{[}\PY{l+m+mi}{0}\PY{p}{]}
    
    \PY{c+c1}{\PYZsh{} Retrieve dimensions from da\PYZsq{}s and x1\PYZsq{}s shapes (≈2 lines)}
    \PY{n}{n\PYZus{}a}\PY{p}{,} \PY{n}{m}\PY{p}{,} \PY{n}{T\PYZus{}x} \PY{o}{=} \PY{n}{da}\PY{o}{.}\PY{n}{shape}
    \PY{n}{n\PYZus{}x}\PY{p}{,} \PY{n}{m} \PY{o}{=} \PY{n}{x1}\PY{o}{.}\PY{n}{shape} 
    
    \PY{c+c1}{\PYZsh{} initialize the gradients with the right sizes (≈6 lines)}
    \PY{n}{dx} \PY{o}{=} \PY{n}{np}\PY{o}{.}\PY{n}{zeros}\PY{p}{(}\PY{p}{(}\PY{n}{n\PYZus{}x}\PY{p}{,} \PY{n}{m}\PY{p}{,} \PY{n}{T\PYZus{}x}\PY{p}{)}\PY{p}{)}
    \PY{n}{dWax} \PY{o}{=} \PY{n}{np}\PY{o}{.}\PY{n}{zeros}\PY{p}{(}\PY{p}{(}\PY{n}{n\PYZus{}a}\PY{p}{,} \PY{n}{n\PYZus{}x}\PY{p}{)}\PY{p}{)}
    \PY{n}{dWaa} \PY{o}{=} \PY{n}{np}\PY{o}{.}\PY{n}{zeros}\PY{p}{(}\PY{p}{(}\PY{n}{n\PYZus{}a}\PY{p}{,} \PY{n}{n\PYZus{}a}\PY{p}{)}\PY{p}{)}
    \PY{n}{dba} \PY{o}{=} \PY{n}{np}\PY{o}{.}\PY{n}{zeros}\PY{p}{(}\PY{p}{(}\PY{n}{n\PYZus{}a}\PY{p}{,} \PY{l+m+mi}{1}\PY{p}{)}\PY{p}{)}
    \PY{n}{da0} \PY{o}{=} \PY{n}{np}\PY{o}{.}\PY{n}{zeros}\PY{p}{(}\PY{p}{(}\PY{n}{n\PYZus{}a}\PY{p}{,} \PY{n}{m}\PY{p}{)}\PY{p}{)}
    \PY{n}{da\PYZus{}prevt} \PY{o}{=} \PY{n}{da}\PY{p}{[}\PY{p}{:}\PY{p}{,}\PY{p}{:}\PY{p}{,}\PY{n}{T\PYZus{}x} \PY{o}{\PYZhy{}} \PY{l+m+mi}{1}\PY{p}{]}
    
    \PY{c+c1}{\PYZsh{} Loop through all the time steps}
    \PY{k}{for} \PY{n}{t} \PY{o+ow}{in} \PY{n+nb}{reversed}\PY{p}{(}\PY{n+nb}{range}\PY{p}{(}\PY{n}{T\PYZus{}x}\PY{p}{)}\PY{p}{)}\PY{p}{:}
        \PY{c+c1}{\PYZsh{} Compute gradients at time step t. Choose wisely the \PYZdq{}da\PYZus{}next\PYZdq{} and the \PYZdq{}cache\PYZdq{} to use in the backward propagation step. (≈1 line)}
        \PY{n}{gradients} \PY{o}{=} \PY{n}{rnn\PYZus{}cell\PYZus{}backward}\PY{p}{(}\PY{n}{da\PYZus{}prevt}\PY{p}{,} \PY{n}{caches}\PY{p}{[}\PY{n}{t} \PY{o}{\PYZhy{}} \PY{l+m+mi}{1}\PY{p}{]}\PY{p}{)}
        \PY{c+c1}{\PYZsh{} Retrieve derivatives from gradients (≈ 1 line)}
        \PY{n}{dxt}\PY{p}{,} \PY{n}{da\PYZus{}prevt}\PY{p}{,} \PY{n}{dWaxt}\PY{p}{,} \PY{n}{dWaat}\PY{p}{,} \PY{n}{dbat} \PY{o}{=} \PY{n}{gradients}\PY{p}{[}\PY{l+s+s2}{\PYZdq{}}\PY{l+s+s2}{dxt}\PY{l+s+s2}{\PYZdq{}}\PY{p}{]}\PY{p}{,} \PY{n}{gradients}\PY{p}{[}\PY{l+s+s2}{\PYZdq{}}\PY{l+s+s2}{da\PYZus{}prev}\PY{l+s+s2}{\PYZdq{}}\PY{p}{]}\PY{p}{,} \PY{n}{gradients}\PY{p}{[}\PY{l+s+s2}{\PYZdq{}}\PY{l+s+s2}{dWax}\PY{l+s+s2}{\PYZdq{}}\PY{p}{]}\PY{p}{,} \PY{n}{gradients}\PY{p}{[}\PY{l+s+s2}{\PYZdq{}}\PY{l+s+s2}{dWaa}\PY{l+s+s2}{\PYZdq{}}\PY{p}{]}\PY{p}{,} \PY{n}{gradients}\PY{p}{[}\PY{l+s+s2}{\PYZdq{}}\PY{l+s+s2}{dba}\PY{l+s+s2}{\PYZdq{}}\PY{p}{]}
        \PY{c+c1}{\PYZsh{} Increment global derivatives w.r.t parameters by adding their derivative at time\PYZhy{}step t (≈4 lines)}
        \PY{n}{dx}\PY{p}{[}\PY{p}{:}\PY{p}{,} \PY{p}{:}\PY{p}{,} \PY{n}{t}\PY{p}{]} \PY{o}{=} \PY{n}{dxt} 
        \PY{n}{dWax} \PY{o}{+}\PY{o}{=} \PY{n}{dWaxt}
        \PY{n}{dWaa} \PY{o}{+}\PY{o}{=} \PY{n}{dWaat} 
        \PY{n}{dba} \PY{o}{+}\PY{o}{=} \PY{n}{dbat}
        
    \PY{c+c1}{\PYZsh{} Set da0 to the gradient of a which has been backpropagated through all time\PYZhy{}steps (≈1 line) }
    \PY{n}{da0} \PY{o}{=} \PY{k+kc}{None}
    \PY{c+c1}{\PYZsh{}\PYZsh{}\PYZsh{} END CODE HERE \PYZsh{}\PYZsh{}\PYZsh{}}

    \PY{c+c1}{\PYZsh{} Store the gradients in a python dictionary}
    \PY{n}{gradients} \PY{o}{=} \PY{p}{\PYZob{}}\PY{l+s+s2}{\PYZdq{}}\PY{l+s+s2}{dx}\PY{l+s+s2}{\PYZdq{}}\PY{p}{:} \PY{n}{dx}\PY{p}{,} \PY{l+s+s2}{\PYZdq{}}\PY{l+s+s2}{da0}\PY{l+s+s2}{\PYZdq{}}\PY{p}{:} \PY{n}{da0}\PY{p}{,} \PY{l+s+s2}{\PYZdq{}}\PY{l+s+s2}{dWax}\PY{l+s+s2}{\PYZdq{}}\PY{p}{:} \PY{n}{dWax}\PY{p}{,} \PY{l+s+s2}{\PYZdq{}}\PY{l+s+s2}{dWaa}\PY{l+s+s2}{\PYZdq{}}\PY{p}{:} \PY{n}{dWaa}\PY{p}{,}\PY{l+s+s2}{\PYZdq{}}\PY{l+s+s2}{dba}\PY{l+s+s2}{\PYZdq{}}\PY{p}{:} \PY{n}{dba}\PY{p}{\PYZcb{}}
    
    \PY{k}{return} \PY{n}{gradients}
\end{Verbatim}
\end{tcolorbox}

    \begin{tcolorbox}[breakable, size=fbox, boxrule=1pt, pad at break*=1mm,colback=cellbackground, colframe=cellborder]
\prompt{In}{incolor}{59}{\boxspacing}
\begin{Verbatim}[commandchars=\\\{\}]
\PY{n}{np}\PY{o}{.}\PY{n}{random}\PY{o}{.}\PY{n}{seed}\PY{p}{(}\PY{l+m+mi}{1}\PY{p}{)}
\PY{n}{x\PYZus{}tmp} \PY{o}{=} \PY{n}{np}\PY{o}{.}\PY{n}{random}\PY{o}{.}\PY{n}{randn}\PY{p}{(}\PY{l+m+mi}{3}\PY{p}{,}\PY{l+m+mi}{10}\PY{p}{,}\PY{l+m+mi}{4}\PY{p}{)}
\PY{n}{a0\PYZus{}tmp} \PY{o}{=} \PY{n}{np}\PY{o}{.}\PY{n}{random}\PY{o}{.}\PY{n}{randn}\PY{p}{(}\PY{l+m+mi}{5}\PY{p}{,}\PY{l+m+mi}{10}\PY{p}{)}
\PY{n}{parameters\PYZus{}tmp} \PY{o}{=} \PY{p}{\PYZob{}}\PY{p}{\PYZcb{}}
\PY{n}{parameters\PYZus{}tmp}\PY{p}{[}\PY{l+s+s1}{\PYZsq{}}\PY{l+s+s1}{Wax}\PY{l+s+s1}{\PYZsq{}}\PY{p}{]} \PY{o}{=} \PY{n}{np}\PY{o}{.}\PY{n}{random}\PY{o}{.}\PY{n}{randn}\PY{p}{(}\PY{l+m+mi}{5}\PY{p}{,}\PY{l+m+mi}{3}\PY{p}{)}
\PY{n}{parameters\PYZus{}tmp}\PY{p}{[}\PY{l+s+s1}{\PYZsq{}}\PY{l+s+s1}{Waa}\PY{l+s+s1}{\PYZsq{}}\PY{p}{]} \PY{o}{=} \PY{n}{np}\PY{o}{.}\PY{n}{random}\PY{o}{.}\PY{n}{randn}\PY{p}{(}\PY{l+m+mi}{5}\PY{p}{,}\PY{l+m+mi}{5}\PY{p}{)}
\PY{n}{parameters\PYZus{}tmp}\PY{p}{[}\PY{l+s+s1}{\PYZsq{}}\PY{l+s+s1}{Wya}\PY{l+s+s1}{\PYZsq{}}\PY{p}{]} \PY{o}{=} \PY{n}{np}\PY{o}{.}\PY{n}{random}\PY{o}{.}\PY{n}{randn}\PY{p}{(}\PY{l+m+mi}{2}\PY{p}{,}\PY{l+m+mi}{5}\PY{p}{)}
\PY{n}{parameters\PYZus{}tmp}\PY{p}{[}\PY{l+s+s1}{\PYZsq{}}\PY{l+s+s1}{ba}\PY{l+s+s1}{\PYZsq{}}\PY{p}{]} \PY{o}{=} \PY{n}{np}\PY{o}{.}\PY{n}{random}\PY{o}{.}\PY{n}{randn}\PY{p}{(}\PY{l+m+mi}{5}\PY{p}{,}\PY{l+m+mi}{1}\PY{p}{)}
\PY{n}{parameters\PYZus{}tmp}\PY{p}{[}\PY{l+s+s1}{\PYZsq{}}\PY{l+s+s1}{by}\PY{l+s+s1}{\PYZsq{}}\PY{p}{]} \PY{o}{=} \PY{n}{np}\PY{o}{.}\PY{n}{random}\PY{o}{.}\PY{n}{randn}\PY{p}{(}\PY{l+m+mi}{2}\PY{p}{,}\PY{l+m+mi}{1}\PY{p}{)}

\PY{n}{a\PYZus{}tmp}\PY{p}{,} \PY{n}{y\PYZus{}tmp}\PY{p}{,} \PY{n}{caches\PYZus{}tmp} \PY{o}{=} \PY{n}{rnn\PYZus{}forward}\PY{p}{(}\PY{n}{x\PYZus{}tmp}\PY{p}{,} \PY{n}{a0\PYZus{}tmp}\PY{p}{,} \PY{n}{parameters\PYZus{}tmp}\PY{p}{)}
\PY{n}{da\PYZus{}tmp} \PY{o}{=} \PY{n}{np}\PY{o}{.}\PY{n}{random}\PY{o}{.}\PY{n}{randn}\PY{p}{(}\PY{l+m+mi}{5}\PY{p}{,} \PY{l+m+mi}{10}\PY{p}{,} \PY{l+m+mi}{4}\PY{p}{)}
\PY{n}{gradients\PYZus{}tmp} \PY{o}{=} \PY{n}{rnn\PYZus{}backward}\PY{p}{(}\PY{n}{da\PYZus{}tmp}\PY{p}{,} \PY{n}{caches\PYZus{}tmp}\PY{p}{)}

\PY{n+nb}{print}\PY{p}{(}\PY{l+s+s2}{\PYZdq{}}\PY{l+s+s2}{gradients[}\PY{l+s+se}{\PYZbs{}\PYZdq{}}\PY{l+s+s2}{dx}\PY{l+s+se}{\PYZbs{}\PYZdq{}}\PY{l+s+s2}{][1][2] =}\PY{l+s+s2}{\PYZdq{}}\PY{p}{,} \PY{n}{gradients\PYZus{}tmp}\PY{p}{[}\PY{l+s+s2}{\PYZdq{}}\PY{l+s+s2}{dx}\PY{l+s+s2}{\PYZdq{}}\PY{p}{]}\PY{p}{[}\PY{l+m+mi}{1}\PY{p}{]}\PY{p}{[}\PY{l+m+mi}{2}\PY{p}{]}\PY{p}{)}
\PY{n+nb}{print}\PY{p}{(}\PY{l+s+s2}{\PYZdq{}}\PY{l+s+s2}{gradients[}\PY{l+s+se}{\PYZbs{}\PYZdq{}}\PY{l+s+s2}{dx}\PY{l+s+se}{\PYZbs{}\PYZdq{}}\PY{l+s+s2}{].shape =}\PY{l+s+s2}{\PYZdq{}}\PY{p}{,} \PY{n}{gradients\PYZus{}tmp}\PY{p}{[}\PY{l+s+s2}{\PYZdq{}}\PY{l+s+s2}{dx}\PY{l+s+s2}{\PYZdq{}}\PY{p}{]}\PY{o}{.}\PY{n}{shape}\PY{p}{)}
\PY{n+nb}{print}\PY{p}{(}\PY{l+s+s2}{\PYZdq{}}\PY{l+s+s2}{gradients[}\PY{l+s+se}{\PYZbs{}\PYZdq{}}\PY{l+s+s2}{da0}\PY{l+s+se}{\PYZbs{}\PYZdq{}}\PY{l+s+s2}{][2][3] =}\PY{l+s+s2}{\PYZdq{}}\PY{p}{,} \PY{n}{gradients\PYZus{}tmp}\PY{p}{[}\PY{l+s+s2}{\PYZdq{}}\PY{l+s+s2}{da0}\PY{l+s+s2}{\PYZdq{}}\PY{p}{]}\PY{p}{[}\PY{l+m+mi}{2}\PY{p}{]}\PY{p}{[}\PY{l+m+mi}{3}\PY{p}{]}\PY{p}{)}
\PY{n+nb}{print}\PY{p}{(}\PY{l+s+s2}{\PYZdq{}}\PY{l+s+s2}{gradients[}\PY{l+s+se}{\PYZbs{}\PYZdq{}}\PY{l+s+s2}{da0}\PY{l+s+se}{\PYZbs{}\PYZdq{}}\PY{l+s+s2}{].shape =}\PY{l+s+s2}{\PYZdq{}}\PY{p}{,} \PY{n}{gradients\PYZus{}tmp}\PY{p}{[}\PY{l+s+s2}{\PYZdq{}}\PY{l+s+s2}{da0}\PY{l+s+s2}{\PYZdq{}}\PY{p}{]}\PY{o}{.}\PY{n}{shape}\PY{p}{)}
\PY{n+nb}{print}\PY{p}{(}\PY{l+s+s2}{\PYZdq{}}\PY{l+s+s2}{gradients[}\PY{l+s+se}{\PYZbs{}\PYZdq{}}\PY{l+s+s2}{dWax}\PY{l+s+se}{\PYZbs{}\PYZdq{}}\PY{l+s+s2}{][3][1] =}\PY{l+s+s2}{\PYZdq{}}\PY{p}{,} \PY{n}{gradients\PYZus{}tmp}\PY{p}{[}\PY{l+s+s2}{\PYZdq{}}\PY{l+s+s2}{dWax}\PY{l+s+s2}{\PYZdq{}}\PY{p}{]}\PY{p}{[}\PY{l+m+mi}{3}\PY{p}{]}\PY{p}{[}\PY{l+m+mi}{1}\PY{p}{]}\PY{p}{)}
\PY{n+nb}{print}\PY{p}{(}\PY{l+s+s2}{\PYZdq{}}\PY{l+s+s2}{gradients[}\PY{l+s+se}{\PYZbs{}\PYZdq{}}\PY{l+s+s2}{dWax}\PY{l+s+se}{\PYZbs{}\PYZdq{}}\PY{l+s+s2}{].shape =}\PY{l+s+s2}{\PYZdq{}}\PY{p}{,} \PY{n}{gradients\PYZus{}tmp}\PY{p}{[}\PY{l+s+s2}{\PYZdq{}}\PY{l+s+s2}{dWax}\PY{l+s+s2}{\PYZdq{}}\PY{p}{]}\PY{o}{.}\PY{n}{shape}\PY{p}{)}
\PY{n+nb}{print}\PY{p}{(}\PY{l+s+s2}{\PYZdq{}}\PY{l+s+s2}{gradients[}\PY{l+s+se}{\PYZbs{}\PYZdq{}}\PY{l+s+s2}{dWaa}\PY{l+s+se}{\PYZbs{}\PYZdq{}}\PY{l+s+s2}{][1][2] =}\PY{l+s+s2}{\PYZdq{}}\PY{p}{,} \PY{n}{gradients\PYZus{}tmp}\PY{p}{[}\PY{l+s+s2}{\PYZdq{}}\PY{l+s+s2}{dWaa}\PY{l+s+s2}{\PYZdq{}}\PY{p}{]}\PY{p}{[}\PY{l+m+mi}{1}\PY{p}{]}\PY{p}{[}\PY{l+m+mi}{2}\PY{p}{]}\PY{p}{)}
\PY{n+nb}{print}\PY{p}{(}\PY{l+s+s2}{\PYZdq{}}\PY{l+s+s2}{gradients[}\PY{l+s+se}{\PYZbs{}\PYZdq{}}\PY{l+s+s2}{dWaa}\PY{l+s+se}{\PYZbs{}\PYZdq{}}\PY{l+s+s2}{].shape =}\PY{l+s+s2}{\PYZdq{}}\PY{p}{,} \PY{n}{gradients\PYZus{}tmp}\PY{p}{[}\PY{l+s+s2}{\PYZdq{}}\PY{l+s+s2}{dWaa}\PY{l+s+s2}{\PYZdq{}}\PY{p}{]}\PY{o}{.}\PY{n}{shape}\PY{p}{)}
\PY{n+nb}{print}\PY{p}{(}\PY{l+s+s2}{\PYZdq{}}\PY{l+s+s2}{gradients[}\PY{l+s+se}{\PYZbs{}\PYZdq{}}\PY{l+s+s2}{dba}\PY{l+s+se}{\PYZbs{}\PYZdq{}}\PY{l+s+s2}{][4] =}\PY{l+s+s2}{\PYZdq{}}\PY{p}{,} \PY{n}{gradients\PYZus{}tmp}\PY{p}{[}\PY{l+s+s2}{\PYZdq{}}\PY{l+s+s2}{dba}\PY{l+s+s2}{\PYZdq{}}\PY{p}{]}\PY{p}{[}\PY{l+m+mi}{4}\PY{p}{]}\PY{p}{)}
\PY{n+nb}{print}\PY{p}{(}\PY{l+s+s2}{\PYZdq{}}\PY{l+s+s2}{gradients[}\PY{l+s+se}{\PYZbs{}\PYZdq{}}\PY{l+s+s2}{dba}\PY{l+s+se}{\PYZbs{}\PYZdq{}}\PY{l+s+s2}{].shape =}\PY{l+s+s2}{\PYZdq{}}\PY{p}{,} \PY{n}{gradients\PYZus{}tmp}\PY{p}{[}\PY{l+s+s2}{\PYZdq{}}\PY{l+s+s2}{dba}\PY{l+s+s2}{\PYZdq{}}\PY{p}{]}\PY{o}{.}\PY{n}{shape}\PY{p}{)}
\end{Verbatim}
\end{tcolorbox}

    \begin{Verbatim}[commandchars=\\\{\}]
gradients["dx"][1][2] = [-0.00268429  0.28508385  0.11782283  0.03309917]
gradients["dx"].shape = (3, 10, 4)
    \end{Verbatim}

    \begin{Verbatim}[commandchars=\\\{\}]

        ---------------------------------------------------------------------------

        TypeError                                 Traceback (most recent call last)

        <ipython-input-59-9af401171061> in <module>
         15 print("gradients[\textbackslash{}"dx\textbackslash{}"][1][2] =", gradients\_tmp["dx"][1][2])
         16 print("gradients[\textbackslash{}"dx\textbackslash{}"].shape =", gradients\_tmp["dx"].shape)
    ---> 17 print("gradients[\textbackslash{}"da0\textbackslash{}"][2][3] =", gradients\_tmp["da0"][2][3])
         18 print("gradients[\textbackslash{}"da0\textbackslash{}"].shape =", gradients\_tmp["da0"].shape)
         19 print("gradients[\textbackslash{}"dWax\textbackslash{}"][3][1] =", gradients\_tmp["dWax"][3][1])


        TypeError: 'NoneType' object is not subscriptable

    \end{Verbatim}

    \textbf{Expected Output}:

gradients{[}``dx''{]}{[}1{]}{[}2{]} =

{[}-2.07101689 -0.59255627 0.02466855 0.01483317{]}

\begin{verbatim}
<tr>
    <td>
        <b>gradients["dx"].shape</b> =
    </td>
    <td>
       (3, 10, 4)
    </td>
</tr>
    <tr>
    <td>
        <b>gradients["da0"][2][3]</b> =
    </td>
    <td>
       -0.314942375127
    </td>
</tr>
    <tr>
    <td>
        <b>gradients["da0"].shape</b> =
    </td>
    <td>
       (5, 10)
    </td>
</tr>
     <tr>
    <td>
        <b>gradients["dWax"][3][1]</b> =
    </td>
    <td>
       11.2641044965
    </td>
</tr>
    <tr>
    <td>
        <b>gradients["dWax"].shape</b> =
    </td>
    <td>
       (5, 3)
    </td>
</tr>
    <tr>
    <td>
        <b>gradients["dWaa"][1][2]</b> = 
    </td>
    <td>
       2.30333312658
    </td>
</tr>
    <tr>
    <td>
        <b>gradients["dWaa"].shape</b> =
    </td>
    <td>
       (5, 5)
    </td>
</tr>
    <tr>
    <td>
        <b>gradients["dba"][4]</b> = 
    </td>
    <td>
       [-0.74747722]
    </td>
</tr>
    <tr>
    <td>
        <b>gradients["dba"].shape</b> = 
    </td>
    <td>
       (5, 1)
    </td>
</tr>
\end{verbatim}

    \#\#\# 3.2 - LSTM Backward Pass

    \hypertarget{one-step-backward}{%
\paragraph{1. One Step Backward}\label{one-step-backward}}

The LSTM backward pass is slightly more complicated than the forward
pass.

Figure 8: LSTM Cell Backward. Note the output functions, while part of
the \texttt{lstm\_cell\_forward}, are not included in
\texttt{lstm\_cell\_backward}

The equations for the LSTM backward pass are provided below. (If you
enjoy calculus exercises feel free to try deriving these from scratch
yourself.)

    \hypertarget{gate-derivatives}{%
\paragraph{2. Gate Derivatives}\label{gate-derivatives}}

Note the location of the gate derivatives (\(\gamma\)..) between the
dense layer and the activation function (see graphic above). This is
convenient for computing parameter derivatives in the next step.
\begin{align}
d\gamma_o^{\langle t \rangle} &= da_{next}*\tanh(c_{next}) * \Gamma_o^{\langle t \rangle}*\left(1-\Gamma_o^{\langle t \rangle}\right)\tag{7} \\[8pt]
dp\widetilde{c}^{\langle t \rangle} &= \left(dc_{next}*\Gamma_u^{\langle t \rangle}+ \Gamma_o^{\langle t \rangle}* (1-\tanh^2(c_{next})) * \Gamma_u^{\langle t \rangle} * da_{next} \right) * \left(1-\left(\widetilde c^{\langle t \rangle}\right)^2\right) \tag{8} \\[8pt]
d\gamma_u^{\langle t \rangle} &= \left(dc_{next}*\widetilde{c}^{\langle t \rangle} + \Gamma_o^{\langle t \rangle}* (1-\tanh^2(c_{next})) * \widetilde{c}^{\langle t \rangle} * da_{next}\right)*\Gamma_u^{\langle t \rangle}*\left(1-\Gamma_u^{\langle t \rangle}\right)\tag{9} \\[8pt]
d\gamma_f^{\langle t \rangle} &= \left(dc_{next}* c_{prev} + \Gamma_o^{\langle t \rangle} * (1-\tanh^2(c_{next})) * c_{prev} * da_{next}\right)*\Gamma_f^{\langle t \rangle}*\left(1-\Gamma_f^{\langle t \rangle}\right)\tag{10}
\end{align}

\hypertarget{parameter-derivatives}{%
\paragraph{3. Parameter Derivatives}\label{parameter-derivatives}}

\$ dW\_f = d\gamma\_f\^{}\{\langle t \rangle\}

\begin{bmatrix} a_{prev} \\ x_t\end{bmatrix}

\^{}T \tag{11} \$ \$ dW\_u = d\gamma\_u\^{}\{\langle t \rangle\}

\begin{bmatrix} a_{prev} \\ x_t\end{bmatrix}

\^{}T \tag{12} \$ \$ dW\_c = dp\widetilde c\^{}\{\langle t \rangle\}

\begin{bmatrix} a_{prev} \\ x_t\end{bmatrix}

\^{}T \tag{13} \$ \$ dW\_o = d\gamma\_o\^{}\{\langle t \rangle\}

\begin{bmatrix} a_{prev} \\ x_t\end{bmatrix}

\^{}T \tag{14}\$

To calculate \(db_f, db_u, db_c, db_o\) you just need to sum across all
`m' examples (axis= 1) on
\(d\gamma_f^{\langle t \rangle}, d\gamma_u^{\langle t \rangle}, dp\widetilde c^{\langle t \rangle}, d\gamma_o^{\langle t \rangle}\)
respectively. Note that you should have the \texttt{keepdims\ =\ True}
option.

\(\displaystyle db_f = \sum_{batch}d\gamma_f^{\langle t \rangle}\tag{15}\)
\(\displaystyle db_u = \sum_{batch}d\gamma_u^{\langle t \rangle}\tag{16}\)
\(\displaystyle db_c = \sum_{batch}dp\widetilde c^{\langle t \rangle}\tag{17}\)
\(\displaystyle db_o = \sum_{batch}d\gamma_o^{\langle t \rangle}\tag{18}\)

Finally, you will compute the derivative with respect to the previous
hidden state, previous memory state, and input.

\$ da\_\{prev\} = W\_f\^{}T d\gamma\_f\^{}\{\langle t \rangle\} +
W\_u\^{}T d\gamma\_u\^{}\{\langle t \rangle\}+ W\_c\^{}T
dp\widetilde c\^{}\{\langle t \rangle\} + W\_o\^{}T
d\gamma\_o\^{}\{\langle t \rangle\} \tag{19}\$

Here, to account for concatenation, the weights for equations 19 are the
first n\_a, (i.e.~\(W_f = W_f[:,:n_a]\) etc\ldots)

\$ dc\_\{prev\} = dc\_\{next\}*\Gamma\_f\^{}\{\langle t \rangle\} +
\Gamma\emph{o\^{}\{\langle t \rangle\} * (1-
\tanh\^{}2(c}\{next\}))*\Gamma\_f\^{}\{\langle t \rangle\}*da\_\{next\}
\tag{20}\$

\$ dx\^{}\{\langle t \rangle\} = W\_f\^{}T d\gamma\_f\^{}\{\langle t
\rangle\} + W\_u\^{}T d\gamma\_u\^{}\{\langle t \rangle\}+ W\_c\^{}T
dp\widetilde c\^{}\{\langle t \rangle\} + W\_o\^{}T
d\gamma\_o\^{}\{\langle t \rangle\}\tag{21} \$

where the weights for equation 21 are from n\_a to the end,
(i.e.~\(W_f = W_f[:,n_a:]\) etc\ldots)

\#\#\# Exercise 7 - lstm\_cell\_backward

Implement \texttt{lstm\_cell\_backward} by implementing equations
\(7-21\) below.

\textbf{Note}:

In the code:

\(d\gamma_o^{\langle t \rangle}\) is represented by \texttt{dot},\\
\(dp\widetilde{c}^{\langle t \rangle}\) is represented by
\texttt{dcct},\\
\(d\gamma_u^{\langle t \rangle}\) is represented by \texttt{dit},\\
\(d\gamma_f^{\langle t \rangle}\) is represented by \texttt{dft}

    \begin{tcolorbox}[breakable, size=fbox, boxrule=1pt, pad at break*=1mm,colback=cellbackground, colframe=cellborder]
\prompt{In}{incolor}{ }{\boxspacing}
\begin{Verbatim}[commandchars=\\\{\}]
\PY{c+c1}{\PYZsh{} UNGRADED FUNCTION: lstm\PYZus{}cell\PYZus{}backward}

\PY{k}{def} \PY{n+nf}{lstm\PYZus{}cell\PYZus{}backward}\PY{p}{(}\PY{n}{da\PYZus{}next}\PY{p}{,} \PY{n}{dc\PYZus{}next}\PY{p}{,} \PY{n}{cache}\PY{p}{)}\PY{p}{:}
    \PY{l+s+sd}{\PYZdq{}\PYZdq{}\PYZdq{}}
\PY{l+s+sd}{    Implement the backward pass for the LSTM\PYZhy{}cell (single time\PYZhy{}step).}

\PY{l+s+sd}{    Arguments:}
\PY{l+s+sd}{    da\PYZus{}next \PYZhy{}\PYZhy{} Gradients of next hidden state, of shape (n\PYZus{}a, m)}
\PY{l+s+sd}{    dc\PYZus{}next \PYZhy{}\PYZhy{} Gradients of next cell state, of shape (n\PYZus{}a, m)}
\PY{l+s+sd}{    cache \PYZhy{}\PYZhy{} cache storing information from the forward pass}

\PY{l+s+sd}{    Returns:}
\PY{l+s+sd}{    gradients \PYZhy{}\PYZhy{} python dictionary containing:}
\PY{l+s+sd}{                        dxt \PYZhy{}\PYZhy{} Gradient of input data at time\PYZhy{}step t, of shape (n\PYZus{}x, m)}
\PY{l+s+sd}{                        da\PYZus{}prev \PYZhy{}\PYZhy{} Gradient w.r.t. the previous hidden state, numpy array of shape (n\PYZus{}a, m)}
\PY{l+s+sd}{                        dc\PYZus{}prev \PYZhy{}\PYZhy{} Gradient w.r.t. the previous memory state, of shape (n\PYZus{}a, m, T\PYZus{}x)}
\PY{l+s+sd}{                        dWf \PYZhy{}\PYZhy{} Gradient w.r.t. the weight matrix of the forget gate, numpy array of shape (n\PYZus{}a, n\PYZus{}a + n\PYZus{}x)}
\PY{l+s+sd}{                        dWi \PYZhy{}\PYZhy{} Gradient w.r.t. the weight matrix of the update gate, numpy array of shape (n\PYZus{}a, n\PYZus{}a + n\PYZus{}x)}
\PY{l+s+sd}{                        dWc \PYZhy{}\PYZhy{} Gradient w.r.t. the weight matrix of the memory gate, numpy array of shape (n\PYZus{}a, n\PYZus{}a + n\PYZus{}x)}
\PY{l+s+sd}{                        dWo \PYZhy{}\PYZhy{} Gradient w.r.t. the weight matrix of the output gate, numpy array of shape (n\PYZus{}a, n\PYZus{}a + n\PYZus{}x)}
\PY{l+s+sd}{                        dbf \PYZhy{}\PYZhy{} Gradient w.r.t. biases of the forget gate, of shape (n\PYZus{}a, 1)}
\PY{l+s+sd}{                        dbi \PYZhy{}\PYZhy{} Gradient w.r.t. biases of the update gate, of shape (n\PYZus{}a, 1)}
\PY{l+s+sd}{                        dbc \PYZhy{}\PYZhy{} Gradient w.r.t. biases of the memory gate, of shape (n\PYZus{}a, 1)}
\PY{l+s+sd}{                        dbo \PYZhy{}\PYZhy{} Gradient w.r.t. biases of the output gate, of shape (n\PYZus{}a, 1)}
\PY{l+s+sd}{    \PYZdq{}\PYZdq{}\PYZdq{}}

    \PY{c+c1}{\PYZsh{} Retrieve information from \PYZdq{}cache\PYZdq{}}
    \PY{p}{(}\PY{n}{a\PYZus{}next}\PY{p}{,} \PY{n}{c\PYZus{}next}\PY{p}{,} \PY{n}{a\PYZus{}prev}\PY{p}{,} \PY{n}{c\PYZus{}prev}\PY{p}{,} \PY{n}{ft}\PY{p}{,} \PY{n}{it}\PY{p}{,} \PY{n}{cct}\PY{p}{,} \PY{n}{ot}\PY{p}{,} \PY{n}{xt}\PY{p}{,} \PY{n}{parameters}\PY{p}{)} \PY{o}{=} \PY{n}{cache}
    
    \PY{c+c1}{\PYZsh{}\PYZsh{}\PYZsh{} START CODE HERE \PYZsh{}\PYZsh{}\PYZsh{}}
    \PY{c+c1}{\PYZsh{} Retrieve dimensions from xt\PYZsq{}s and a\PYZus{}next\PYZsq{}s shape (≈2 lines)}
    \PY{n}{n\PYZus{}x}\PY{p}{,} \PY{n}{m} \PY{o}{=} \PY{k+kc}{None}
    \PY{n}{n\PYZus{}a}\PY{p}{,} \PY{n}{m} \PY{o}{=} \PY{k+kc}{None}
    
    \PY{c+c1}{\PYZsh{} Compute gates related derivatives. Their values can be found by looking carefully at equations (7) to (10) (≈4 lines)}
    \PY{n}{dot} \PY{o}{=} \PY{k+kc}{None}
    \PY{n}{dcct} \PY{o}{=} \PY{k+kc}{None}
    \PY{n}{dit} \PY{o}{=} \PY{k+kc}{None}
    \PY{n}{dft} \PY{o}{=} \PY{k+kc}{None}
    
    \PY{c+c1}{\PYZsh{} Compute parameters related derivatives. Use equations (11)\PYZhy{}(18) (≈8 lines)}
    \PY{n}{dWf} \PY{o}{=} \PY{k+kc}{None}
    \PY{n}{dWi} \PY{o}{=} \PY{k+kc}{None}
    \PY{n}{dWc} \PY{o}{=} \PY{k+kc}{None}
    \PY{n}{dWo} \PY{o}{=} \PY{k+kc}{None}
    \PY{n}{dbf} \PY{o}{=} \PY{k+kc}{None}
    \PY{n}{dbi} \PY{o}{=} \PY{k+kc}{None}
    \PY{n}{dbc} \PY{o}{=} \PY{k+kc}{None}
    \PY{n}{dbo} \PY{o}{=} \PY{k+kc}{None}

    \PY{c+c1}{\PYZsh{} Compute derivatives w.r.t previous hidden state, previous memory state and input. Use equations (19)\PYZhy{}(21). (≈3 lines)}
    \PY{n}{da\PYZus{}prev} \PY{o}{=} \PY{k+kc}{None}
    \PY{n}{dc\PYZus{}prev} \PY{o}{=} \PY{k+kc}{None}
    \PY{n}{dxt} \PY{o}{=} \PY{k+kc}{None}
    \PY{c+c1}{\PYZsh{}\PYZsh{}\PYZsh{} END CODE HERE \PYZsh{}\PYZsh{}\PYZsh{}}
    
    
    
    \PY{c+c1}{\PYZsh{} Save gradients in dictionary}
    \PY{n}{gradients} \PY{o}{=} \PY{p}{\PYZob{}}\PY{l+s+s2}{\PYZdq{}}\PY{l+s+s2}{dxt}\PY{l+s+s2}{\PYZdq{}}\PY{p}{:} \PY{n}{dxt}\PY{p}{,} \PY{l+s+s2}{\PYZdq{}}\PY{l+s+s2}{da\PYZus{}prev}\PY{l+s+s2}{\PYZdq{}}\PY{p}{:} \PY{n}{da\PYZus{}prev}\PY{p}{,} \PY{l+s+s2}{\PYZdq{}}\PY{l+s+s2}{dc\PYZus{}prev}\PY{l+s+s2}{\PYZdq{}}\PY{p}{:} \PY{n}{dc\PYZus{}prev}\PY{p}{,} \PY{l+s+s2}{\PYZdq{}}\PY{l+s+s2}{dWf}\PY{l+s+s2}{\PYZdq{}}\PY{p}{:} \PY{n}{dWf}\PY{p}{,}\PY{l+s+s2}{\PYZdq{}}\PY{l+s+s2}{dbf}\PY{l+s+s2}{\PYZdq{}}\PY{p}{:} \PY{n}{dbf}\PY{p}{,} \PY{l+s+s2}{\PYZdq{}}\PY{l+s+s2}{dWi}\PY{l+s+s2}{\PYZdq{}}\PY{p}{:} \PY{n}{dWi}\PY{p}{,}\PY{l+s+s2}{\PYZdq{}}\PY{l+s+s2}{dbi}\PY{l+s+s2}{\PYZdq{}}\PY{p}{:} \PY{n}{dbi}\PY{p}{,}
                \PY{l+s+s2}{\PYZdq{}}\PY{l+s+s2}{dWc}\PY{l+s+s2}{\PYZdq{}}\PY{p}{:} \PY{n}{dWc}\PY{p}{,}\PY{l+s+s2}{\PYZdq{}}\PY{l+s+s2}{dbc}\PY{l+s+s2}{\PYZdq{}}\PY{p}{:} \PY{n}{dbc}\PY{p}{,} \PY{l+s+s2}{\PYZdq{}}\PY{l+s+s2}{dWo}\PY{l+s+s2}{\PYZdq{}}\PY{p}{:} \PY{n}{dWo}\PY{p}{,}\PY{l+s+s2}{\PYZdq{}}\PY{l+s+s2}{dbo}\PY{l+s+s2}{\PYZdq{}}\PY{p}{:} \PY{n}{dbo}\PY{p}{\PYZcb{}}

    \PY{k}{return} \PY{n}{gradients}
\end{Verbatim}
\end{tcolorbox}

    \begin{tcolorbox}[breakable, size=fbox, boxrule=1pt, pad at break*=1mm,colback=cellbackground, colframe=cellborder]
\prompt{In}{incolor}{ }{\boxspacing}
\begin{Verbatim}[commandchars=\\\{\}]
\PY{n}{np}\PY{o}{.}\PY{n}{random}\PY{o}{.}\PY{n}{seed}\PY{p}{(}\PY{l+m+mi}{1}\PY{p}{)}
\PY{n}{xt\PYZus{}tmp} \PY{o}{=} \PY{n}{np}\PY{o}{.}\PY{n}{random}\PY{o}{.}\PY{n}{randn}\PY{p}{(}\PY{l+m+mi}{3}\PY{p}{,}\PY{l+m+mi}{10}\PY{p}{)}
\PY{n}{a\PYZus{}prev\PYZus{}tmp} \PY{o}{=} \PY{n}{np}\PY{o}{.}\PY{n}{random}\PY{o}{.}\PY{n}{randn}\PY{p}{(}\PY{l+m+mi}{5}\PY{p}{,}\PY{l+m+mi}{10}\PY{p}{)}
\PY{n}{c\PYZus{}prev\PYZus{}tmp} \PY{o}{=} \PY{n}{np}\PY{o}{.}\PY{n}{random}\PY{o}{.}\PY{n}{randn}\PY{p}{(}\PY{l+m+mi}{5}\PY{p}{,}\PY{l+m+mi}{10}\PY{p}{)}
\PY{n}{parameters\PYZus{}tmp} \PY{o}{=} \PY{p}{\PYZob{}}\PY{p}{\PYZcb{}}
\PY{n}{parameters\PYZus{}tmp}\PY{p}{[}\PY{l+s+s1}{\PYZsq{}}\PY{l+s+s1}{Wf}\PY{l+s+s1}{\PYZsq{}}\PY{p}{]} \PY{o}{=} \PY{n}{np}\PY{o}{.}\PY{n}{random}\PY{o}{.}\PY{n}{randn}\PY{p}{(}\PY{l+m+mi}{5}\PY{p}{,} \PY{l+m+mi}{5}\PY{o}{+}\PY{l+m+mi}{3}\PY{p}{)}
\PY{n}{parameters\PYZus{}tmp}\PY{p}{[}\PY{l+s+s1}{\PYZsq{}}\PY{l+s+s1}{bf}\PY{l+s+s1}{\PYZsq{}}\PY{p}{]} \PY{o}{=} \PY{n}{np}\PY{o}{.}\PY{n}{random}\PY{o}{.}\PY{n}{randn}\PY{p}{(}\PY{l+m+mi}{5}\PY{p}{,}\PY{l+m+mi}{1}\PY{p}{)}
\PY{n}{parameters\PYZus{}tmp}\PY{p}{[}\PY{l+s+s1}{\PYZsq{}}\PY{l+s+s1}{Wi}\PY{l+s+s1}{\PYZsq{}}\PY{p}{]} \PY{o}{=} \PY{n}{np}\PY{o}{.}\PY{n}{random}\PY{o}{.}\PY{n}{randn}\PY{p}{(}\PY{l+m+mi}{5}\PY{p}{,} \PY{l+m+mi}{5}\PY{o}{+}\PY{l+m+mi}{3}\PY{p}{)}
\PY{n}{parameters\PYZus{}tmp}\PY{p}{[}\PY{l+s+s1}{\PYZsq{}}\PY{l+s+s1}{bi}\PY{l+s+s1}{\PYZsq{}}\PY{p}{]} \PY{o}{=} \PY{n}{np}\PY{o}{.}\PY{n}{random}\PY{o}{.}\PY{n}{randn}\PY{p}{(}\PY{l+m+mi}{5}\PY{p}{,}\PY{l+m+mi}{1}\PY{p}{)}
\PY{n}{parameters\PYZus{}tmp}\PY{p}{[}\PY{l+s+s1}{\PYZsq{}}\PY{l+s+s1}{Wo}\PY{l+s+s1}{\PYZsq{}}\PY{p}{]} \PY{o}{=} \PY{n}{np}\PY{o}{.}\PY{n}{random}\PY{o}{.}\PY{n}{randn}\PY{p}{(}\PY{l+m+mi}{5}\PY{p}{,} \PY{l+m+mi}{5}\PY{o}{+}\PY{l+m+mi}{3}\PY{p}{)}
\PY{n}{parameters\PYZus{}tmp}\PY{p}{[}\PY{l+s+s1}{\PYZsq{}}\PY{l+s+s1}{bo}\PY{l+s+s1}{\PYZsq{}}\PY{p}{]} \PY{o}{=} \PY{n}{np}\PY{o}{.}\PY{n}{random}\PY{o}{.}\PY{n}{randn}\PY{p}{(}\PY{l+m+mi}{5}\PY{p}{,}\PY{l+m+mi}{1}\PY{p}{)}
\PY{n}{parameters\PYZus{}tmp}\PY{p}{[}\PY{l+s+s1}{\PYZsq{}}\PY{l+s+s1}{Wc}\PY{l+s+s1}{\PYZsq{}}\PY{p}{]} \PY{o}{=} \PY{n}{np}\PY{o}{.}\PY{n}{random}\PY{o}{.}\PY{n}{randn}\PY{p}{(}\PY{l+m+mi}{5}\PY{p}{,} \PY{l+m+mi}{5}\PY{o}{+}\PY{l+m+mi}{3}\PY{p}{)}
\PY{n}{parameters\PYZus{}tmp}\PY{p}{[}\PY{l+s+s1}{\PYZsq{}}\PY{l+s+s1}{bc}\PY{l+s+s1}{\PYZsq{}}\PY{p}{]} \PY{o}{=} \PY{n}{np}\PY{o}{.}\PY{n}{random}\PY{o}{.}\PY{n}{randn}\PY{p}{(}\PY{l+m+mi}{5}\PY{p}{,}\PY{l+m+mi}{1}\PY{p}{)}
\PY{n}{parameters\PYZus{}tmp}\PY{p}{[}\PY{l+s+s1}{\PYZsq{}}\PY{l+s+s1}{Wy}\PY{l+s+s1}{\PYZsq{}}\PY{p}{]} \PY{o}{=} \PY{n}{np}\PY{o}{.}\PY{n}{random}\PY{o}{.}\PY{n}{randn}\PY{p}{(}\PY{l+m+mi}{2}\PY{p}{,}\PY{l+m+mi}{5}\PY{p}{)}
\PY{n}{parameters\PYZus{}tmp}\PY{p}{[}\PY{l+s+s1}{\PYZsq{}}\PY{l+s+s1}{by}\PY{l+s+s1}{\PYZsq{}}\PY{p}{]} \PY{o}{=} \PY{n}{np}\PY{o}{.}\PY{n}{random}\PY{o}{.}\PY{n}{randn}\PY{p}{(}\PY{l+m+mi}{2}\PY{p}{,}\PY{l+m+mi}{1}\PY{p}{)}

\PY{n}{a\PYZus{}next\PYZus{}tmp}\PY{p}{,} \PY{n}{c\PYZus{}next\PYZus{}tmp}\PY{p}{,} \PY{n}{yt\PYZus{}tmp}\PY{p}{,} \PY{n}{cache\PYZus{}tmp} \PY{o}{=} \PY{n}{lstm\PYZus{}cell\PYZus{}forward}\PY{p}{(}\PY{n}{xt\PYZus{}tmp}\PY{p}{,} \PY{n}{a\PYZus{}prev\PYZus{}tmp}\PY{p}{,} \PY{n}{c\PYZus{}prev\PYZus{}tmp}\PY{p}{,} \PY{n}{parameters\PYZus{}tmp}\PY{p}{)}

\PY{n}{da\PYZus{}next\PYZus{}tmp} \PY{o}{=} \PY{n}{np}\PY{o}{.}\PY{n}{random}\PY{o}{.}\PY{n}{randn}\PY{p}{(}\PY{l+m+mi}{5}\PY{p}{,}\PY{l+m+mi}{10}\PY{p}{)}
\PY{n}{dc\PYZus{}next\PYZus{}tmp} \PY{o}{=} \PY{n}{np}\PY{o}{.}\PY{n}{random}\PY{o}{.}\PY{n}{randn}\PY{p}{(}\PY{l+m+mi}{5}\PY{p}{,}\PY{l+m+mi}{10}\PY{p}{)}
\PY{n}{gradients\PYZus{}tmp} \PY{o}{=} \PY{n}{lstm\PYZus{}cell\PYZus{}backward}\PY{p}{(}\PY{n}{da\PYZus{}next\PYZus{}tmp}\PY{p}{,} \PY{n}{dc\PYZus{}next\PYZus{}tmp}\PY{p}{,} \PY{n}{cache\PYZus{}tmp}\PY{p}{)}
\PY{n+nb}{print}\PY{p}{(}\PY{l+s+s2}{\PYZdq{}}\PY{l+s+s2}{gradients[}\PY{l+s+se}{\PYZbs{}\PYZdq{}}\PY{l+s+s2}{dxt}\PY{l+s+se}{\PYZbs{}\PYZdq{}}\PY{l+s+s2}{][1][2] =}\PY{l+s+s2}{\PYZdq{}}\PY{p}{,} \PY{n}{gradients\PYZus{}tmp}\PY{p}{[}\PY{l+s+s2}{\PYZdq{}}\PY{l+s+s2}{dxt}\PY{l+s+s2}{\PYZdq{}}\PY{p}{]}\PY{p}{[}\PY{l+m+mi}{1}\PY{p}{]}\PY{p}{[}\PY{l+m+mi}{2}\PY{p}{]}\PY{p}{)}
\PY{n+nb}{print}\PY{p}{(}\PY{l+s+s2}{\PYZdq{}}\PY{l+s+s2}{gradients[}\PY{l+s+se}{\PYZbs{}\PYZdq{}}\PY{l+s+s2}{dxt}\PY{l+s+se}{\PYZbs{}\PYZdq{}}\PY{l+s+s2}{].shape =}\PY{l+s+s2}{\PYZdq{}}\PY{p}{,} \PY{n}{gradients\PYZus{}tmp}\PY{p}{[}\PY{l+s+s2}{\PYZdq{}}\PY{l+s+s2}{dxt}\PY{l+s+s2}{\PYZdq{}}\PY{p}{]}\PY{o}{.}\PY{n}{shape}\PY{p}{)}
\PY{n+nb}{print}\PY{p}{(}\PY{l+s+s2}{\PYZdq{}}\PY{l+s+s2}{gradients[}\PY{l+s+se}{\PYZbs{}\PYZdq{}}\PY{l+s+s2}{da\PYZus{}prev}\PY{l+s+se}{\PYZbs{}\PYZdq{}}\PY{l+s+s2}{][2][3] =}\PY{l+s+s2}{\PYZdq{}}\PY{p}{,} \PY{n}{gradients\PYZus{}tmp}\PY{p}{[}\PY{l+s+s2}{\PYZdq{}}\PY{l+s+s2}{da\PYZus{}prev}\PY{l+s+s2}{\PYZdq{}}\PY{p}{]}\PY{p}{[}\PY{l+m+mi}{2}\PY{p}{]}\PY{p}{[}\PY{l+m+mi}{3}\PY{p}{]}\PY{p}{)}
\PY{n+nb}{print}\PY{p}{(}\PY{l+s+s2}{\PYZdq{}}\PY{l+s+s2}{gradients[}\PY{l+s+se}{\PYZbs{}\PYZdq{}}\PY{l+s+s2}{da\PYZus{}prev}\PY{l+s+se}{\PYZbs{}\PYZdq{}}\PY{l+s+s2}{].shape =}\PY{l+s+s2}{\PYZdq{}}\PY{p}{,} \PY{n}{gradients\PYZus{}tmp}\PY{p}{[}\PY{l+s+s2}{\PYZdq{}}\PY{l+s+s2}{da\PYZus{}prev}\PY{l+s+s2}{\PYZdq{}}\PY{p}{]}\PY{o}{.}\PY{n}{shape}\PY{p}{)}
\PY{n+nb}{print}\PY{p}{(}\PY{l+s+s2}{\PYZdq{}}\PY{l+s+s2}{gradients[}\PY{l+s+se}{\PYZbs{}\PYZdq{}}\PY{l+s+s2}{dc\PYZus{}prev}\PY{l+s+se}{\PYZbs{}\PYZdq{}}\PY{l+s+s2}{][2][3] =}\PY{l+s+s2}{\PYZdq{}}\PY{p}{,} \PY{n}{gradients\PYZus{}tmp}\PY{p}{[}\PY{l+s+s2}{\PYZdq{}}\PY{l+s+s2}{dc\PYZus{}prev}\PY{l+s+s2}{\PYZdq{}}\PY{p}{]}\PY{p}{[}\PY{l+m+mi}{2}\PY{p}{]}\PY{p}{[}\PY{l+m+mi}{3}\PY{p}{]}\PY{p}{)}
\PY{n+nb}{print}\PY{p}{(}\PY{l+s+s2}{\PYZdq{}}\PY{l+s+s2}{gradients[}\PY{l+s+se}{\PYZbs{}\PYZdq{}}\PY{l+s+s2}{dc\PYZus{}prev}\PY{l+s+se}{\PYZbs{}\PYZdq{}}\PY{l+s+s2}{].shape =}\PY{l+s+s2}{\PYZdq{}}\PY{p}{,} \PY{n}{gradients\PYZus{}tmp}\PY{p}{[}\PY{l+s+s2}{\PYZdq{}}\PY{l+s+s2}{dc\PYZus{}prev}\PY{l+s+s2}{\PYZdq{}}\PY{p}{]}\PY{o}{.}\PY{n}{shape}\PY{p}{)}
\PY{n+nb}{print}\PY{p}{(}\PY{l+s+s2}{\PYZdq{}}\PY{l+s+s2}{gradients[}\PY{l+s+se}{\PYZbs{}\PYZdq{}}\PY{l+s+s2}{dWf}\PY{l+s+se}{\PYZbs{}\PYZdq{}}\PY{l+s+s2}{][3][1] =}\PY{l+s+s2}{\PYZdq{}}\PY{p}{,} \PY{n}{gradients\PYZus{}tmp}\PY{p}{[}\PY{l+s+s2}{\PYZdq{}}\PY{l+s+s2}{dWf}\PY{l+s+s2}{\PYZdq{}}\PY{p}{]}\PY{p}{[}\PY{l+m+mi}{3}\PY{p}{]}\PY{p}{[}\PY{l+m+mi}{1}\PY{p}{]}\PY{p}{)}
\PY{n+nb}{print}\PY{p}{(}\PY{l+s+s2}{\PYZdq{}}\PY{l+s+s2}{gradients[}\PY{l+s+se}{\PYZbs{}\PYZdq{}}\PY{l+s+s2}{dWf}\PY{l+s+se}{\PYZbs{}\PYZdq{}}\PY{l+s+s2}{].shape =}\PY{l+s+s2}{\PYZdq{}}\PY{p}{,} \PY{n}{gradients\PYZus{}tmp}\PY{p}{[}\PY{l+s+s2}{\PYZdq{}}\PY{l+s+s2}{dWf}\PY{l+s+s2}{\PYZdq{}}\PY{p}{]}\PY{o}{.}\PY{n}{shape}\PY{p}{)}
\PY{n+nb}{print}\PY{p}{(}\PY{l+s+s2}{\PYZdq{}}\PY{l+s+s2}{gradients[}\PY{l+s+se}{\PYZbs{}\PYZdq{}}\PY{l+s+s2}{dWi}\PY{l+s+se}{\PYZbs{}\PYZdq{}}\PY{l+s+s2}{][1][2] =}\PY{l+s+s2}{\PYZdq{}}\PY{p}{,} \PY{n}{gradients\PYZus{}tmp}\PY{p}{[}\PY{l+s+s2}{\PYZdq{}}\PY{l+s+s2}{dWi}\PY{l+s+s2}{\PYZdq{}}\PY{p}{]}\PY{p}{[}\PY{l+m+mi}{1}\PY{p}{]}\PY{p}{[}\PY{l+m+mi}{2}\PY{p}{]}\PY{p}{)}
\PY{n+nb}{print}\PY{p}{(}\PY{l+s+s2}{\PYZdq{}}\PY{l+s+s2}{gradients[}\PY{l+s+se}{\PYZbs{}\PYZdq{}}\PY{l+s+s2}{dWi}\PY{l+s+se}{\PYZbs{}\PYZdq{}}\PY{l+s+s2}{].shape =}\PY{l+s+s2}{\PYZdq{}}\PY{p}{,} \PY{n}{gradients\PYZus{}tmp}\PY{p}{[}\PY{l+s+s2}{\PYZdq{}}\PY{l+s+s2}{dWi}\PY{l+s+s2}{\PYZdq{}}\PY{p}{]}\PY{o}{.}\PY{n}{shape}\PY{p}{)}
\PY{n+nb}{print}\PY{p}{(}\PY{l+s+s2}{\PYZdq{}}\PY{l+s+s2}{gradients[}\PY{l+s+se}{\PYZbs{}\PYZdq{}}\PY{l+s+s2}{dWc}\PY{l+s+se}{\PYZbs{}\PYZdq{}}\PY{l+s+s2}{][3][1] =}\PY{l+s+s2}{\PYZdq{}}\PY{p}{,} \PY{n}{gradients\PYZus{}tmp}\PY{p}{[}\PY{l+s+s2}{\PYZdq{}}\PY{l+s+s2}{dWc}\PY{l+s+s2}{\PYZdq{}}\PY{p}{]}\PY{p}{[}\PY{l+m+mi}{3}\PY{p}{]}\PY{p}{[}\PY{l+m+mi}{1}\PY{p}{]}\PY{p}{)}
\PY{n+nb}{print}\PY{p}{(}\PY{l+s+s2}{\PYZdq{}}\PY{l+s+s2}{gradients[}\PY{l+s+se}{\PYZbs{}\PYZdq{}}\PY{l+s+s2}{dWc}\PY{l+s+se}{\PYZbs{}\PYZdq{}}\PY{l+s+s2}{].shape =}\PY{l+s+s2}{\PYZdq{}}\PY{p}{,} \PY{n}{gradients\PYZus{}tmp}\PY{p}{[}\PY{l+s+s2}{\PYZdq{}}\PY{l+s+s2}{dWc}\PY{l+s+s2}{\PYZdq{}}\PY{p}{]}\PY{o}{.}\PY{n}{shape}\PY{p}{)}
\PY{n+nb}{print}\PY{p}{(}\PY{l+s+s2}{\PYZdq{}}\PY{l+s+s2}{gradients[}\PY{l+s+se}{\PYZbs{}\PYZdq{}}\PY{l+s+s2}{dWo}\PY{l+s+se}{\PYZbs{}\PYZdq{}}\PY{l+s+s2}{][1][2] =}\PY{l+s+s2}{\PYZdq{}}\PY{p}{,} \PY{n}{gradients\PYZus{}tmp}\PY{p}{[}\PY{l+s+s2}{\PYZdq{}}\PY{l+s+s2}{dWo}\PY{l+s+s2}{\PYZdq{}}\PY{p}{]}\PY{p}{[}\PY{l+m+mi}{1}\PY{p}{]}\PY{p}{[}\PY{l+m+mi}{2}\PY{p}{]}\PY{p}{)}
\PY{n+nb}{print}\PY{p}{(}\PY{l+s+s2}{\PYZdq{}}\PY{l+s+s2}{gradients[}\PY{l+s+se}{\PYZbs{}\PYZdq{}}\PY{l+s+s2}{dWo}\PY{l+s+se}{\PYZbs{}\PYZdq{}}\PY{l+s+s2}{].shape =}\PY{l+s+s2}{\PYZdq{}}\PY{p}{,} \PY{n}{gradients\PYZus{}tmp}\PY{p}{[}\PY{l+s+s2}{\PYZdq{}}\PY{l+s+s2}{dWo}\PY{l+s+s2}{\PYZdq{}}\PY{p}{]}\PY{o}{.}\PY{n}{shape}\PY{p}{)}
\PY{n+nb}{print}\PY{p}{(}\PY{l+s+s2}{\PYZdq{}}\PY{l+s+s2}{gradients[}\PY{l+s+se}{\PYZbs{}\PYZdq{}}\PY{l+s+s2}{dbf}\PY{l+s+se}{\PYZbs{}\PYZdq{}}\PY{l+s+s2}{][4] =}\PY{l+s+s2}{\PYZdq{}}\PY{p}{,} \PY{n}{gradients\PYZus{}tmp}\PY{p}{[}\PY{l+s+s2}{\PYZdq{}}\PY{l+s+s2}{dbf}\PY{l+s+s2}{\PYZdq{}}\PY{p}{]}\PY{p}{[}\PY{l+m+mi}{4}\PY{p}{]}\PY{p}{)}
\PY{n+nb}{print}\PY{p}{(}\PY{l+s+s2}{\PYZdq{}}\PY{l+s+s2}{gradients[}\PY{l+s+se}{\PYZbs{}\PYZdq{}}\PY{l+s+s2}{dbf}\PY{l+s+se}{\PYZbs{}\PYZdq{}}\PY{l+s+s2}{].shape =}\PY{l+s+s2}{\PYZdq{}}\PY{p}{,} \PY{n}{gradients\PYZus{}tmp}\PY{p}{[}\PY{l+s+s2}{\PYZdq{}}\PY{l+s+s2}{dbf}\PY{l+s+s2}{\PYZdq{}}\PY{p}{]}\PY{o}{.}\PY{n}{shape}\PY{p}{)}
\PY{n+nb}{print}\PY{p}{(}\PY{l+s+s2}{\PYZdq{}}\PY{l+s+s2}{gradients[}\PY{l+s+se}{\PYZbs{}\PYZdq{}}\PY{l+s+s2}{dbi}\PY{l+s+se}{\PYZbs{}\PYZdq{}}\PY{l+s+s2}{][4] =}\PY{l+s+s2}{\PYZdq{}}\PY{p}{,} \PY{n}{gradients\PYZus{}tmp}\PY{p}{[}\PY{l+s+s2}{\PYZdq{}}\PY{l+s+s2}{dbi}\PY{l+s+s2}{\PYZdq{}}\PY{p}{]}\PY{p}{[}\PY{l+m+mi}{4}\PY{p}{]}\PY{p}{)}
\PY{n+nb}{print}\PY{p}{(}\PY{l+s+s2}{\PYZdq{}}\PY{l+s+s2}{gradients[}\PY{l+s+se}{\PYZbs{}\PYZdq{}}\PY{l+s+s2}{dbi}\PY{l+s+se}{\PYZbs{}\PYZdq{}}\PY{l+s+s2}{].shape =}\PY{l+s+s2}{\PYZdq{}}\PY{p}{,} \PY{n}{gradients\PYZus{}tmp}\PY{p}{[}\PY{l+s+s2}{\PYZdq{}}\PY{l+s+s2}{dbi}\PY{l+s+s2}{\PYZdq{}}\PY{p}{]}\PY{o}{.}\PY{n}{shape}\PY{p}{)}
\PY{n+nb}{print}\PY{p}{(}\PY{l+s+s2}{\PYZdq{}}\PY{l+s+s2}{gradients[}\PY{l+s+se}{\PYZbs{}\PYZdq{}}\PY{l+s+s2}{dbc}\PY{l+s+se}{\PYZbs{}\PYZdq{}}\PY{l+s+s2}{][4] =}\PY{l+s+s2}{\PYZdq{}}\PY{p}{,} \PY{n}{gradients\PYZus{}tmp}\PY{p}{[}\PY{l+s+s2}{\PYZdq{}}\PY{l+s+s2}{dbc}\PY{l+s+s2}{\PYZdq{}}\PY{p}{]}\PY{p}{[}\PY{l+m+mi}{4}\PY{p}{]}\PY{p}{)}
\PY{n+nb}{print}\PY{p}{(}\PY{l+s+s2}{\PYZdq{}}\PY{l+s+s2}{gradients[}\PY{l+s+se}{\PYZbs{}\PYZdq{}}\PY{l+s+s2}{dbc}\PY{l+s+se}{\PYZbs{}\PYZdq{}}\PY{l+s+s2}{].shape =}\PY{l+s+s2}{\PYZdq{}}\PY{p}{,} \PY{n}{gradients\PYZus{}tmp}\PY{p}{[}\PY{l+s+s2}{\PYZdq{}}\PY{l+s+s2}{dbc}\PY{l+s+s2}{\PYZdq{}}\PY{p}{]}\PY{o}{.}\PY{n}{shape}\PY{p}{)}
\PY{n+nb}{print}\PY{p}{(}\PY{l+s+s2}{\PYZdq{}}\PY{l+s+s2}{gradients[}\PY{l+s+se}{\PYZbs{}\PYZdq{}}\PY{l+s+s2}{dbo}\PY{l+s+se}{\PYZbs{}\PYZdq{}}\PY{l+s+s2}{][4] =}\PY{l+s+s2}{\PYZdq{}}\PY{p}{,} \PY{n}{gradients\PYZus{}tmp}\PY{p}{[}\PY{l+s+s2}{\PYZdq{}}\PY{l+s+s2}{dbo}\PY{l+s+s2}{\PYZdq{}}\PY{p}{]}\PY{p}{[}\PY{l+m+mi}{4}\PY{p}{]}\PY{p}{)}
\PY{n+nb}{print}\PY{p}{(}\PY{l+s+s2}{\PYZdq{}}\PY{l+s+s2}{gradients[}\PY{l+s+se}{\PYZbs{}\PYZdq{}}\PY{l+s+s2}{dbo}\PY{l+s+se}{\PYZbs{}\PYZdq{}}\PY{l+s+s2}{].shape =}\PY{l+s+s2}{\PYZdq{}}\PY{p}{,} \PY{n}{gradients\PYZus{}tmp}\PY{p}{[}\PY{l+s+s2}{\PYZdq{}}\PY{l+s+s2}{dbo}\PY{l+s+s2}{\PYZdq{}}\PY{p}{]}\PY{o}{.}\PY{n}{shape}\PY{p}{)}
\end{Verbatim}
\end{tcolorbox}

    \textbf{Expected Output}:

gradients{[}``dxt''{]}{[}1{]}{[}2{]} =

3.23055911511

\begin{verbatim}
<tr>
    <td>
        <b>gradients["dxt"].shape</b> =
    </td>
    <td>
       (3, 10)
    </td>
</tr>
    <tr>
    <td>
        <b>gradients["da_prev"][2][3]</b> =
    </td>
    <td>
       -0.0639621419711
    </td>
</tr>
    <tr>
    <td>
        <b>gradients["da_prev"].shape</b> =
    </td>
    <td>
       (5, 10)
    </td>
</tr>
     <tr>
    <td>
        <b>gradients["dc_prev"][2][3]</b> =
    </td>
    <td>
       0.797522038797
    </td>
</tr>
    <tr>
    <td>
        <b>gradients["dc_prev"].shape</b> =
    </td>
    <td>
       (5, 10)
    </td>
</tr>
    <tr>
    <td>
        <b>gradients["dWf"][3][1]</b> = 
    </td>
    <td>
       -0.147954838164
    </td>
</tr>
    <tr>
    <td>
        <b>gradients["dWf"].shape</b> =
    </td>
    <td>
       (5, 8)
    </td>
</tr>
    <tr>
    <td>
        <b>gradients["dWi"][1][2]</b> = 
    </td>
    <td>
       1.05749805523
    </td>
</tr>
    <tr>
    <td>
        <b>gradients["dWi"].shape</b> = 
    </td>
    <td>
       (5, 8)
    </td>
</tr>
<tr>
    <td>
        <b>gradients["dWc"][3][1]</b> = 
    </td>
    <td>
       2.30456216369
    </td>
</tr>
    <tr>
    <td>
        <b>gradients["dWc"].shape</b> = 
    </td>
    <td>
       (5, 8)
    </td>
</tr>
<tr>
    <td>
        <b>gradients["dWo"][1][2]</b> = 
    </td>
    <td>
       0.331311595289
    </td>
</tr>
    <tr>
    <td>
        <b>gradients["dWo"].shape</b> = 
    </td>
    <td>
       (5, 8)
    </td>
</tr>
<tr>
    <td>
        <b>gradients["dbf"][4]</b> = 
    </td>
    <td>
       [ 0.18864637]
    </td>
</tr>
    <tr>
    <td>
        <b>gradients["dbf"].shape</b> = 
    </td>
    <td>
       (5, 1)
    </td>
</tr>
<tr>
    <td>
        <b>gradients["dbi"][4]</b> = 
    </td>
    <td>
       [-0.40142491]
    </td>
</tr>
    <tr>
    <td>
        <b>gradients["dbi"].shape</b> = 
    </td>
    <td>
       (5, 1)
    </td>
</tr>
    <tr>
    <td>
        <b>gradients["dbc"][4]</b> = 
    </td>
    <td>
       [ 0.25587763]
    </td>
</tr>
    <tr>
    <td>
        <b>gradients["dbc"].shape</b> = 
    </td>
    <td>
       (5, 1)
    </td>
</tr>
    <tr>
    <td>
        <b>gradients["dbo"][4]</b> = 
    </td>
    <td>
       [ 0.13893342]
    </td>
</tr>
    <tr>
    <td>
        <b>gradients["dbo"].shape</b> = 
    </td>
    <td>
       (5, 1)
    </td>
</tr>
\end{verbatim}

    \#\#\# 3.3 Backward Pass through the LSTM RNN

This part is very similar to the \texttt{rnn\_backward} function you
implemented above. You will first create variables of the same dimension
as your return variables. You will then iterate over all the time steps
starting from the end and call the one step function you implemented for
LSTM at each iteration. You will then update the parameters by summing
them individually. Finally return a dictionary with the new gradients.

\#\#\# Exercise 8 - lstm\_backward

Implement the \texttt{lstm\_backward} function.

\textbf{Instructions}: Create a for loop starting from \(T_x\) and going
backward. For each step, call \texttt{lstm\_cell\_backward} and update
your old gradients by adding the new gradients to them. Note that
\texttt{dxt} is not updated, but is stored.

    \begin{tcolorbox}[breakable, size=fbox, boxrule=1pt, pad at break*=1mm,colback=cellbackground, colframe=cellborder]
\prompt{In}{incolor}{ }{\boxspacing}
\begin{Verbatim}[commandchars=\\\{\}]
\PY{c+c1}{\PYZsh{} UNGRADED FUNCTION: lstm\PYZus{}backward}

\PY{k}{def} \PY{n+nf}{lstm\PYZus{}backward}\PY{p}{(}\PY{n}{da}\PY{p}{,} \PY{n}{caches}\PY{p}{)}\PY{p}{:}
    
    \PY{l+s+sd}{\PYZdq{}\PYZdq{}\PYZdq{}}
\PY{l+s+sd}{    Implement the backward pass for the RNN with LSTM\PYZhy{}cell (over a whole sequence).}

\PY{l+s+sd}{    Arguments:}
\PY{l+s+sd}{    da \PYZhy{}\PYZhy{} Gradients w.r.t the hidden states, numpy\PYZhy{}array of shape (n\PYZus{}a, m, T\PYZus{}x)}
\PY{l+s+sd}{    caches \PYZhy{}\PYZhy{} cache storing information from the forward pass (lstm\PYZus{}forward)}

\PY{l+s+sd}{    Returns:}
\PY{l+s+sd}{    gradients \PYZhy{}\PYZhy{} python dictionary containing:}
\PY{l+s+sd}{                        dx \PYZhy{}\PYZhy{} Gradient of inputs, of shape (n\PYZus{}x, m, T\PYZus{}x)}
\PY{l+s+sd}{                        da0 \PYZhy{}\PYZhy{} Gradient w.r.t. the previous hidden state, numpy array of shape (n\PYZus{}a, m)}
\PY{l+s+sd}{                        dWf \PYZhy{}\PYZhy{} Gradient w.r.t. the weight matrix of the forget gate, numpy array of shape (n\PYZus{}a, n\PYZus{}a + n\PYZus{}x)}
\PY{l+s+sd}{                        dWi \PYZhy{}\PYZhy{} Gradient w.r.t. the weight matrix of the update gate, numpy array of shape (n\PYZus{}a, n\PYZus{}a + n\PYZus{}x)}
\PY{l+s+sd}{                        dWc \PYZhy{}\PYZhy{} Gradient w.r.t. the weight matrix of the memory gate, numpy array of shape (n\PYZus{}a, n\PYZus{}a + n\PYZus{}x)}
\PY{l+s+sd}{                        dWo \PYZhy{}\PYZhy{} Gradient w.r.t. the weight matrix of the output gate, numpy array of shape (n\PYZus{}a, n\PYZus{}a + n\PYZus{}x)}
\PY{l+s+sd}{                        dbf \PYZhy{}\PYZhy{} Gradient w.r.t. biases of the forget gate, of shape (n\PYZus{}a, 1)}
\PY{l+s+sd}{                        dbi \PYZhy{}\PYZhy{} Gradient w.r.t. biases of the update gate, of shape (n\PYZus{}a, 1)}
\PY{l+s+sd}{                        dbc \PYZhy{}\PYZhy{} Gradient w.r.t. biases of the memory gate, of shape (n\PYZus{}a, 1)}
\PY{l+s+sd}{                        dbo \PYZhy{}\PYZhy{} Gradient w.r.t. biases of the output gate, of shape (n\PYZus{}a, 1)}
\PY{l+s+sd}{    \PYZdq{}\PYZdq{}\PYZdq{}}

    \PY{c+c1}{\PYZsh{} Retrieve values from the first cache (t=1) of caches.}
    \PY{p}{(}\PY{n}{caches}\PY{p}{,} \PY{n}{x}\PY{p}{)} \PY{o}{=} \PY{n}{caches}
    \PY{p}{(}\PY{n}{a1}\PY{p}{,} \PY{n}{c1}\PY{p}{,} \PY{n}{a0}\PY{p}{,} \PY{n}{c0}\PY{p}{,} \PY{n}{f1}\PY{p}{,} \PY{n}{i1}\PY{p}{,} \PY{n}{cc1}\PY{p}{,} \PY{n}{o1}\PY{p}{,} \PY{n}{x1}\PY{p}{,} \PY{n}{parameters}\PY{p}{)} \PY{o}{=} \PY{n}{caches}\PY{p}{[}\PY{l+m+mi}{0}\PY{p}{]}
    
    \PY{c+c1}{\PYZsh{}\PYZsh{}\PYZsh{} START CODE HERE \PYZsh{}\PYZsh{}\PYZsh{}}
    \PY{c+c1}{\PYZsh{} Retrieve dimensions from da\PYZsq{}s and x1\PYZsq{}s shapes (≈2 lines)}
    \PY{n}{n\PYZus{}a}\PY{p}{,} \PY{n}{m}\PY{p}{,} \PY{n}{T\PYZus{}x} \PY{o}{=} \PY{k+kc}{None}
    \PY{n}{n\PYZus{}x}\PY{p}{,} \PY{n}{m} \PY{o}{=} \PY{k+kc}{None}
    
    \PY{c+c1}{\PYZsh{} initialize the gradients with the right sizes (≈12 lines)}
    \PY{n}{dx} \PY{o}{=} \PY{k+kc}{None}
    \PY{n}{da0} \PY{o}{=} \PY{k+kc}{None}
    \PY{n}{da\PYZus{}prevt} \PY{o}{=} \PY{k+kc}{None}
    \PY{n}{dc\PYZus{}prevt} \PY{o}{=} \PY{k+kc}{None}
    \PY{n}{dWf} \PY{o}{=} \PY{k+kc}{None}
    \PY{n}{dWi} \PY{o}{=} \PY{k+kc}{None}
    \PY{n}{dWc} \PY{o}{=} \PY{k+kc}{None}
    \PY{n}{dWo} \PY{o}{=} \PY{k+kc}{None}
    \PY{n}{dbf} \PY{o}{=} \PY{k+kc}{None}
    \PY{n}{dbi} \PY{o}{=} \PY{k+kc}{None}
    \PY{n}{dbc} \PY{o}{=} \PY{k+kc}{None}
    \PY{n}{dbo} \PY{o}{=} \PY{k+kc}{None}
    
    \PY{c+c1}{\PYZsh{} loop back over the whole sequence}
    \PY{k}{for} \PY{n}{t} \PY{o+ow}{in} \PY{n+nb}{reversed}\PY{p}{(}\PY{n+nb}{range}\PY{p}{(}\PY{k+kc}{None}\PY{p}{)}\PY{p}{)}\PY{p}{:}
        \PY{c+c1}{\PYZsh{} Compute all gradients using lstm\PYZus{}cell\PYZus{}backward. Choose wisely the \PYZdq{}da\PYZus{}next\PYZdq{} (same as done for Ex 6).}
        \PY{n}{gradients} \PY{o}{=} \PY{k+kc}{None}
        \PY{c+c1}{\PYZsh{} Store or add the gradient to the parameters\PYZsq{} previous step\PYZsq{}s gradient}
        \PY{n}{da\PYZus{}prevt} \PY{o}{=} \PY{k+kc}{None}
        \PY{n}{dc\PYZus{}prevt} \PY{o}{=} \PY{k+kc}{None}
        \PY{n}{dx}\PY{p}{[}\PY{p}{:}\PY{p}{,}\PY{p}{:}\PY{p}{,}\PY{n}{t}\PY{p}{]} \PY{o}{=} \PY{k+kc}{None}
        \PY{n}{dWf} \PY{o}{+}\PY{o}{=} \PY{k+kc}{None}
        \PY{n}{dWi} \PY{o}{+}\PY{o}{=} \PY{k+kc}{None}
        \PY{n}{dWc} \PY{o}{+}\PY{o}{=} \PY{k+kc}{None}
        \PY{n}{dWo} \PY{o}{+}\PY{o}{=} \PY{k+kc}{None}
        \PY{n}{dbf} \PY{o}{+}\PY{o}{=} \PY{k+kc}{None}
        \PY{n}{dbi} \PY{o}{+}\PY{o}{=} \PY{k+kc}{None}
        \PY{n}{dbc} \PY{o}{+}\PY{o}{=} \PY{k+kc}{None}
        \PY{n}{dbo} \PY{o}{+}\PY{o}{=} \PY{k+kc}{None}
    \PY{c+c1}{\PYZsh{} Set the first activation\PYZsq{}s gradient to the backpropagated gradient da\PYZus{}prev.}
    \PY{n}{da0} \PY{o}{=} \PY{k+kc}{None}
    
    \PY{c+c1}{\PYZsh{}\PYZsh{}\PYZsh{} END CODE HERE \PYZsh{}\PYZsh{}\PYZsh{}}

    \PY{c+c1}{\PYZsh{} Store the gradients in a python dictionary}
    \PY{n}{gradients} \PY{o}{=} \PY{p}{\PYZob{}}\PY{l+s+s2}{\PYZdq{}}\PY{l+s+s2}{dx}\PY{l+s+s2}{\PYZdq{}}\PY{p}{:} \PY{n}{dx}\PY{p}{,} \PY{l+s+s2}{\PYZdq{}}\PY{l+s+s2}{da0}\PY{l+s+s2}{\PYZdq{}}\PY{p}{:} \PY{n}{da0}\PY{p}{,} \PY{l+s+s2}{\PYZdq{}}\PY{l+s+s2}{dWf}\PY{l+s+s2}{\PYZdq{}}\PY{p}{:} \PY{n}{dWf}\PY{p}{,}\PY{l+s+s2}{\PYZdq{}}\PY{l+s+s2}{dbf}\PY{l+s+s2}{\PYZdq{}}\PY{p}{:} \PY{n}{dbf}\PY{p}{,} \PY{l+s+s2}{\PYZdq{}}\PY{l+s+s2}{dWi}\PY{l+s+s2}{\PYZdq{}}\PY{p}{:} \PY{n}{dWi}\PY{p}{,}\PY{l+s+s2}{\PYZdq{}}\PY{l+s+s2}{dbi}\PY{l+s+s2}{\PYZdq{}}\PY{p}{:} \PY{n}{dbi}\PY{p}{,}
                \PY{l+s+s2}{\PYZdq{}}\PY{l+s+s2}{dWc}\PY{l+s+s2}{\PYZdq{}}\PY{p}{:} \PY{n}{dWc}\PY{p}{,}\PY{l+s+s2}{\PYZdq{}}\PY{l+s+s2}{dbc}\PY{l+s+s2}{\PYZdq{}}\PY{p}{:} \PY{n}{dbc}\PY{p}{,} \PY{l+s+s2}{\PYZdq{}}\PY{l+s+s2}{dWo}\PY{l+s+s2}{\PYZdq{}}\PY{p}{:} \PY{n}{dWo}\PY{p}{,}\PY{l+s+s2}{\PYZdq{}}\PY{l+s+s2}{dbo}\PY{l+s+s2}{\PYZdq{}}\PY{p}{:} \PY{n}{dbo}\PY{p}{\PYZcb{}}
    
    \PY{k}{return} \PY{n}{gradients}
\end{Verbatim}
\end{tcolorbox}

    \begin{tcolorbox}[breakable, size=fbox, boxrule=1pt, pad at break*=1mm,colback=cellbackground, colframe=cellborder]
\prompt{In}{incolor}{ }{\boxspacing}
\begin{Verbatim}[commandchars=\\\{\}]
\PY{n}{np}\PY{o}{.}\PY{n}{random}\PY{o}{.}\PY{n}{seed}\PY{p}{(}\PY{l+m+mi}{1}\PY{p}{)}
\PY{n}{x\PYZus{}tmp} \PY{o}{=} \PY{n}{np}\PY{o}{.}\PY{n}{random}\PY{o}{.}\PY{n}{randn}\PY{p}{(}\PY{l+m+mi}{3}\PY{p}{,}\PY{l+m+mi}{10}\PY{p}{,}\PY{l+m+mi}{7}\PY{p}{)}
\PY{n}{a0\PYZus{}tmp} \PY{o}{=} \PY{n}{np}\PY{o}{.}\PY{n}{random}\PY{o}{.}\PY{n}{randn}\PY{p}{(}\PY{l+m+mi}{5}\PY{p}{,}\PY{l+m+mi}{10}\PY{p}{)}

\PY{n}{parameters\PYZus{}tmp} \PY{o}{=} \PY{p}{\PYZob{}}\PY{p}{\PYZcb{}}
\PY{n}{parameters\PYZus{}tmp}\PY{p}{[}\PY{l+s+s1}{\PYZsq{}}\PY{l+s+s1}{Wf}\PY{l+s+s1}{\PYZsq{}}\PY{p}{]} \PY{o}{=} \PY{n}{np}\PY{o}{.}\PY{n}{random}\PY{o}{.}\PY{n}{randn}\PY{p}{(}\PY{l+m+mi}{5}\PY{p}{,} \PY{l+m+mi}{5}\PY{o}{+}\PY{l+m+mi}{3}\PY{p}{)}
\PY{n}{parameters\PYZus{}tmp}\PY{p}{[}\PY{l+s+s1}{\PYZsq{}}\PY{l+s+s1}{bf}\PY{l+s+s1}{\PYZsq{}}\PY{p}{]} \PY{o}{=} \PY{n}{np}\PY{o}{.}\PY{n}{random}\PY{o}{.}\PY{n}{randn}\PY{p}{(}\PY{l+m+mi}{5}\PY{p}{,}\PY{l+m+mi}{1}\PY{p}{)}
\PY{n}{parameters\PYZus{}tmp}\PY{p}{[}\PY{l+s+s1}{\PYZsq{}}\PY{l+s+s1}{Wi}\PY{l+s+s1}{\PYZsq{}}\PY{p}{]} \PY{o}{=} \PY{n}{np}\PY{o}{.}\PY{n}{random}\PY{o}{.}\PY{n}{randn}\PY{p}{(}\PY{l+m+mi}{5}\PY{p}{,} \PY{l+m+mi}{5}\PY{o}{+}\PY{l+m+mi}{3}\PY{p}{)}
\PY{n}{parameters\PYZus{}tmp}\PY{p}{[}\PY{l+s+s1}{\PYZsq{}}\PY{l+s+s1}{bi}\PY{l+s+s1}{\PYZsq{}}\PY{p}{]} \PY{o}{=} \PY{n}{np}\PY{o}{.}\PY{n}{random}\PY{o}{.}\PY{n}{randn}\PY{p}{(}\PY{l+m+mi}{5}\PY{p}{,}\PY{l+m+mi}{1}\PY{p}{)}
\PY{n}{parameters\PYZus{}tmp}\PY{p}{[}\PY{l+s+s1}{\PYZsq{}}\PY{l+s+s1}{Wo}\PY{l+s+s1}{\PYZsq{}}\PY{p}{]} \PY{o}{=} \PY{n}{np}\PY{o}{.}\PY{n}{random}\PY{o}{.}\PY{n}{randn}\PY{p}{(}\PY{l+m+mi}{5}\PY{p}{,} \PY{l+m+mi}{5}\PY{o}{+}\PY{l+m+mi}{3}\PY{p}{)}
\PY{n}{parameters\PYZus{}tmp}\PY{p}{[}\PY{l+s+s1}{\PYZsq{}}\PY{l+s+s1}{bo}\PY{l+s+s1}{\PYZsq{}}\PY{p}{]} \PY{o}{=} \PY{n}{np}\PY{o}{.}\PY{n}{random}\PY{o}{.}\PY{n}{randn}\PY{p}{(}\PY{l+m+mi}{5}\PY{p}{,}\PY{l+m+mi}{1}\PY{p}{)}
\PY{n}{parameters\PYZus{}tmp}\PY{p}{[}\PY{l+s+s1}{\PYZsq{}}\PY{l+s+s1}{Wc}\PY{l+s+s1}{\PYZsq{}}\PY{p}{]} \PY{o}{=} \PY{n}{np}\PY{o}{.}\PY{n}{random}\PY{o}{.}\PY{n}{randn}\PY{p}{(}\PY{l+m+mi}{5}\PY{p}{,} \PY{l+m+mi}{5}\PY{o}{+}\PY{l+m+mi}{3}\PY{p}{)}
\PY{n}{parameters\PYZus{}tmp}\PY{p}{[}\PY{l+s+s1}{\PYZsq{}}\PY{l+s+s1}{bc}\PY{l+s+s1}{\PYZsq{}}\PY{p}{]} \PY{o}{=} \PY{n}{np}\PY{o}{.}\PY{n}{random}\PY{o}{.}\PY{n}{randn}\PY{p}{(}\PY{l+m+mi}{5}\PY{p}{,}\PY{l+m+mi}{1}\PY{p}{)}
\PY{n}{parameters\PYZus{}tmp}\PY{p}{[}\PY{l+s+s1}{\PYZsq{}}\PY{l+s+s1}{Wy}\PY{l+s+s1}{\PYZsq{}}\PY{p}{]} \PY{o}{=} \PY{n}{np}\PY{o}{.}\PY{n}{zeros}\PY{p}{(}\PY{p}{(}\PY{l+m+mi}{2}\PY{p}{,}\PY{l+m+mi}{5}\PY{p}{)}\PY{p}{)}       \PY{c+c1}{\PYZsh{} unused, but needed for lstm\PYZus{}forward}
\PY{n}{parameters\PYZus{}tmp}\PY{p}{[}\PY{l+s+s1}{\PYZsq{}}\PY{l+s+s1}{by}\PY{l+s+s1}{\PYZsq{}}\PY{p}{]} \PY{o}{=} \PY{n}{np}\PY{o}{.}\PY{n}{zeros}\PY{p}{(}\PY{p}{(}\PY{l+m+mi}{2}\PY{p}{,}\PY{l+m+mi}{1}\PY{p}{)}\PY{p}{)}       \PY{c+c1}{\PYZsh{} unused, but needed for lstm\PYZus{}forward}

\PY{n}{a\PYZus{}tmp}\PY{p}{,} \PY{n}{y\PYZus{}tmp}\PY{p}{,} \PY{n}{c\PYZus{}tmp}\PY{p}{,} \PY{n}{caches\PYZus{}tmp} \PY{o}{=} \PY{n}{lstm\PYZus{}forward}\PY{p}{(}\PY{n}{x\PYZus{}tmp}\PY{p}{,} \PY{n}{a0\PYZus{}tmp}\PY{p}{,} \PY{n}{parameters\PYZus{}tmp}\PY{p}{)}

\PY{n}{da\PYZus{}tmp} \PY{o}{=} \PY{n}{np}\PY{o}{.}\PY{n}{random}\PY{o}{.}\PY{n}{randn}\PY{p}{(}\PY{l+m+mi}{5}\PY{p}{,} \PY{l+m+mi}{10}\PY{p}{,} \PY{l+m+mi}{4}\PY{p}{)}
\PY{n}{gradients\PYZus{}tmp} \PY{o}{=} \PY{n}{lstm\PYZus{}backward}\PY{p}{(}\PY{n}{da\PYZus{}tmp}\PY{p}{,} \PY{n}{caches\PYZus{}tmp}\PY{p}{)}

\PY{n+nb}{print}\PY{p}{(}\PY{l+s+s2}{\PYZdq{}}\PY{l+s+s2}{gradients[}\PY{l+s+se}{\PYZbs{}\PYZdq{}}\PY{l+s+s2}{dx}\PY{l+s+se}{\PYZbs{}\PYZdq{}}\PY{l+s+s2}{][1][2] =}\PY{l+s+s2}{\PYZdq{}}\PY{p}{,} \PY{n}{gradients\PYZus{}tmp}\PY{p}{[}\PY{l+s+s2}{\PYZdq{}}\PY{l+s+s2}{dx}\PY{l+s+s2}{\PYZdq{}}\PY{p}{]}\PY{p}{[}\PY{l+m+mi}{1}\PY{p}{]}\PY{p}{[}\PY{l+m+mi}{2}\PY{p}{]}\PY{p}{)}
\PY{n+nb}{print}\PY{p}{(}\PY{l+s+s2}{\PYZdq{}}\PY{l+s+s2}{gradients[}\PY{l+s+se}{\PYZbs{}\PYZdq{}}\PY{l+s+s2}{dx}\PY{l+s+se}{\PYZbs{}\PYZdq{}}\PY{l+s+s2}{].shape =}\PY{l+s+s2}{\PYZdq{}}\PY{p}{,} \PY{n}{gradients\PYZus{}tmp}\PY{p}{[}\PY{l+s+s2}{\PYZdq{}}\PY{l+s+s2}{dx}\PY{l+s+s2}{\PYZdq{}}\PY{p}{]}\PY{o}{.}\PY{n}{shape}\PY{p}{)}
\PY{n+nb}{print}\PY{p}{(}\PY{l+s+s2}{\PYZdq{}}\PY{l+s+s2}{gradients[}\PY{l+s+se}{\PYZbs{}\PYZdq{}}\PY{l+s+s2}{da0}\PY{l+s+se}{\PYZbs{}\PYZdq{}}\PY{l+s+s2}{][2][3] =}\PY{l+s+s2}{\PYZdq{}}\PY{p}{,} \PY{n}{gradients\PYZus{}tmp}\PY{p}{[}\PY{l+s+s2}{\PYZdq{}}\PY{l+s+s2}{da0}\PY{l+s+s2}{\PYZdq{}}\PY{p}{]}\PY{p}{[}\PY{l+m+mi}{2}\PY{p}{]}\PY{p}{[}\PY{l+m+mi}{3}\PY{p}{]}\PY{p}{)}
\PY{n+nb}{print}\PY{p}{(}\PY{l+s+s2}{\PYZdq{}}\PY{l+s+s2}{gradients[}\PY{l+s+se}{\PYZbs{}\PYZdq{}}\PY{l+s+s2}{da0}\PY{l+s+se}{\PYZbs{}\PYZdq{}}\PY{l+s+s2}{].shape =}\PY{l+s+s2}{\PYZdq{}}\PY{p}{,} \PY{n}{gradients\PYZus{}tmp}\PY{p}{[}\PY{l+s+s2}{\PYZdq{}}\PY{l+s+s2}{da0}\PY{l+s+s2}{\PYZdq{}}\PY{p}{]}\PY{o}{.}\PY{n}{shape}\PY{p}{)}
\PY{n+nb}{print}\PY{p}{(}\PY{l+s+s2}{\PYZdq{}}\PY{l+s+s2}{gradients[}\PY{l+s+se}{\PYZbs{}\PYZdq{}}\PY{l+s+s2}{dWf}\PY{l+s+se}{\PYZbs{}\PYZdq{}}\PY{l+s+s2}{][3][1] =}\PY{l+s+s2}{\PYZdq{}}\PY{p}{,} \PY{n}{gradients\PYZus{}tmp}\PY{p}{[}\PY{l+s+s2}{\PYZdq{}}\PY{l+s+s2}{dWf}\PY{l+s+s2}{\PYZdq{}}\PY{p}{]}\PY{p}{[}\PY{l+m+mi}{3}\PY{p}{]}\PY{p}{[}\PY{l+m+mi}{1}\PY{p}{]}\PY{p}{)}
\PY{n+nb}{print}\PY{p}{(}\PY{l+s+s2}{\PYZdq{}}\PY{l+s+s2}{gradients[}\PY{l+s+se}{\PYZbs{}\PYZdq{}}\PY{l+s+s2}{dWf}\PY{l+s+se}{\PYZbs{}\PYZdq{}}\PY{l+s+s2}{].shape =}\PY{l+s+s2}{\PYZdq{}}\PY{p}{,} \PY{n}{gradients\PYZus{}tmp}\PY{p}{[}\PY{l+s+s2}{\PYZdq{}}\PY{l+s+s2}{dWf}\PY{l+s+s2}{\PYZdq{}}\PY{p}{]}\PY{o}{.}\PY{n}{shape}\PY{p}{)}
\PY{n+nb}{print}\PY{p}{(}\PY{l+s+s2}{\PYZdq{}}\PY{l+s+s2}{gradients[}\PY{l+s+se}{\PYZbs{}\PYZdq{}}\PY{l+s+s2}{dWi}\PY{l+s+se}{\PYZbs{}\PYZdq{}}\PY{l+s+s2}{][1][2] =}\PY{l+s+s2}{\PYZdq{}}\PY{p}{,} \PY{n}{gradients\PYZus{}tmp}\PY{p}{[}\PY{l+s+s2}{\PYZdq{}}\PY{l+s+s2}{dWi}\PY{l+s+s2}{\PYZdq{}}\PY{p}{]}\PY{p}{[}\PY{l+m+mi}{1}\PY{p}{]}\PY{p}{[}\PY{l+m+mi}{2}\PY{p}{]}\PY{p}{)}
\PY{n+nb}{print}\PY{p}{(}\PY{l+s+s2}{\PYZdq{}}\PY{l+s+s2}{gradients[}\PY{l+s+se}{\PYZbs{}\PYZdq{}}\PY{l+s+s2}{dWi}\PY{l+s+se}{\PYZbs{}\PYZdq{}}\PY{l+s+s2}{].shape =}\PY{l+s+s2}{\PYZdq{}}\PY{p}{,} \PY{n}{gradients\PYZus{}tmp}\PY{p}{[}\PY{l+s+s2}{\PYZdq{}}\PY{l+s+s2}{dWi}\PY{l+s+s2}{\PYZdq{}}\PY{p}{]}\PY{o}{.}\PY{n}{shape}\PY{p}{)}
\PY{n+nb}{print}\PY{p}{(}\PY{l+s+s2}{\PYZdq{}}\PY{l+s+s2}{gradients[}\PY{l+s+se}{\PYZbs{}\PYZdq{}}\PY{l+s+s2}{dWc}\PY{l+s+se}{\PYZbs{}\PYZdq{}}\PY{l+s+s2}{][3][1] =}\PY{l+s+s2}{\PYZdq{}}\PY{p}{,} \PY{n}{gradients\PYZus{}tmp}\PY{p}{[}\PY{l+s+s2}{\PYZdq{}}\PY{l+s+s2}{dWc}\PY{l+s+s2}{\PYZdq{}}\PY{p}{]}\PY{p}{[}\PY{l+m+mi}{3}\PY{p}{]}\PY{p}{[}\PY{l+m+mi}{1}\PY{p}{]}\PY{p}{)}
\PY{n+nb}{print}\PY{p}{(}\PY{l+s+s2}{\PYZdq{}}\PY{l+s+s2}{gradients[}\PY{l+s+se}{\PYZbs{}\PYZdq{}}\PY{l+s+s2}{dWc}\PY{l+s+se}{\PYZbs{}\PYZdq{}}\PY{l+s+s2}{].shape =}\PY{l+s+s2}{\PYZdq{}}\PY{p}{,} \PY{n}{gradients\PYZus{}tmp}\PY{p}{[}\PY{l+s+s2}{\PYZdq{}}\PY{l+s+s2}{dWc}\PY{l+s+s2}{\PYZdq{}}\PY{p}{]}\PY{o}{.}\PY{n}{shape}\PY{p}{)}
\PY{n+nb}{print}\PY{p}{(}\PY{l+s+s2}{\PYZdq{}}\PY{l+s+s2}{gradients[}\PY{l+s+se}{\PYZbs{}\PYZdq{}}\PY{l+s+s2}{dWo}\PY{l+s+se}{\PYZbs{}\PYZdq{}}\PY{l+s+s2}{][1][2] =}\PY{l+s+s2}{\PYZdq{}}\PY{p}{,} \PY{n}{gradients\PYZus{}tmp}\PY{p}{[}\PY{l+s+s2}{\PYZdq{}}\PY{l+s+s2}{dWo}\PY{l+s+s2}{\PYZdq{}}\PY{p}{]}\PY{p}{[}\PY{l+m+mi}{1}\PY{p}{]}\PY{p}{[}\PY{l+m+mi}{2}\PY{p}{]}\PY{p}{)}
\PY{n+nb}{print}\PY{p}{(}\PY{l+s+s2}{\PYZdq{}}\PY{l+s+s2}{gradients[}\PY{l+s+se}{\PYZbs{}\PYZdq{}}\PY{l+s+s2}{dWo}\PY{l+s+se}{\PYZbs{}\PYZdq{}}\PY{l+s+s2}{].shape =}\PY{l+s+s2}{\PYZdq{}}\PY{p}{,} \PY{n}{gradients\PYZus{}tmp}\PY{p}{[}\PY{l+s+s2}{\PYZdq{}}\PY{l+s+s2}{dWo}\PY{l+s+s2}{\PYZdq{}}\PY{p}{]}\PY{o}{.}\PY{n}{shape}\PY{p}{)}
\PY{n+nb}{print}\PY{p}{(}\PY{l+s+s2}{\PYZdq{}}\PY{l+s+s2}{gradients[}\PY{l+s+se}{\PYZbs{}\PYZdq{}}\PY{l+s+s2}{dbf}\PY{l+s+se}{\PYZbs{}\PYZdq{}}\PY{l+s+s2}{][4] =}\PY{l+s+s2}{\PYZdq{}}\PY{p}{,} \PY{n}{gradients\PYZus{}tmp}\PY{p}{[}\PY{l+s+s2}{\PYZdq{}}\PY{l+s+s2}{dbf}\PY{l+s+s2}{\PYZdq{}}\PY{p}{]}\PY{p}{[}\PY{l+m+mi}{4}\PY{p}{]}\PY{p}{)}
\PY{n+nb}{print}\PY{p}{(}\PY{l+s+s2}{\PYZdq{}}\PY{l+s+s2}{gradients[}\PY{l+s+se}{\PYZbs{}\PYZdq{}}\PY{l+s+s2}{dbf}\PY{l+s+se}{\PYZbs{}\PYZdq{}}\PY{l+s+s2}{].shape =}\PY{l+s+s2}{\PYZdq{}}\PY{p}{,} \PY{n}{gradients\PYZus{}tmp}\PY{p}{[}\PY{l+s+s2}{\PYZdq{}}\PY{l+s+s2}{dbf}\PY{l+s+s2}{\PYZdq{}}\PY{p}{]}\PY{o}{.}\PY{n}{shape}\PY{p}{)}
\PY{n+nb}{print}\PY{p}{(}\PY{l+s+s2}{\PYZdq{}}\PY{l+s+s2}{gradients[}\PY{l+s+se}{\PYZbs{}\PYZdq{}}\PY{l+s+s2}{dbi}\PY{l+s+se}{\PYZbs{}\PYZdq{}}\PY{l+s+s2}{][4] =}\PY{l+s+s2}{\PYZdq{}}\PY{p}{,} \PY{n}{gradients\PYZus{}tmp}\PY{p}{[}\PY{l+s+s2}{\PYZdq{}}\PY{l+s+s2}{dbi}\PY{l+s+s2}{\PYZdq{}}\PY{p}{]}\PY{p}{[}\PY{l+m+mi}{4}\PY{p}{]}\PY{p}{)}
\PY{n+nb}{print}\PY{p}{(}\PY{l+s+s2}{\PYZdq{}}\PY{l+s+s2}{gradients[}\PY{l+s+se}{\PYZbs{}\PYZdq{}}\PY{l+s+s2}{dbi}\PY{l+s+se}{\PYZbs{}\PYZdq{}}\PY{l+s+s2}{].shape =}\PY{l+s+s2}{\PYZdq{}}\PY{p}{,} \PY{n}{gradients\PYZus{}tmp}\PY{p}{[}\PY{l+s+s2}{\PYZdq{}}\PY{l+s+s2}{dbi}\PY{l+s+s2}{\PYZdq{}}\PY{p}{]}\PY{o}{.}\PY{n}{shape}\PY{p}{)}
\PY{n+nb}{print}\PY{p}{(}\PY{l+s+s2}{\PYZdq{}}\PY{l+s+s2}{gradients[}\PY{l+s+se}{\PYZbs{}\PYZdq{}}\PY{l+s+s2}{dbc}\PY{l+s+se}{\PYZbs{}\PYZdq{}}\PY{l+s+s2}{][4] =}\PY{l+s+s2}{\PYZdq{}}\PY{p}{,} \PY{n}{gradients\PYZus{}tmp}\PY{p}{[}\PY{l+s+s2}{\PYZdq{}}\PY{l+s+s2}{dbc}\PY{l+s+s2}{\PYZdq{}}\PY{p}{]}\PY{p}{[}\PY{l+m+mi}{4}\PY{p}{]}\PY{p}{)}
\PY{n+nb}{print}\PY{p}{(}\PY{l+s+s2}{\PYZdq{}}\PY{l+s+s2}{gradients[}\PY{l+s+se}{\PYZbs{}\PYZdq{}}\PY{l+s+s2}{dbc}\PY{l+s+se}{\PYZbs{}\PYZdq{}}\PY{l+s+s2}{].shape =}\PY{l+s+s2}{\PYZdq{}}\PY{p}{,} \PY{n}{gradients\PYZus{}tmp}\PY{p}{[}\PY{l+s+s2}{\PYZdq{}}\PY{l+s+s2}{dbc}\PY{l+s+s2}{\PYZdq{}}\PY{p}{]}\PY{o}{.}\PY{n}{shape}\PY{p}{)}
\PY{n+nb}{print}\PY{p}{(}\PY{l+s+s2}{\PYZdq{}}\PY{l+s+s2}{gradients[}\PY{l+s+se}{\PYZbs{}\PYZdq{}}\PY{l+s+s2}{dbo}\PY{l+s+se}{\PYZbs{}\PYZdq{}}\PY{l+s+s2}{][4] =}\PY{l+s+s2}{\PYZdq{}}\PY{p}{,} \PY{n}{gradients\PYZus{}tmp}\PY{p}{[}\PY{l+s+s2}{\PYZdq{}}\PY{l+s+s2}{dbo}\PY{l+s+s2}{\PYZdq{}}\PY{p}{]}\PY{p}{[}\PY{l+m+mi}{4}\PY{p}{]}\PY{p}{)}
\PY{n+nb}{print}\PY{p}{(}\PY{l+s+s2}{\PYZdq{}}\PY{l+s+s2}{gradients[}\PY{l+s+se}{\PYZbs{}\PYZdq{}}\PY{l+s+s2}{dbo}\PY{l+s+se}{\PYZbs{}\PYZdq{}}\PY{l+s+s2}{].shape =}\PY{l+s+s2}{\PYZdq{}}\PY{p}{,} \PY{n}{gradients\PYZus{}tmp}\PY{p}{[}\PY{l+s+s2}{\PYZdq{}}\PY{l+s+s2}{dbo}\PY{l+s+s2}{\PYZdq{}}\PY{p}{]}\PY{o}{.}\PY{n}{shape}\PY{p}{)}
\end{Verbatim}
\end{tcolorbox}

    \textbf{Expected Output}:

gradients{[}``dx''{]}{[}1{]}{[}2{]} =

{[}0.00218254 0.28205375 -0.48292508 -0.43281115{]}

\begin{verbatim}
<tr>
    <td>
        <b>gradients["dx"].shape</b> =
    </td>
    <td>
       (3, 10, 4)
    </td>
</tr>
    <tr>
    <td>
        <b>gradients["da0"][2][3]</b> =
    </td>
    <td>
       0.312770310257
    </td>
</tr>
    <tr>
    <td>
        <b>gradients["da0"].shape</b> =
    </td>
    <td>
       (5, 10)
    </td>
</tr>
    <tr>
    <td>
        <b>gradients["dWf"][3][1]</b> = 
    </td>
    <td>
       -0.0809802310938
    </td>
</tr>
    <tr>
    <td>
        <b>gradients["dWf"].shape</b> =
    </td>
    <td>
       (5, 8)
    </td>
</tr>
    <tr>
    <td>
        <b>gradients["dWi"][1][2]</b> = 
    </td>
    <td>
       0.40512433093
    </td>
</tr>
    <tr>
    <td>
        <b>gradients["dWi"].shape</b> = 
    </td>
    <td>
       (5, 8)
    </td>
</tr>
<tr>
    <td>
        <b>gradients["dWc"][3][1]</b> = 
    </td>
    <td>
       -0.0793746735512
    </td>
</tr>
    <tr>
    <td>
        <b>gradients["dWc"].shape</b> = 
    </td>
    <td>
       (5, 8)
    </td>
</tr>
<tr>
    <td>
        <b>gradients["dWo"][1][2]</b> = 
    </td>
    <td>
       0.038948775763
    </td>
</tr>
    <tr>
    <td>
        <b>gradients["dWo"].shape</b> = 
    </td>
    <td>
       (5, 8)
    </td>
</tr>
<tr>
    <td>
        <b>gradients["dbf"][4]</b> = 
    </td>
    <td>
       [-0.15745657]
    </td>
</tr>
    <tr>
    <td>
        <b>gradients["dbf"].shape</b> = 
    </td>
    <td>
       (5, 1)
    </td>
</tr>
<tr>
    <td>
        <b>gradients["dbi"][4]</b> = 
    </td>
    <td>
       [-0.50848333]
    </td>
</tr>
    <tr>
    <td>
        <b>gradients["dbi"].shape</b> = 
    </td>
    <td>
       (5, 1)
    </td>
</tr>
    <tr>
    <td>
        <b>gradients["dbc"][4]</b> = 
    </td>
    <td>
       [-0.42510818]
    </td>
</tr>
    <tr>
    <td>
        <b>gradients["dbc"].shape</b> = 
    </td>
    <td>
       (5, 1)
    </td>
</tr>
    <tr>
    <td>
        <b>gradients["dbo"][4]</b> = 
    </td>
    <td>
       [ -0.17958196]
    </td>
</tr>
    <tr>
    <td>
        <b>gradients["dbo"].shape</b> = 
    </td>
    <td>
       (5, 1)
    </td>
</tr>
\end{verbatim}

    \hypertarget{congratulations-on-completing-this-assignment}{%
\subsubsection{Congratulations on completing this
assignment!}\label{congratulations-on-completing-this-assignment}}

You now understand how recurrent neural networks work! In the next
exercise, you'll use an RNN to build a character-level language model.
See you there!


    % Add a bibliography block to the postdoc
    
    
    
\end{document}
